%============================================================
%  Bd. 33 - Gruppen und Ringe
%============================================================

\documentclass[main.tex]{subfiles}

\ifSubfilesClassLoaded{
  \BandSetupStandalone{B33}
}{
}

\title{Bd. 33 - Gruppen und Ringe}
\author{Martin Kunze}
\date{}
\setcounter{file}{33}

\hypersetup{
  pdftitle={Bd. 33 - Gruppen und Ringe},
  pdfauthor={Martin Kunze}
}

\begin{document}

\newcommand{\BandXXIXProofRefStack}[1]{%
  \ensuremath{\begin{gathered}#1\end{gathered}}}

\title{Bd. 33 - Gruppen und Ringe}
\author{Martin Kunze}
\date{}
\hypersetup{pageanchor=false}
\maketitle
\hypersetup{pageanchor=true}
\setcounter{tocdepth}{3}
\tableofcontents
\setcounter{file}{33}
\renewcommand{\FormulaBandID}{33}

\chapter{Einleitung}

Die Bände~22--30 entwickeln Halbgruppen und ihre Spezialisierungen,
Band~31 Monoide und Band~32 Halbringe.
Der vorliegende Band ergänzt zunächst die Existenz inverser Elemente und
gelangt so zu Gruppen und abelschen Gruppen. Danach wird die Konstruktion
der ganzen Zahlen aus Band~17 verallgemeinert: Jedes kommutative kürzbare
Monoid wird durch formale Differenzen in eine abelsche Gruppe eingebettet.
Anschließend werden zwei Operationen auf demselben Träger miteinander
verbunden: Die Addition bildet eine abelsche Gruppe, die Multiplikation ein
Monoid, und beide Operationen erfüllen die Distributivgesetze. Dies führt zu
Ringen mit Eins und zu kommutativen Ringen mit Eins.

Die in Band~17 konstruierten ganzen Zahlen bilden das grundlegende Beispiel.
Ihre bereits aus der Quotientenkonstruktion hergeleiteten Rechengesetze
werden am Ende dieses Bandes zu den entsprechenden algebraischen
Strukturaussagen zusammengefasst.

Wie in den Bänden~22--24 gehört jede Operation obligatorisch zur
Strukturbezeichnung. Daher schreiben wir beispielsweise
\(\GroupAlg{G}{\circ}{e}\) und
\(\RingAlg{R}{\oplus}{\odot}{0_R}{1_R}\), niemals bloß
\(\mathsf{Grp}(G)\) oder \(\mathsf{Ring}(R)\).

\chapter{Gruppen}

\section{Gruppen als Monoide mit Inversen}

\FormulaAxiomDeltaK[Existenz eines beidseitigen Inversen]{%
  \begin{aligned}[t]
    &\MonoidAlg{G}{\circ}{e}\dsep x\in G\vdash{}\\
    &\exists y\in G\,
      \bigl(x\circ y=e\land y\circ x=e\bigr)
  \end{aligned}
}{GroupInverseExistenceLaw}{%
  \DeltaRow{Mengen}{G\dsep e\dsep x\dsep y}
  \DeltaRow{Funktionen}{\circ}
}

\FormulaDefDeltaK[Begriff der Gruppe]{%
  \GroupAlg{G}{\circ}{e}%
}{GroupDef}{%
  \DeltaRow{Mengen}{G\dsep e\dsep x\dsep y}
  \DeltaRow{Funktionen}{\circ}
  \DeltaRow{\textbf{Axiome}}{}
  \DeltaRow{Monoid}{\MonoidAlg{G}{\circ}{e}}
    [\FormulaRefAuto{MonoidDef}[def]]
  \DeltaRow{Beidseitige Inverse}{%
    \begin{aligned}[t]
      &x\in G\vdash{}\\[-2pt]
      &\exists y\in G\,
        \bigl(x\circ y=e\land y\circ x=e\bigr)
    \end{aligned}}
    [\FormulaRefAuto{GroupInverseExistenceLaw}[ax]]
  \DeltaRow{\textbf{Neues Symbol}}{}
  \DeltaPrem{Gruppe}{\GroupAlg{G}{\circ}{e}}
}

\FormulaAxiomDeltaK[Monoidanteil einer Gruppe]{%
  \GroupAlg{G}{\circ}{e}\vdash\MonoidAlg{G}{\circ}{e}%
}{GroupBaseMonoidAxiom}{%
  \DeltaRow{Mengen}{G\dsep e}
  \DeltaRow{Funktionen}{\circ}
}

\FormulaAxiomDeltaK[Inverse in einer Gruppe]{%
  \begin{aligned}[t]
    &\GroupAlg{G}{\circ}{e}\dsep x\in G\vdash{}\\
    &\exists y\in G\,
      \bigl(x\circ y=e\land y\circ x=e\bigr)
  \end{aligned}
}{GroupInverseExistenceAxiom}{%
  \DeltaRow{Mengen}{G\dsep e\dsep x\dsep y}
  \DeltaRow{Funktionen}{\circ}
}

\section{Grundlegende Gruppengesetze}

\FormulaThmDeltaK[Halbgruppenanteil einer Gruppe]{%
  \GroupAlg{G}{\circ}{e}\vdash\SemigroupAlg{G}{\circ}%
}{GroupBaseSemigroup}{%
  \DeltaRow{Mengen}{G\dsep e}
  \DeltaRow{Funktionen}{\circ}
}
\begin{tabproof}
  \proofstep{1}{\GroupAlg{G}{\circ}{e}}{\rA}
  \proofstep{1}{\MonoidAlg{G}{\circ}{e}}{%
    \FormulaRefAuto{GroupBaseMonoidAxiom}[ax]{1}}
  \proofstep{1}{\SemigroupAlg{G}{\circ}}{%
    \FormulaRefAuto{MonoidBaseSemigroupAxiom}[ax]{2}}
\end{tabproof}

\FormulaThmDeltaK[Abgeschlossenheit der Gruppenoperation]{%
  \GroupAlg{G}{\circ}{e}\dsep x,y\in G\vdash x\circ y\in G%
}{GroupClosure}{%
  \DeltaRow{Mengen}{G\dsep e\dsep x\dsep y}
  \DeltaRow{Funktionen}{\circ}
}
\begin{tabproof}
  \proofstep{1}{\GroupAlg{G}{\circ}{e}}{\rA}
  \proofstep{2}{x\in G}{\rA}
  \proofstep{3}{y\in G}{\rA}
  \proofstep{1}{\SemigroupAlg{G}{\circ}}{%
    \FormulaRefAuto{GroupBaseSemigroup}[thm]{1}}
  \proofstep{1,2,3}{x\circ y\in G}{%
    \FormulaRefAuto{SemigroupClosureAxiom}[ax]{4,2,3}}
\end{tabproof}

\FormulaThmDeltaK[Assoziativität der Gruppenoperation]{%
  \begin{aligned}[t]
    &\GroupAlg{G}{\circ}{e}\dsep x,y,z\in G\vdash{}\\
    &(x\circ y)\circ z=x\circ(y\circ z)
  \end{aligned}
}{GroupAssociativity}{%
  \DeltaRow{Mengen}{G\dsep e\dsep x\dsep y\dsep z}
  \DeltaRow{Funktionen}{\circ}
}
\begin{tabproofwide}
  \proofstepwidestar[1]{\GroupAlg{G}{\circ}{e}}{\rA}
  \proofstepwidestar[2]{x\in G}{\rA}
  \proofstepwidestar[3]{y\in G}{\rA}
  \proofstepwidestar[4]{z\in G}{\rA}
  \proofstepwidestar[1]{\SemigroupAlg{G}{\circ}}{%
    \FormulaRefAuto{GroupBaseSemigroup}[thm]{1}}
  \proofstepwidestar[1,2,3,4]{(x\circ y)\circ z=x\circ(y\circ z)}{%
    \FormulaRefAuto{SemigroupAssociativityAxiom}[ax]{5,2,3,4}}
\end{tabproofwide}

\FormulaThmDeltaK[Neutrales Element liegt im Gruppenträger]{%
  \GroupAlg{G}{\circ}{e}\vdash e\in G%
}{GroupIdentityCarrier}{%
  \DeltaRow{Mengen}{G\dsep e}
  \DeltaRow{Funktionen}{\circ}
}
\begin{tabproof}
  \proofstep{1}{\GroupAlg{G}{\circ}{e}}{\rA}
  \proofstep{1}{\MonoidAlg{G}{\circ}{e}}{%
    \FormulaRefAuto{GroupBaseMonoidAxiom}[ax]{1}}
  \proofstep{1}{e\in G}{%
    \FormulaRefAuto{MonoidIdentityCarrierAxiom}[ax]{2}}
\end{tabproof}

\FormulaThmDeltaK[Linksneutralität in einer Gruppe]{%
  \GroupAlg{G}{\circ}{e}\dsep x\in G\vdash e\circ x=x%
}{GroupLeftIdentity}{%
  \DeltaRow{Mengen}{G\dsep e\dsep x}
  \DeltaRow{Funktionen}{\circ}
}
\begin{tabproof}
  \proofstep{1}{\GroupAlg{G}{\circ}{e}}{\rA}
  \proofstep{2}{x\in G}{\rA}
  \proofstep{1}{\MonoidAlg{G}{\circ}{e}}{%
    \FormulaRefAuto{GroupBaseMonoidAxiom}[ax]{1}}
  \proofstep{1,2}{e\circ x=x}{%
    \FormulaRefAuto{MonoidLeftIdentityAxiom}[ax]{3,2}}
\end{tabproof}

\FormulaThmDeltaK[Rechtsneutralität in einer Gruppe]{%
  \GroupAlg{G}{\circ}{e}\dsep x\in G\vdash x\circ e=x%
}{GroupRightIdentity}{%
  \DeltaRow{Mengen}{G\dsep e\dsep x}
  \DeltaRow{Funktionen}{\circ}
}
\begin{tabproof}
  \proofstep{1}{\GroupAlg{G}{\circ}{e}}{\rA}
  \proofstep{2}{x\in G}{\rA}
  \proofstep{1}{\MonoidAlg{G}{\circ}{e}}{%
    \FormulaRefAuto{GroupBaseMonoidAxiom}[ax]{1}}
  \proofstep{1,2}{x\circ e=x}{%
    \FormulaRefAuto{MonoidRightIdentityAxiom}[ax]{3,2}}
\end{tabproof}

\section{Eindeutigkeit des Inversen}

\FormulaThmDeltaK[Eindeutigkeit aus Links- und Rechtsinverse]{%
  \begin{aligned}[t]
    &\GroupAlg{G}{\circ}{e}\dsep x,y,z\in G\dsep
      y\circ x=e\dsep x\circ z=e\\
    &\vdash y=z
  \end{aligned}
}{GroupInverseUnique}{%
  \DeltaRow{Mengen}{G\dsep e\dsep x\dsep y\dsep z}
  \DeltaRow{Funktionen}{\circ}
}
\begin{tabproofsplitwide}
  \proofpartwideindR[Einsetzen der Rechtsinverse]{%
    \begin{aligned}[t]
      &\GroupAlg{G}{\circ}{e}\dsep x,y,z\in G\dsep
        y\circ x=e\dsep x\circ z=e\\
      &\vdash y=y\circ(x\circ z)
    \end{aligned}
  }{%
    \GroupAlg{G}{\circ}{e}\dsep x,y,z\in G\dsep
    y\circ x=e\dsep x\circ z=e
    \vdash y=y\circ(x\circ z)%
  }[GroupInverseUniqueRightInverseSubstitution]
    \proofstepwidestar[1]{\GroupAlg{G}{\circ}{e}}{\rA}
    \proofstepwidestar[2]{x\in G}{\rA}
    \proofstepwidestar[3]{y\in G}{\rA}
    \proofstepwidestar[4]{z\in G}{\rA}
    \proofstepwidestar[5]{y\circ x=e}{\rA}
    \proofstepwidestar[6]{x\circ z=e}{\rA}
    \proofstepwide[1,3]{y}{=}{y\circ e}{%
      \FormulaRefAuto{a=b\vdash b=a}{%
        \FormulaRefAuto{GroupRightIdentity}[thm]{1,3}}}
    \proofstepwide[6]{y\circ e}{=}{y\circ(x\circ z)}{\rIE{6}}
    \proofstepwidestar[1,2,3,4,5,6]{%
      y=y\circ(x\circ z)}{\rChain{7,8}}
  \closeproofpartwideindR

  \proofpartwideindR[Reduktion mit der Linksinversen]{%
    \begin{aligned}[t]
      &\GroupAlg{G}{\circ}{e}\dsep x,y,z\in G\dsep
        y\circ x=e\dsep x\circ z=e\\
      &\vdash y\circ(x\circ z)=z
    \end{aligned}
  }{%
    \GroupAlg{G}{\circ}{e}\dsep x,y,z\in G\dsep
    y\circ x=e\dsep x\circ z=e
    \vdash y\circ(x\circ z)=z%
  }[GroupInverseUniqueLeftInverseReduction]
    \proofstepwidestar[1]{\GroupAlg{G}{\circ}{e}}{\rA}
    \proofstepwidestar[2]{x\in G}{\rA}
    \proofstepwidestar[3]{y\in G}{\rA}
    \proofstepwidestar[4]{z\in G}{\rA}
    \proofstepwidestar[5]{y\circ x=e}{\rA}
    \proofstepwidestar[6]{x\circ z=e}{\rA}
    \proofstepwide[1,2,3,4]{y\circ(x\circ z)}{=}{%
      (y\circ x)\circ z}{%
      \FormulaRefAuto{a=b\vdash b=a}{%
        \FormulaRefAuto{GroupAssociativity}[thm]{1,3,2,4}}}
    \proofstepwide[5]{(y\circ x)\circ z}{=}{e\circ z}{\rIE{5}}
    \proofstepwide[1,4]{e\circ z}{=}{z}{%
      \FormulaRefAuto{GroupLeftIdentity}[thm]{1,4}}
    \proofstepwidestar[1,2,3,4,5,6]{%
      y\circ(x\circ z)=z}{\rChain{7,8,9}}
  \closeproofpartwideindR

  \proofpartwideindR[Verknüpfung beider Teilrechnungen]{%
    \begin{aligned}[t]
      &\GroupAlg{G}{\circ}{e}\dsep x,y,z\in G\dsep
        y\circ x=e\dsep x\circ z=e\\
      &\vdash y=z
    \end{aligned}
  }{%
    \GroupAlg{G}{\circ}{e}\dsep x,y,z\in G\dsep
    y\circ x=e\dsep x\circ z=e
    \vdash y=z%
  }[GroupInverseUniqueConclusion]
    \proofstepwidestar[1]{\GroupAlg{G}{\circ}{e}}{\rA}
    \proofstepwidestar[2]{x\in G}{\rA}
    \proofstepwidestar[3]{y\in G}{\rA}
    \proofstepwidestar[4]{z\in G}{\rA}
    \proofstepwidestar[5]{y\circ x=e}{\rA}
    \proofstepwidestar[6]{x\circ z=e}{\rA}
    \proofstepwidestar[1,2,3,4,5,6]{%
      y=y\circ(x\circ z)}{%
      \FormulaRefAuto{GroupInverseUniqueRightInverseSubstitution}{%
        1,2,3,4,5,6}}
    \proofstepwidestar[1,2,3,4,5,6]{%
      y\circ(x\circ z)=z}{%
      \FormulaRefAuto{GroupInverseUniqueLeftInverseReduction}{%
        1,2,3,4,5,6}}
    \proofstepwidestar[1,2,3,4,5,6]{y=z}{\rChain{7,8}}
  \closeproofpartwideindR
\end{tabproofsplitwide}

\FormulaDefDeltaK[Inverses Element]{%
  \begin{aligned}[t]
    &\GroupAlg{G}{\circ}{e}\dsep x\in G\vdash{}\\
    &x^{-1}\coloneqq
      \iota y\,
      \bigl(y\in G\land x\circ y=e\land y\circ x=e\bigr)
  \end{aligned}
}{GroupInverseElement}{%
  \DeltaRow{Mengen}{G\dsep e\dsep x\dsep y}
  \DeltaRow{Funktionen}{\circ}
  \DeltaRow{Träger}{x^{-1}\in G}
  \DeltaRow{Rechtsinverse}{x\circ x^{-1}=e}
  \DeltaRow{Linksinverse}{x^{-1}\circ x=e}
  \DeltaRow{Eindeutigkeit}{\FormulaRefAuto{GroupInverseUnique}[thm]}
  \DeltaRow{\textbf{Neues Symbol}}{}
  \DeltaPrem{Inverses Element}{x^{-1}}
}

\FormulaThmDeltaK[Inverses Element liegt im Gruppenträger]{%
  \GroupAlg{G}{\circ}{e}\dsep x\in G\vdash x^{-1}\in G%
}{GroupInverseCarrier}{%
  \DeltaRow{Mengen}{G\dsep e\dsep x}
  \DeltaRow{Funktionen}{\circ}
}
\begin{tabproof}
  \proofstep{1}{\GroupAlg{G}{\circ}{e}}{\rA}
  \proofstep{2}{x\in G}{\rA}
  \proofstep{1,2}{x^{-1}\in G}{%
    \FormulaRefAuto{GroupInverseElement}[def]{1,2}}
\end{tabproof}

\FormulaThmDeltaK[Rechtsinverse Gleichung]{%
  \GroupAlg{G}{\circ}{e}\dsep x\in G\vdash x\circ x^{-1}=e%
}{GroupInverseRight}{%
  \DeltaRow{Mengen}{G\dsep e\dsep x}
  \DeltaRow{Funktionen}{\circ}
}
\begin{tabproof}
  \proofstep{1}{\GroupAlg{G}{\circ}{e}}{\rA}
  \proofstep{2}{x\in G}{\rA}
  \proofstep{1,2}{x\circ x^{-1}=e}{%
    \FormulaRefAuto{GroupInverseElement}[def]{1,2}}
\end{tabproof}

\FormulaThmDeltaK[Linksinverse Gleichung]{%
  \GroupAlg{G}{\circ}{e}\dsep x\in G\vdash x^{-1}\circ x=e%
}{GroupInverseLeft}{%
  \DeltaRow{Mengen}{G\dsep e\dsep x}
  \DeltaRow{Funktionen}{\circ}
}
\begin{tabproof}
  \proofstep{1}{\GroupAlg{G}{\circ}{e}}{\rA}
  \proofstep{2}{x\in G}{\rA}
  \proofstep{1,2}{x^{-1}\circ x=e}{%
    \FormulaRefAuto{GroupInverseElement}[def]{1,2}}
\end{tabproof}

\section{Kürzungsgesetze}

\FormulaThmDeltaK[Linkskürzung in Gruppen]{%
  \begin{aligned}[t]
    &\GroupAlg{G}{\circ}{e}\dsep a,b,c\in G\dsep
      a\circ b=a\circ c\\
    &\vdash b=c
  \end{aligned}
}{GroupLeftCancellation}{%
  \DeltaRow{Mengen}{G\dsep e\dsep a\dsep b\dsep c}
  \DeltaRow{Funktionen}{\circ}
}
\begin{tabproofsplitwide}
  \proofpartwideindR[Linksmultiplikation mit dem Inversen]{%
    \begin{aligned}[t]
      &\GroupAlg{G}{\circ}{e}\dsep a,b,c\in G\dsep
        a\circ b=a\circ c\vdash{}\\[-2pt]
      &(a^{-1}\circ a)\circ b=(a^{-1}\circ a)\circ c
    \end{aligned}
  }{%
    \begin{aligned}[t]
      &\GroupAlg{G}{\circ}{e}\dsep a,b,c\in G\dsep
        a\circ b=a\circ c\vdash{}\\[-2pt]
      &(a^{-1}\circ a)\circ b=(a^{-1}\circ a)\circ c
    \end{aligned}
  }[GroupLeftCancellationInverseMultiplication]
    \proofstepwidestar[1]{\GroupAlg{G}{\circ}{e}}{\rA}
    \proofstepwidestar[2]{a\in G}{\rA}
    \proofstepwidestar[3]{b\in G}{\rA}
    \proofstepwidestar[4]{c\in G}{\rA}
    \proofstepwidestar[5]{a\circ b=a\circ c}{\rA}
    \proofstepwidestar[1,2]{a^{-1}\in G}{%
      \FormulaRefAuto{GroupInverseCarrier}[thm]{1,2}}
    \proofstepwidestar[5,6]{%
      a^{-1}\circ(a\circ b)=a^{-1}\circ(a\circ c)}{\rIE{5}}
    \proofstepwide[1,2,3,6]{%
      (a^{-1}\circ a)\circ b}{=}{a^{-1}\circ(a\circ b)}{%
      \FormulaRefAuto{GroupAssociativity}[thm]{1,6,2,3}}
    \proofstepwide[1,2,4,6]{%
      a^{-1}\circ(a\circ c)}{=}{(a^{-1}\circ a)\circ c}{%
      \FormulaRefAuto{a=b\vdash b=a}{%
        \FormulaRefAuto{GroupAssociativity}[thm]{1,6,2,4}}}
    \proofstepwidestar[1,2,3,4,5,6]{%
      (a^{-1}\circ a)\circ b=(a^{-1}\circ a)\circ c}{%
      \rChain{8,7,9}}
  \closeproofpartwideindR

  \proofpartwideindR[Reduktion auf das neutrale Element]{%
    \begin{aligned}[t]
      &\GroupAlg{G}{\circ}{e}\dsep a,b,c\in G\dsep
        a\circ b=a\circ c\vdash{}\\[-2pt]
      &e\circ b=e\circ c
    \end{aligned}
  }{%
    \begin{aligned}[t]
      &\GroupAlg{G}{\circ}{e}\dsep a,b,c\in G\dsep
        a\circ b=a\circ c\vdash{}\\[-2pt]
      &e\circ b=e\circ c
    \end{aligned}
  }[GroupLeftCancellationIdentityEquation]
    \proofstepwidestar[1]{\GroupAlg{G}{\circ}{e}}{\rA}
    \proofstepwidestar[2]{a\in G}{\rA}
    \proofstepwidestar[3]{b\in G}{\rA}
    \proofstepwidestar[4]{c\in G}{\rA}
    \proofstepwidestar[5]{a\circ b=a\circ c}{\rA}
    \proofstepwidestar[1,2,3,4,5]{%
      (a^{-1}\circ a)\circ b=(a^{-1}\circ a)\circ c}{%
      \FormulaRefAuto{GroupLeftCancellationInverseMultiplication}[thm]{%
        1,2,3,4,5}}
    \proofstepwidestar[1,2]{a^{-1}\circ a=e}{%
      \FormulaRefAuto{GroupInverseLeft}[thm]{1,2}}
    \proofstepwidestar[1,2]{e=a^{-1}\circ a}{%
      \FormulaRefAuto{a=b\vdash b=a}{7}}
    \proofstepwide[1,2,3]{e\circ b}{=}{(a^{-1}\circ a)\circ b}{\rIE{8}}
    \proofstepwide[1,2,4]{(a^{-1}\circ a)\circ c}{=}{e\circ c}{\rIE{7}}
    \proofstepwidestar[1,2,3,4,5]{e\circ b=e\circ c}{%
      \rChain{9,6,10}}
  \closeproofpartwideindR

  \proofpartwideindR[Entfernen des neutralen Elements]{%
    \begin{aligned}[t]
      &\GroupAlg{G}{\circ}{e}\dsep a,b,c\in G\dsep
        a\circ b=a\circ c\vdash b=c
    \end{aligned}
  }{%
    \begin{aligned}[t]
      &\GroupAlg{G}{\circ}{e}\dsep a,b,c\in G\dsep
        a\circ b=a\circ c\vdash b=c
    \end{aligned}
  }[GroupLeftCancellationConclusion]
    \proofstepwidestar[1]{\GroupAlg{G}{\circ}{e}}{\rA}
    \proofstepwidestar[2]{a\in G}{\rA}
    \proofstepwidestar[3]{b\in G}{\rA}
    \proofstepwidestar[4]{c\in G}{\rA}
    \proofstepwidestar[5]{a\circ b=a\circ c}{\rA}
    \proofstepwidestar[1,2,3,4,5]{e\circ b=e\circ c}{%
      \FormulaRefAuto{GroupLeftCancellationIdentityEquation}[thm]{1,2,3,4,5}}
    \proofstepwide[1,3]{b}{=}{e\circ b}{%
      \FormulaRefAuto{a=b\vdash b=a}{%
        \FormulaRefAuto{GroupLeftIdentity}[thm]{1,3}}}
    \proofstepwide[1,4]{e\circ c}{=}{c}{%
      \FormulaRefAuto{GroupLeftIdentity}[thm]{1,4}}
    \proofstepwidestar[1,2,3,4,5]{b=c}{\rChain{7,6,8}}
  \closeproofpartwideindR
\end{tabproofsplitwide}

\FormulaThmDeltaK[Rechtskürzung in Gruppen]{%
  \begin{aligned}[t]
    &\GroupAlg{G}{\circ}{e}\dsep a,b,c\in G\dsep
      b\circ a=c\circ a\\
    &\vdash b=c
  \end{aligned}
}{GroupRightCancellation}{%
  \DeltaRow{Mengen}{G\dsep e\dsep a\dsep b\dsep c}
  \DeltaRow{Funktionen}{\circ}
}
\begin{tabproofsplitwide}
  \proofpartwideindR[Rechtsmultiplikation mit dem Inversen]{%
    \begin{aligned}[t]
      &\GroupAlg{G}{\circ}{e}\dsep a,b,c\in G\dsep
        b\circ a=c\circ a\vdash{}\\[-2pt]
      &b\circ(a\circ a^{-1})=c\circ(a\circ a^{-1})
    \end{aligned}
  }{%
    \begin{aligned}[t]
      &\GroupAlg{G}{\circ}{e}\dsep a,b,c\in G\dsep
        b\circ a=c\circ a\vdash{}\\[-2pt]
      &b\circ(a\circ a^{-1})=c\circ(a\circ a^{-1})
    \end{aligned}
  }[GroupRightCancellationInverseMultiplication]
    \proofstepwidestar[1]{\GroupAlg{G}{\circ}{e}}{\rA}
    \proofstepwidestar[2]{a\in G}{\rA}
    \proofstepwidestar[3]{b\in G}{\rA}
    \proofstepwidestar[4]{c\in G}{\rA}
    \proofstepwidestar[5]{b\circ a=c\circ a}{\rA}
    \proofstepwidestar[1,2]{a^{-1}\in G}{%
      \FormulaRefAuto{GroupInverseCarrier}[thm]{1,2}}
    \proofstepwidestar[5,6]{%
      (b\circ a)\circ a^{-1}=(c\circ a)\circ a^{-1}}{\rIE{5}}
    \proofstepwide[1,2,3,6]{%
      b\circ(a\circ a^{-1})}{=}{(b\circ a)\circ a^{-1}}{%
      \FormulaRefAuto{a=b\vdash b=a}{%
        \FormulaRefAuto{GroupAssociativity}[thm]{1,3,2,6}}}
    \proofstepwide[1,2,4,6]{%
      (c\circ a)\circ a^{-1}}{=}{c\circ(a\circ a^{-1})}{%
      \FormulaRefAuto{GroupAssociativity}[thm]{1,4,2,6}}
    \proofstepwidestar[1,2,3,4,5,6]{%
      b\circ(a\circ a^{-1})=c\circ(a\circ a^{-1})}{%
      \rChain{8,7,9}}
  \closeproofpartwideindR

  \proofpartwideindR[Reduktion auf das neutrale Element]{%
    \begin{aligned}[t]
      &\GroupAlg{G}{\circ}{e}\dsep a,b,c\in G\dsep
        b\circ a=c\circ a\vdash{}\\[-2pt]
      &b\circ e=c\circ e
    \end{aligned}
  }{%
    \begin{aligned}[t]
      &\GroupAlg{G}{\circ}{e}\dsep a,b,c\in G\dsep
        b\circ a=c\circ a\vdash{}\\[-2pt]
      &b\circ e=c\circ e
    \end{aligned}
  }[GroupRightCancellationIdentityEquation]
    \proofstepwidestar[1]{\GroupAlg{G}{\circ}{e}}{\rA}
    \proofstepwidestar[2]{a\in G}{\rA}
    \proofstepwidestar[3]{b\in G}{\rA}
    \proofstepwidestar[4]{c\in G}{\rA}
    \proofstepwidestar[5]{b\circ a=c\circ a}{\rA}
    \proofstepwidestar[1,2,3,4,5]{%
      b\circ(a\circ a^{-1})=c\circ(a\circ a^{-1})}{%
      \FormulaRefAuto{GroupRightCancellationInverseMultiplication}[thm]{%
        1,2,3,4,5}}
    \proofstepwidestar[1,2]{a\circ a^{-1}=e}{%
      \FormulaRefAuto{GroupInverseRight}[thm]{1,2}}
    \proofstepwidestar[1,2]{e=a\circ a^{-1}}{%
      \FormulaRefAuto{a=b\vdash b=a}{7}}
    \proofstepwide[1,2,3]{b\circ e}{=}{b\circ(a\circ a^{-1})}{\rIE{8}}
    \proofstepwide[1,2,4]{c\circ(a\circ a^{-1})}{=}{c\circ e}{\rIE{7}}
    \proofstepwidestar[1,2,3,4,5]{b\circ e=c\circ e}{%
      \rChain{9,6,10}}
  \closeproofpartwideindR

  \proofpartwideindR[Entfernen des neutralen Elements]{%
    \begin{aligned}[t]
      &\GroupAlg{G}{\circ}{e}\dsep a,b,c\in G\dsep
        b\circ a=c\circ a\vdash b=c
    \end{aligned}
  }{%
    \begin{aligned}[t]
      &\GroupAlg{G}{\circ}{e}\dsep a,b,c\in G\dsep
        b\circ a=c\circ a\vdash b=c
    \end{aligned}
  }[GroupRightCancellationConclusion]
    \proofstepwidestar[1]{\GroupAlg{G}{\circ}{e}}{\rA}
    \proofstepwidestar[2]{a\in G}{\rA}
    \proofstepwidestar[3]{b\in G}{\rA}
    \proofstepwidestar[4]{c\in G}{\rA}
    \proofstepwidestar[5]{b\circ a=c\circ a}{\rA}
    \proofstepwidestar[1,2,3,4,5]{b\circ e=c\circ e}{%
      \FormulaRefAuto{GroupRightCancellationIdentityEquation}[thm]{1,2,3,4,5}}
    \proofstepwide[1,3]{b}{=}{b\circ e}{%
      \FormulaRefAuto{a=b\vdash b=a}{%
        \FormulaRefAuto{GroupRightIdentity}[thm]{1,3}}}
    \proofstepwide[1,4]{c\circ e}{=}{c}{%
      \FormulaRefAuto{GroupRightIdentity}[thm]{1,4}}
    \proofstepwidestar[1,2,3,4,5]{b=c}{\rChain{7,6,8}}
  \closeproofpartwideindR
\end{tabproofsplitwide}

\section{Inverse eines Produkts}

\FormulaThmDeltaK[Inverse eines Produkts]{%
  \begin{aligned}[t]
    &\GroupAlg{G}{\circ}{e}\dsep x,y\in G\vdash{}\\
    &(x\circ y)^{-1}=y^{-1}\circ x^{-1}
  \end{aligned}
}{GroupProductInverse}{%
  \DeltaRow{Mengen}{G\dsep e\dsep x\dsep y}
  \DeltaRow{Funktionen}{\circ}
}
\begin{tabproofsplitwide}
  \proofpartwideindR[Rechtsinverse des Produkts]{%
    \begin{aligned}[t]
      &\GroupAlg{G}{\circ}{e}\dsep x,y\in G\vdash{}\\
      &(x\circ y)\circ(y^{-1}\circ x^{-1})=e
    \end{aligned}
  }{%
    \GroupAlg{G}{\circ}{e}\dsep x,y\in G
    \vdash (x\circ y)\circ(y^{-1}\circ x^{-1})=e%
  }[GroupProductInverseRightCandidate]
    \proofstepwidestar[1]{\GroupAlg{G}{\circ}{e}}{\rA}
    \proofstepwidestar[2]{x\in G}{\rA}
    \proofstepwidestar[3]{y\in G}{\rA}
    \proofstepwidestar[1,2]{x^{-1}\in G}{%
      \FormulaRefAuto{GroupInverseCarrier}[thm]{1,2}}
    \proofstepwidestar[1,3]{y^{-1}\in G}{%
      \FormulaRefAuto{GroupInverseCarrier}[thm]{1,3}}
    \proofstepwide[1,2,3,4,5]{%
      (x\circ y)\circ(y^{-1}\circ x^{-1})}{=}{%
      x\circ\bigl((y\circ y^{-1})\circ x^{-1}\bigr)}{%
      \FormulaRefAuto{GroupAssociativity}[thm]{1,2,3,5,4}}
    \proofstepwide[1,3]{%
      x\circ\bigl((y\circ y^{-1})\circ x^{-1}\bigr)}{=}{%
      x\circ(e\circ x^{-1})}{%
      \FormulaRefAuto{GroupInverseRight}[thm]{1,3}}
    \proofstepwide[1,2,4]{x\circ(e\circ x^{-1})}{=}{x\circ x^{-1}}{%
      \FormulaRefAuto{GroupLeftIdentity}[thm]{1,4}}
    \proofstepwide[1,2]{x\circ x^{-1}}{=}{e}{%
      \FormulaRefAuto{GroupInverseRight}[thm]{1,2}}
    \proofstepwidestar[1,2,3]{%
      (x\circ y)\circ(y^{-1}\circ x^{-1})=e}{%
      \rChain{6,7,8,9}}
  \closeproofpartwideindR

  \proofpartwideindR[Linksinverse des Produkts]{%
    \begin{aligned}[t]
      &\GroupAlg{G}{\circ}{e}\dsep x,y\in G\vdash{}\\
      &(y^{-1}\circ x^{-1})\circ(x\circ y)=e
    \end{aligned}
  }{%
    \GroupAlg{G}{\circ}{e}\dsep x,y\in G
    \vdash (y^{-1}\circ x^{-1})\circ(x\circ y)=e%
  }[GroupProductInverseLeftCandidate]
    \proofstepwidestar[1]{\GroupAlg{G}{\circ}{e}}{\rA}
    \proofstepwidestar[2]{x\in G}{\rA}
    \proofstepwidestar[3]{y\in G}{\rA}
    \proofstepwidestar[1,2]{x^{-1}\in G}{%
      \FormulaRefAuto{GroupInverseCarrier}[thm]{1,2}}
    \proofstepwidestar[1,3]{y^{-1}\in G}{%
      \FormulaRefAuto{GroupInverseCarrier}[thm]{1,3}}
    \proofstepwide[1,2,3,4,5]{%
      (y^{-1}\circ x^{-1})\circ(x\circ y)}{=}{%
      y^{-1}\circ\bigl((x^{-1}\circ x)\circ y\bigr)}{%
      \FormulaRefAuto{GroupAssociativity}[thm]{1,5,4,2,3}}
    \proofstepwide[1,2]{%
      y^{-1}\circ\bigl((x^{-1}\circ x)\circ y\bigr)}{=}{%
      y^{-1}\circ(e\circ y)}{%
      \FormulaRefAuto{GroupInverseLeft}[thm]{1,2}}
    \proofstepwide[1,3,5]{y^{-1}\circ(e\circ y)}{=}{y^{-1}\circ y}{%
      \FormulaRefAuto{GroupLeftIdentity}[thm]{1,3}}
    \proofstepwide[1,3]{y^{-1}\circ y}{=}{e}{%
      \FormulaRefAuto{GroupInverseLeft}[thm]{1,3}}
    \proofstepwidestar[1,2,3]{%
      (y^{-1}\circ x^{-1})\circ(x\circ y)=e}{%
      \rChain{6,7,8,9}}
  \closeproofpartwideindR

  \proofpartwide{Eindeutigkeit}
    \proofstepwidestar[1]{\GroupAlg{G}{\circ}{e}}{\rA}
    \proofstepwidestar[2]{x\in G}{\rA}
    \proofstepwidestar[3]{y\in G}{\rA}
    \proofstepwidestar[1,2]{x^{-1}\in G}{%
      \FormulaRefAuto{GroupInverseCarrier}[thm]{1,2}}
    \proofstepwidestar[1,3]{y^{-1}\in G}{%
      \FormulaRefAuto{GroupInverseCarrier}[thm]{1,3}}
    \proofstepwidestar[1,2,3]{x\circ y\in G}{%
      \FormulaRefAuto{GroupClosure}[thm]{1,2,3}}
    \proofstepwidestar[1,4,5]{y^{-1}\circ x^{-1}\in G}{%
      \FormulaRefAuto{GroupClosure}[thm]{1,5,4}}
    \proofstepwidestar[1,6]{(x\circ y)^{-1}\in G}{%
      \FormulaRefAuto{GroupInverseCarrier}[thm]{1,6}}
    \proofstepwidestar[1,2,3]{%
      (y^{-1}\circ x^{-1})\circ(x\circ y)=e}{%
      \FormulaRefAuto{GroupProductInverseLeftCandidate}[thm]{1,2,3}}
    \proofstepwidestar[1,6]{%
      (x\circ y)\circ(x\circ y)^{-1}=e}{%
      \FormulaRefAuto{GroupInverseRight}[thm]{1,6}}
    \proofstepwidestar[1,2,3,6,7,8,9,10]{%
      y^{-1}\circ x^{-1}=(x\circ y)^{-1}}{%
      \FormulaRefAuto{GroupInverseUnique}[thm]{1,6,7,8,9,10}}
    \proofstepwidestar[1,2,3]{%
      (x\circ y)^{-1}=y^{-1}\circ x^{-1}}{%
      \FormulaRefAuto{a=b\vdash b=a}{11}}
  \closeproofpartwide
\end{tabproofsplitwide}

\chapter{Abelsche Gruppen}

\FormulaAxiomDeltaK[Kommutativität einer Gruppenoperation]{%
  \GroupAlg{G}{\circ}{e}\dsep x,y\in G
  \vdash x\circ y=y\circ x%
}{AbelianGroupCommutativityLaw}{%
  \DeltaRow{Mengen}{G\dsep e\dsep x\dsep y}
  \DeltaRow{Funktionen}{\circ}
}

\FormulaDefDeltaK[Begriff der abelschen Gruppe]{%
  \AbelianGroupAlg{G}{\circ}{e}%
}{AbelianGroupDef}{%
  \DeltaRow{Mengen}{G\dsep e\dsep x\dsep y}
  \DeltaRow{Funktionen}{\circ}
  \DeltaRow{\textbf{Axiome}}{}
  \DeltaRow{Gruppe}{\GroupAlg{G}{\circ}{e}}
    [\FormulaRefAuto{GroupDef}[def]]
  \DeltaRow{Kommutativität}{%
    x,y\in G\vdash x\circ y=y\circ x}
    [\FormulaRefAuto{AbelianGroupCommutativityLaw}[ax]]
  \DeltaRow{\textbf{Neues Symbol}}{}
  \DeltaPrem{Abelsche Gruppe}{\AbelianGroupAlg{G}{\circ}{e}}
}

\FormulaAxiomDeltaK[Gruppenanteil einer abelschen Gruppe]{%
  \AbelianGroupAlg{G}{\circ}{e}\vdash\GroupAlg{G}{\circ}{e}%
}{AbelianGroupBaseGroupAxiom}{%
  \DeltaRow{Mengen}{G\dsep e}
  \DeltaRow{Funktionen}{\circ}
}

\FormulaAxiomDeltaK[Kommutativität in einer abelschen Gruppe]{%
  \AbelianGroupAlg{G}{\circ}{e}\dsep x,y\in G
  \vdash x\circ y=y\circ x%
}{AbelianGroupCommutativityAxiom}{%
  \DeltaRow{Mengen}{G\dsep e\dsep x\dsep y}
  \DeltaRow{Funktionen}{\circ}
}

\FormulaThmDeltaK[Abelsche Gruppen sind kommutative Monoide]{%
  \AbelianGroupAlg{G}{\circ}{e}
  \vdash\CommutativeMonoidAlg{G}{\circ}{e}%
}{AbelianGroupIsCommutativeMonoid}{%
  \DeltaRow{Mengen}{G\dsep e\dsep x\dsep y}
  \DeltaRow{Funktionen}{\circ}
}
\begin{tabproof}
  \proofstep{1}{\AbelianGroupAlg{G}{\circ}{e}}{\rA}
  \proofstep{1}{\GroupAlg{G}{\circ}{e}}{%
    \FormulaRefAuto{AbelianGroupBaseGroupAxiom}[ax]{1}}
  \proofstep{1}{\MonoidAlg{G}{\circ}{e}}{%
    \FormulaRefAuto{GroupBaseMonoidAxiom}[ax]{2}}
  \proofstep{4}{x\in G}{\rA}
  \proofstep{5}{y\in G}{\rA}
  \proofstep{1,4,5}{x\circ y=y\circ x}{%
    \FormulaRefAuto{AbelianGroupCommutativityAxiom}[ax]{1,4,5}}
  \proofstep{1}{\CommutativeMonoidAlg{G}{\circ}{e}}{%
    \FormulaRefAuto{CommutativeMonoidDef}[def]{3,6}}
\end{tabproof}

\chapter{Homomorphismen von Gruppen}

\section{Gruppenhomomorphismen}

Bei Gruppen genügt die Erhaltung der Gruppenoperation. Anders als bei
beliebigen Monoiden folgen die Erhaltung des neutralen Elements und der
Inversen aus der Existenz von Inversen im Ziel.

\FormulaDefDeltaK[Begriff des Gruppenhomomorphismus]{%
  \GroupHom{f}{G,\circ,e}{H,\diamond,u}%
}{GroupHomDef}{%
  \DeltaRow{Mengen}{G\dsep H\dsep e\dsep u\dsep x\dsep y}
  \DeltaRow{Funktionen}{f\dsep\circ\dsep\diamond}
  \DeltaRow{\textbf{Axiome}}{}
  \DeltaPrem{Quellgruppe}{\GroupAlg{G}{\circ}{e}}
  \DeltaPrem{Zielgruppe}{\GroupAlg{H}{\diamond}{u}}
  \DeltaPrem{Halbgruppenhomomorphismus}{%
    \SemigroupHom{f}{G,\circ}{H,\diamond}}
  \DeltaRow{\textbf{Neues Symbol}}{}
  \DeltaPrem{Gruppenhomomorphismus}{%
    \GroupHom{f}{G,\circ,e}{H,\diamond,u}}
}

\FormulaAxiomDeltaK[Trägerabbildung eines Gruppenhomomorphismus]{%
  \GroupHom{f}{G,\circ,e}{H,\diamond,u}
  \vdash f\colon G\to H%
}{GroupHomFunctionAxiom}{%
  \DeltaRow{Mengen}{G\dsep H\dsep e\dsep u}
  \DeltaRow{Funktionen}{f\dsep\circ\dsep\diamond}
}
\par\smallskip

\FormulaAxiomDeltaK[Operationserhaltung eines Gruppenhomomorphismus]{%
  \GroupHom{f}{G,\circ,e}{H,\diamond,u}\dsep x,y\in G
  \vdash f(x\circ y)=f(x)\diamond f(y)%
}{GroupHomOperationAxiom}{%
  \DeltaRow{Mengen}{G\dsep H\dsep e\dsep u\dsep x\dsep y}
  \DeltaRow{Funktionen}{f\dsep\circ\dsep\diamond}
}

\FormulaThmDeltaK[Gruppenhomomorphismen erhalten das neutrale Element]{%
  \GroupHom{f}{G,\circ,e}{H,\diamond,u}
  \vdash f(e)=u%
}{GroupHomIdentityPreservation}{%
  \DeltaRow{Mengen}{G\dsep H\dsep e\dsep u}
  \DeltaRow{Funktionen}{f\dsep\circ\dsep\diamond}
}
\begin{tabproofwide}
  \proofstepwidestar[1]{%
    \GroupHom{f}{G,\circ,e}{H,\diamond,u}}{\rA}
  \proofstepwidestar[1]{\GroupAlg{G}{\circ}{e}}{%
    \FormulaRefAuto{GroupHomDef}[def]{1}}
  \proofstepwidestar[1]{\GroupAlg{H}{\diamond}{u}}{%
    \FormulaRefAuto{GroupHomDef}[def]{1}}
  \proofstepwidestar[1]{f\colon G\to H}{%
    \FormulaRefAuto{GroupHomFunctionAxiom}[ax]{1}}
  \proofstepwidestar[1]{e\in G}{%
    \FormulaRefAuto{GroupIdentityCarrier}[thm]{2}}
  \proofstepwidestar[1]{f(e)\in H}{%
    \FormulaRefAuto{F\colon A\to B\dsep a\in A\vdash F(a)\in B}{4,5}}
  \proofstepwidestar[1]{e\circ e=e}{%
    \FormulaRefAuto{GroupLeftIdentity}[thm]{2,5}}
  \proofstepwidestar[1]{f(e\circ e)=f(e)}{\rIE{7}}
  \proofstepwidestar[1]{f(e\circ e)=f(e)\diamond f(e)}{%
    \FormulaRefAuto{GroupHomOperationAxiom}[ax]{1,5,5}}
  \proofstepwidestar[1]{f(e)\diamond f(e)=f(e)}{\rIE{9,8}}
  \proofstepwidestar[1]{f(e)\diamond u=f(e)}{%
    \FormulaRefAuto{GroupRightIdentity}[thm]{3,6}}
  \proofstepwidestar[1]{%
    f(e)\diamond u=f(e)\diamond f(e)}{\rIE{11,%
      \FormulaRefAuto{a=b\vdash b=a}{10}}}
  \proofstepwidestar[1]{u=f(e)}{%
    \FormulaRefAuto{GroupLeftCancellation}[thm]{3,6,
      \FormulaRefAuto{GroupIdentityCarrier}[thm]{3},6,12}}
  \proofstepwidestar[1]{f(e)=u}{%
    \FormulaRefAuto{a=b\vdash b=a}{13}}
\end{tabproofwide}

\FormulaThmDeltaK[Gruppenhomomorphismen erhalten Inverse]{%
  \GroupHom{f}{G,\circ,e}{H,\diamond,u}\dsep x\in G
  \vdash f(x^{-1})=f(x)^{-1}%
}{GroupHomInversePreservation}{%
  \DeltaRow{Mengen}{G\dsep H\dsep e\dsep u\dsep x}
  \DeltaRow{Funktionen}{f\dsep\circ\dsep\diamond}
}
\begin{tabproofwide}
  \proofstepwidestar[1]{%
    \GroupHom{f}{G,\circ,e}{H,\diamond,u}}{\rA}
  \proofstepwidestar[2]{x\in G}{\rA}
  \proofstepwidestar[1]{\GroupAlg{G}{\circ}{e}}{%
    \FormulaRefAuto{GroupHomDef}[def]{1}}
  \proofstepwidestar[1]{\GroupAlg{H}{\diamond}{u}}{%
    \FormulaRefAuto{GroupHomDef}[def]{1}}
  \proofstepwidestar[1]{f\colon G\to H}{%
    \FormulaRefAuto{GroupHomFunctionAxiom}[ax]{1}}
  \proofstepwidestar[1,2]{x^{-1}\in G}{%
    \FormulaRefAuto{GroupInverseCarrier}[thm]{3,2}}
  \proofstepwidestar[1,2]{f(x)\in H}{%
    \FormulaRefAuto{F\colon A\to B\dsep a\in A\vdash F(a)\in B}{5,2}}
  \proofstepwidestar[1,2]{f(x^{-1})\in H}{%
    \FormulaRefAuto{F\colon A\to B\dsep a\in A\vdash F(a)\in B}{5,6}}
  \proofstepwidestar[1,2]{%
    f(x)\diamond f(x^{-1})=f(x\circ x^{-1})}{%
    \FormulaRefAuto{a=b\vdash b=a}{%
      \FormulaRefAuto{GroupHomOperationAxiom}[ax]{1,2,6}}}
  \proofstepwidestar[1,2]{f(x\circ x^{-1})=f(e)}{%
    \FormulaRefAuto{GroupInverseRight}[thm]{3,2}}
  \proofstepwidestar[1]{f(e)=u}{%
    \FormulaRefAuto{GroupHomIdentityPreservation}[thm]{1}}
  \proofstepwidestar[1,2]{f(x)\diamond f(x^{-1})=u}{%
    \rChain{9,10,11}}
  \proofstepwidestar[1,2]{%
    f(x^{-1})\diamond f(x)=f(x^{-1}\circ x)}{%
    \FormulaRefAuto{a=b\vdash b=a}{%
      \FormulaRefAuto{GroupHomOperationAxiom}[ax]{1,6,2}}}
  \proofstepwidestar[1,2]{f(x^{-1}\circ x)=f(e)}{%
    \FormulaRefAuto{GroupInverseLeft}[thm]{3,2}}
  \proofstepwidestar[1,2]{f(x^{-1})\diamond f(x)=u}{%
    \rChain{13,14,11}}
  \proofstepwidestar[1,2]{f(x)^{-1}\in H}{%
    \FormulaRefAuto{GroupInverseCarrier}[thm]{4,7}}
  \proofstepwidestar[1,2]{%
    f(x)\diamond f(x)^{-1}=u}{%
    \FormulaRefAuto{GroupInverseRight}[thm]{4,7}}
  \proofstepwidestar[1,2]{%
    f(x^{-1})=f(x)^{-1}}{%
    \FormulaRefAuto{GroupInverseUnique}[thm]{4,7,8,16,15,17}}
\end{tabproofwide}

\FormulaThmDeltaK[Ein Gruppenhomomorphismus ist ein Monoidhomomorphismus]{%
  \GroupHom{f}{G,\circ,e}{H,\diamond,u}
  \vdash\MonoidHom{f}{G,\circ,e}{H,\diamond,u}%
}{GroupHomIsMonoidHom}{%
  \DeltaRow{Mengen}{G\dsep H\dsep e\dsep u}
  \DeltaRow{Funktionen}{f\dsep\circ\dsep\diamond}
}
\begin{tabproof}
  \proofstep{1}{\GroupHom{f}{G,\circ,e}{H,\diamond,u}}{\rA}
  \proofstep{1}{\GroupAlg{G}{\circ}{e}}{%
    \FormulaRefAuto{GroupHomDef}[def]{1}}
  \proofstep{1}{\GroupAlg{H}{\diamond}{u}}{%
    \FormulaRefAuto{GroupHomDef}[def]{1}}
  \proofstep{1}{\MonoidAlg{G}{\circ}{e}}{%
    \FormulaRefAuto{GroupBaseMonoidAxiom}[ax]{2}}
  \proofstep{1}{\MonoidAlg{H}{\diamond}{u}}{%
    \FormulaRefAuto{GroupBaseMonoidAxiom}[ax]{3}}
  \proofstep{1}{\SemigroupHom{f}{G,\circ}{H,\diamond}}{%
    \FormulaRefAuto{GroupHomDef}[def]{1}}
  \proofstep{1}{f(e)=u}{%
    \FormulaRefAuto{GroupHomIdentityPreservation}[thm]{1}}
  \proofstep{1}{\MonoidHom{f}{G,\circ,e}{H,\diamond,u}}{%
    \FormulaRefAuto{MonoidHomDef}[def]{4,5,6,7}}
\end{tabproof}

\section{Gruppenisomorphismen und Komposition}

\FormulaDefDeltaK[Begriff des Gruppenisomorphismus]{%
  \GroupIso{f}{G,\circ,e}{H,\diamond,u}%
}{GroupIsoDef}{%
  \DeltaRow{Mengen}{G\dsep H\dsep e\dsep u}
  \DeltaRow{Funktionen}{f\dsep\circ\dsep\diamond}
  \DeltaRow{\textbf{Axiome}}{}
  \DeltaPrem{Gruppenhomomorphismus}{%
    \GroupHom{f}{G,\circ,e}{H,\diamond,u}}
  \DeltaPrem{Bijektivität}{f\colon G\bij H}
  \DeltaRow{\textbf{Neues Symbol}}{}
  \DeltaPrem{Gruppenisomorphismus}{%
    \GroupIso{f}{G,\circ,e}{H,\diamond,u}}
}

\FormulaThmDeltaK[Komposition von Gruppenhomomorphismen]{%
  \begin{aligned}[t]
    &\GroupHom{f}{G,\circ,e}{H,\diamond,u}\dsep
      \GroupHom{g}{H,\diamond,u}{K,\bullet,v}\\
    &\vdash\GroupHom{g\circ f}{G,\circ,e}{K,\bullet,v}
  \end{aligned}
}{GroupHomComposition}{%
  \DeltaRow{Mengen}{G\dsep H\dsep K\dsep e\dsep u\dsep v}
  \DeltaRow{Funktionen}{f\dsep g\dsep\circ\dsep\diamond\dsep\bullet}
}
\begin{tabproof}
  \proofstep{1}{\GroupHom{f}{G,\circ,e}{H,\diamond,u}}{\rA}
  \proofstep{2}{\GroupHom{g}{H,\diamond,u}{K,\bullet,v}}{\rA}
  \proofstep{1}{\SemigroupHom{f}{G,\circ}{H,\diamond}}{%
    \FormulaRefAuto{GroupHomDef}[def]{1}}
  \proofstep{2}{\SemigroupHom{g}{H,\diamond}{K,\bullet}}{%
    \FormulaRefAuto{GroupHomDef}[def]{2}}
  \proofstep{1,2}{\SemigroupHom{g\circ f}{G,\circ}{K,\bullet}}{%
    \FormulaRefAuto{SemigroupHomComposition}[thm]{3,4}}
  \proofstep{1}{\GroupAlg{G}{\circ}{e}}{%
    \FormulaRefAuto{GroupHomDef}[def]{1}}
  \proofstep{2}{\GroupAlg{K}{\bullet}{v}}{%
    \FormulaRefAuto{GroupHomDef}[def]{2}}
  \proofstep{1,2}{\GroupHom{g\circ f}{G,\circ,e}{K,\bullet,v}}{%
    \FormulaRefAuto{GroupHomDef}[def]{6,7,5}}
\end{tabproof}

\FormulaThmDeltaK[Die Identität ist ein Gruppenisomorphismus]{%
  \GroupAlg{G}{\circ}{e}
  \vdash\GroupIso{\Id_G}{G,\circ,e}{G,\circ,e}%
}{GroupIdentityIsomorphism}{%
  \DeltaRow{Mengen}{G\dsep e}
  \DeltaRow{Funktionen}{\circ}
}
\begin{tabproof}
  \proofstep{1}{\GroupAlg{G}{\circ}{e}}{\rA}
  \proofstep{1}{\SemigroupAlg{G}{\circ}}{%
    \FormulaRefAuto{GroupBaseSemigroup}[thm]{1}}
  \proofstep{1}{\SemigroupIso{\Id_G}{G,\circ}{G,\circ}}{%
    \FormulaRefAuto{SemigroupIdentityIsomorphism}[thm]{2}}
  \proofstep{1}{\SemigroupHom{\Id_G}{G,\circ}{G,\circ}}{%
    \FormulaRefAuto{SemigroupIsoHomAxiom}[ax]{3}}
  \proofstep{1}{\GroupHom{\Id_G}{G,\circ,e}{G,\circ,e}}{%
    \FormulaRefAuto{GroupHomDef}[def]{1,1,4}}
  \proofstep{}{\Id_G\colon G\bij G}{%
    \FormulaRefAuto{\Id_A\colon A\bij A}[thm]}
  \proofstep{1}{\GroupIso{\Id_G}{G,\circ,e}{G,\circ,e}}{%
    \FormulaRefAuto{GroupIsoDef}[def]{5,6}}
\end{tabproof}

\section{Spezielle Definitionen}

Für abelsche Gruppen bleibt die Erhaltungsgleichung unverändert. Die
Kommutativität ist eine Eigenschaft von Quell- und Zielgruppe, keine weitere
Forderung an die Abbildung.

\FormulaDefDeltaK[Homomorphismen abelscher Gruppen]{%
  \begin{aligned}[t]
    &\AbelianGroupHom{f}{G,\circ,e}{H,\diamond,u}\\[-2pt]
    &\quad\coloneqq
      \AbelianGroupAlg{G}{\circ}{e}\land
      \AbelianGroupAlg{H}{\diamond}{u}\\[-2pt]
    &\qquad\land\GroupHom{f}{G,\circ,e}{H,\diamond,u}
  \end{aligned}
}{AbelianGroupHomDef}{%
  \DeltaRow{Mengen}{G\dsep H\dsep e\dsep u}
  \DeltaRow{Funktionen}{f\dsep\circ\dsep\diamond}
  \DeltaRow{Erhaltungsgesetz}{%
    x,y\in G\vdash f(x\circ y)=f(x)\diamond f(y)}
  \DeltaRow{\textbf{Neues Symbol}}{}
  \DeltaPrem{Symbol}{\mathsf{AGrpHom}}
}

\FormulaAxiomDeltaK[Quellstruktur eines Homomorphismus abelscher Gruppen]{%
  \AbelianGroupHom{f}{G,\circ,e}{H,\diamond,u}
  \vdash \AbelianGroupAlg{G}{\circ}{e}%
}{AbelianGroupHomSourceAxiom}{%
  \DeltaRow{Mengen}{G\dsep H\dsep e\dsep u}
  \DeltaRow{Funktionen}{f\dsep\circ\dsep\diamond}
}

\FormulaAxiomDeltaK[Zielstruktur eines Homomorphismus abelscher Gruppen]{%
  \AbelianGroupHom{f}{G,\circ,e}{H,\diamond,u}
  \vdash \AbelianGroupAlg{H}{\diamond}{u}%
}{AbelianGroupHomTargetAxiom}{%
  \DeltaRow{Mengen}{G\dsep H\dsep e\dsep u}
  \DeltaRow{Funktionen}{f\dsep\circ\dsep\diamond}
}

\FormulaAxiomDeltaK[Gruppenhomomorphismusanteil eines Homomorphismus abelscher Gruppen]{%
  \AbelianGroupHom{f}{G,\circ,e}{H,\diamond,u}
  \vdash \GroupHom{f}{G,\circ,e}{H,\diamond,u}%
}{AbelianGroupHomBaseHomAxiom}{%
  \DeltaRow{Mengen}{G\dsep H\dsep e\dsep u}
  \DeltaRow{Funktionen}{f\dsep\circ\dsep\diamond}
}

\FormulaDefDeltaK[Isomorphismen abelscher Gruppen]{%
  \begin{aligned}[t]
    &\AbelianGroupIso{f}{G,\circ,e}{H,\diamond,u}\\[-2pt]
    &\quad\coloneqq
      \AbelianGroupHom{f}{G,\circ,e}{H,\diamond,u}
      \land f\colon G\bij H
  \end{aligned}
}{AbelianGroupIsoDef}{%
  \DeltaRow{Mengen}{G\dsep H\dsep e\dsep u}
  \DeltaRow{Funktionen}{f\dsep\circ\dsep\diamond}
  \DeltaRow{\textbf{Neues Symbol}}{}
  \DeltaPrem{Symbol}{\mathsf{AGrpIso}}
}

\FormulaAxiomDeltaK[Homomorphismusanteil eines Isomorphismus abelscher Gruppen]{%
  \AbelianGroupIso{f}{G,\circ,e}{H,\diamond,u}
  \vdash \AbelianGroupHom{f}{G,\circ,e}{H,\diamond,u}%
}{AbelianGroupIsoHomAxiom}{%
  \DeltaRow{Mengen}{G\dsep H\dsep e\dsep u}
  \DeltaRow{Funktionen}{f\dsep\circ\dsep\diamond}
}

\FormulaAxiomDeltaK[Bijektivität eines Isomorphismus abelscher Gruppen]{%
  \AbelianGroupIso{f}{G,\circ,e}{H,\diamond,u}
  \vdash f\colon G\bij H%
}{AbelianGroupIsoBijectiveAxiom}{%
  \DeltaRow{Mengen}{G\dsep H\dsep e\dsep u}
  \DeltaRow{Funktionen}{f\dsep\circ\dsep\diamond}
}

Ein Gruppenisomorphismus transportiert die Kommutativität in beiden
Richtungen. Deshalb stimmt diese Spezialisierung mit der üblichen
Isomorphie abelscher Gruppen überein.

\chapter{Formale Differenzen und der Einbettungssatz}

\section{Differenzenpaare}

Sei im Folgenden \(\CommutativeMonoidAlg{M}{\star}{e}\), und sei
\(\star\) kürzbar. Die Konstruktion aus Band~17 wird nun mit der
vollständigen Ausgangsstruktur \((M,\star,e)\) im Index wiederholt.

\FormulaDefDeltaK[Träger der formalen Differenzenpaare]{%
  \CommutativeMonoidAlg{M}{\star}{e}\vdash
  \FormalDifferenceCarrier{M,\star,e}\coloneqq M\times M
}{GroupCompletionPairCarrier}{%
  \DeltaRow{Mengen}{M\dsep e\dsep\FormalDifferenceCarrier{M,\star,e}}
  \DeltaRow{Funktionen}{\star}
  \DeltaRow{Neue Symbole}{\FormalDifferenceCarrier{M,\star,e}}
}

\FormulaDefDeltaK[Äquivalenz formaler Differenzenpaare]{%
  \begin{aligned}[t]
    &\CommutativeMonoidAlg{M}{\star}{e}\dsep
      a,b,c,d\in M\vdash{}\\[-2pt]
    &(a,b)\FormalDifferenceRel{M,\star,e}(c,d)
      \coloneqq a\star d=b\star c
  \end{aligned}
}{GroupCompletionPairRelationDef}{%
  \DeltaRow{Mengen}{M\dsep e\dsep a\dsep b\dsep c\dsep d}
  \DeltaRow{Funktionen}{\star}
  \DeltaRow{Relationsoperatoren}{\FormalDifferenceRel{M,\star,e}}
  \DeltaRow{Trägermenge}{\FormalDifferenceCarrier{M,\star,e}}
  \DeltaRow{Neue Symbole}{\FormalDifferenceRel{M,\star,e}}
}

\FormulaThmDeltaK[Charakterisierung der Relation formaler Differenzen]{%
  \begin{aligned}[t]
    &\CommutativeMonoidAlg{M}{\star}{e}\dsep
      a,b,c,d\in M\vdash{}\\[-2pt]
    &(a,b)\FormalDifferenceRel{M,\star,e}(c,d)
      \leftrightarrow a\star d=b\star c
  \end{aligned}
}{GroupCompletionPairRelationChar}{%
  \DeltaRow{Mengen}{M\dsep e\dsep a\dsep b\dsep c\dsep d}
  \DeltaRow{Funktionen}{\star}
  \DeltaRow{Relationsoperatoren}{\FormalDifferenceRel{M,\star,e}}
}
\begin{tabproof}
  \proofstep{1}{\CommutativeMonoidAlg{M}{\star}{e}}{\rA}
  \proofstep{2}{a,b,c,d\in M}{\rA}
  \proofstep{1,2}{%
    (a,b)\FormalDifferenceRel{M,\star,e}(c,d)
    \leftrightarrow a\star d=b\star c}{%
    \FormulaRefAuto{GroupCompletionPairRelationDef}[def]{1,2}}
\end{tabproof}

\FormulaThmDeltaK[Reflexivität auf formalen Differenzenpaaren]{%
  \begin{aligned}[t]
    &\CommutativeMonoidAlg{M}{\star}{e}\dsep a,b\in M\vdash{}\\[-2pt]
    &(a,b)\FormalDifferenceRel{M,\star,e}(a,b)
  \end{aligned}
}{GroupCompletionPairRelationReflexive}{%
  \DeltaRow{Mengen}{M\dsep e\dsep a\dsep b}
  \DeltaRow{Funktionen}{\star}
  \DeltaRow{Relationsoperatoren}{\FormalDifferenceRel{M,\star,e}}
}
\begin{tabproof}
  \proofstep{1}{\CommutativeMonoidAlg{M}{\star}{e}}{\rA}
  \proofstep{2}{a,b\in M}{\rA}
  \proofstep{1,2}{a\star b=b\star a}{%
    \FormulaRefAuto{CommutativeMonoidCommutativityAxiom}[ax]{1,2}}
  \proofstep{1,2}{(a,b)\FormalDifferenceRel{M,\star,e}(a,b)}{%
    \FormulaRefAuto{GroupCompletionPairRelationChar}{1,2,3}}
\end{tabproof}

\FormulaThmDeltaK[Symmetrie auf formalen Differenzenpaaren]{%
  \begin{aligned}[t]
    &\CommutativeMonoidAlg{M}{\star}{e}\dsep a,b,c,d\in M\dsep
      (a,b)\FormalDifferenceRel{M,\star,e}(c,d)\vdash{}\\[-2pt]
    &(c,d)\FormalDifferenceRel{M,\star,e}(a,b)
  \end{aligned}
}{GroupCompletionPairRelationSymmetric}{%
  \DeltaRow{Mengen}{M\dsep e\dsep a\dsep b\dsep c\dsep d}
  \DeltaRow{Funktionen}{\star}
  \DeltaRow{Relationsoperatoren}{\FormalDifferenceRel{M,\star,e}}
}
\begin{tabproofwide}
  \proofstepwidestar[1]{\CommutativeMonoidAlg{M}{\star}{e}}{\rA}
  \proofstepwidestar[2]{a,b,c,d\in M}{\rA}
  \proofstepwidestar[3]{%
    (a,b)\FormalDifferenceRel{M,\star,e}(c,d)}{\rA}
  \proofstepwide[1,2,3]{c\star b}{=}{b\star c}{%
    \FormulaRefAuto{CommutativeMonoidCommutativityAxiom}[ax]{1,2}}
  \proofstepwide[1,2,3]{}{=}{a\star d}{%
    \FormulaRefAuto{a=b\vdash b=a}{%
      \FormulaRefAuto{GroupCompletionPairRelationChar}{1,2,3}}}
  \proofstepwide[1,2,3]{}{=}{d\star a}{%
    \FormulaRefAuto{CommutativeMonoidCommutativityAxiom}[ax]{1,2}}
  \proofstepwidestar[1,2,3]{c\star b=d\star a}{\rChain{4,6}}
  \proofstepwidestar[1,2,3]{%
    (c,d)\FormalDifferenceRel{M,\star,e}(a,b)}{%
    \FormulaRefAuto{GroupCompletionPairRelationChar}{1,2,7}}
\end{tabproofwide}

\FormulaThmDeltaK[Transitiver Kreuzproduktvergleich]{%
  \begin{aligned}[t]
    &\CommutativeMonoidAlg{M}{\star}{e}\dsep
      a,b,c,d,f,g\in M\dsep
      a\star d=b\star c\dsep c\star g=d\star f\vdash{}\\[-2pt]
    &(a\star g)\star(c\star d)
      =(b\star f)\star(c\star d)
  \end{aligned}
}{GroupCompletionPairRelationTransitiveRearrangement}{%
  \DeltaRow{Mengen}{M\dsep e\dsep a\dsep b\dsep c\dsep d\dsep f\dsep g}
  \DeltaRow{Funktionen}{\star}
}
\begin{tabproofwide}
  \proofstepwidestar[1]{\CommutativeMonoidAlg{M}{\star}{e}}{\rA}
  \proofstepwidestar[2]{a,b,c,d,f,g\in M}{\rA}
  \proofstepwidestar[3]{a\star d=b\star c}{\rA}
  \proofstepwidestar[4]{c\star g=d\star f}{\rA}
  \proofstepwidestar[1]{\CommutativeSemigroupAlg{M}{\star}}{%
    \FormulaRefAuto{CommutativeMonoidIsCommutativeSemigroup}{1}}
  \proofstepwide[1,2,3,4]{(a\star g)\star(c\star d)}{=}{%
    (a\star g)\star(d\star c)}{%
    \FormulaRefAuto{CommutativeMonoidCommutativityAxiom}[ax]{1,2}}
  \proofstepwide[1,2,3,4]{}{=}{(a\star d)\star(g\star c)}{%
    \FormulaRefAuto{CommutativeSemigroupFourTermExchange}{5,2}}
  \proofstepwide[1,2,3,4]{}{=}{(a\star d)\star(c\star g)}{%
    \FormulaRefAuto{CommutativeMonoidCommutativityAxiom}[ax]{1,2}}
  \proofstepwide[1,2,3,4]{}{=}{(b\star c)\star(d\star f)}{\rIE{3,4}}
  \proofstepwide[1,2,3,4]{}{=}{(b\star c)\star(f\star d)}{%
    \FormulaRefAuto{CommutativeMonoidCommutativityAxiom}[ax]{1,2}}
  \proofstepwide[1,2,3,4]{}{=}{(b\star f)\star(c\star d)}{%
    \FormulaRefAuto{CommutativeSemigroupFourTermExchange}{5,2}}
  \proofstepwidestar[1,2,3,4]{%
    (a\star g)\star(c\star d)
      =(b\star f)\star(c\star d)}{\rChain{6,11}}
\end{tabproofwide}

\FormulaThmDeltaK[Transitivität auf formalen Differenzenpaaren]{%
  \begin{aligned}[t]
    &\CommutativeMonoidAlg{M}{\star}{e}\dsep
      \mathsf{K\ddot{u}rz}(M,\star)\dsep a,b,c,d,f,g\in M\dsep{}\\[-2pt]
    &(a,b)\FormalDifferenceRel{M,\star,e}(c,d)\dsep
      (c,d)\FormalDifferenceRel{M,\star,e}(f,g)\vdash
      (a,b)\FormalDifferenceRel{M,\star,e}(f,g)
  \end{aligned}
}{GroupCompletionPairRelationTransitive}{%
  \DeltaRow{Mengen}{M\dsep e\dsep a\dsep b\dsep c\dsep d\dsep f\dsep g}
  \DeltaRow{Funktionen}{\star}
  \DeltaRow{Relationsoperatoren}{\FormalDifferenceRel{M,\star,e}}
}
\begin{tabproofwide}
  \proofstepwidestar[1]{\CommutativeMonoidAlg{M}{\star}{e}}{\rA}
  \proofstepwidestar[2]{\mathsf{K\ddot{u}rz}(M,\star)}{\rA}
  \proofstepwidestar[3]{a,b,c,d,f,g\in M}{\rA}
  \proofstepwidestar[4]{%
    (a,b)\FormalDifferenceRel{M,\star,e}(c,d)}{\rA}
  \proofstepwidestar[5]{%
    (c,d)\FormalDifferenceRel{M,\star,e}(f,g)}{\rA}
  \proofstepwidestar[1,3,4]{a\star d=b\star c}{%
    \FormulaRefAuto{GroupCompletionPairRelationChar}{1,3,4}}
  \proofstepwidestar[1,3,5]{c\star g=d\star f}{%
    \FormulaRefAuto{GroupCompletionPairRelationChar}{1,3,5}}
  \proofstepwidestar[1,3,4,5]{%
    (a\star g)\star(c\star d)
      =(b\star f)\star(c\star d)}{%
    \FormulaRefAuto{GroupCompletionPairRelationTransitiveRearrangement}{%
      1,3,6,7}}
  \proofstepwidestar[1]{\MonoidAlg{M}{\star}{e}}{%
    \FormulaRefAuto{CommutativeMonoidBaseMonoidAxiom}[ax]{1}}
  \proofstepwidestar[1]{\SemigroupAlg{M}{\star}}{%
    \FormulaRefAuto{MonoidBaseSemigroupAxiom}[ax]{9}}
  \proofstepwidestar[2]{\mathsf{RK\ddot{u}rz}(M,\star)}{%
    \FormulaRefAuto{CancellativeRightAxiom}[ax]{2}}
  \proofstepwidestar[1,3]{a\star g\in M}{%
    \FormulaRefAuto{SemigroupClosureAxiom}[ax]{10,3}}
  \proofstepwidestar[1,3]{b\star f\in M}{%
    \FormulaRefAuto{SemigroupClosureAxiom}[ax]{10,3}}
  \proofstepwidestar[1,3]{c\star d\in M}{%
    \FormulaRefAuto{SemigroupClosureAxiom}[ax]{10,3}}
  \proofstepwidestar[1,2,3,4,5]{a\star g=b\star f}{%
    \FormulaRefAuto{RightCancellationAxiom}[ax]{11,12,13,14,8}}
  \proofstepwidestar[1,2,3,4,5]{%
    (a,b)\FormalDifferenceRel{M,\star,e}(f,g)}{%
    \FormulaRefAuto{GroupCompletionPairRelationChar}{1,3,15}}
\end{tabproofwide}

\FormulaThmDeltaK[Darstellung eines formalen Differenzenpaares]{%
  \begin{aligned}[t]
    &\CommutativeMonoidAlg{M}{\star}{e}\dsep
      x\in\FormalDifferenceCarrier{M,\star,e}\vdash{}\\[-2pt]
    &\exists a\in M\,\exists b\in M\,x=(a,b)
  \end{aligned}
}{GroupCompletionPairRepresentative}{%
  \DeltaRow{Mengen}{M\dsep e\dsep x\dsep a\dsep b}
  \DeltaRow{Funktionen}{\star}
}
\begin{tabproofwide}
  \proofstepwidestar[1]{\CommutativeMonoidAlg{M}{\star}{e}}{\rA}
  \proofstepwidestar[2]{x\in\FormalDifferenceCarrier{M,\star,e}}{\rA}
  \proofstepwidestar[1]{\FormalDifferenceCarrier{M,\star,e}=M\times M}{%
    \FormulaRefAuto{GroupCompletionPairCarrier}[def]{1}}
  \proofstepwidestar[1,2]{x\in M\times M}{\rIE{3,2}}
  \proofstepwidestar[1,2]{\exists a\in M\,\exists b\in M\,x=(a,b)}{%
    \FormulaRefAuto{x\in A\times B\eqvdash
      \exists a\in A\,\exists b\in B\,x=(a,b)}{4}}
\end{tabproofwide}

\FormulaThmDeltaK[Die Relation formaler Differenzen ist eine Äquivalenzrelation]{%
  \begin{aligned}[t]
    &\CommutativeMonoidAlg{M}{\star}{e}\dsep
      \mathsf{K\ddot{u}rz}(M,\star)\vdash{}\\[-2pt]
    &\EqRel{\FormalDifferenceCarrier{M,\star,e},
      \FormalDifferenceRel{M,\star,e}}
  \end{aligned}
}{GroupCompletionPairEquivalence}{%
  \DeltaRow{Mengen}{M\dsep e\dsep x\dsep y\dsep z\dsep
    a\dsep b\dsep c\dsep d\dsep f\dsep g}
  \DeltaRow{Funktionen}{\star}
  \DeltaRow{Relationsoperatoren}{\FormalDifferenceRel{M,\star,e}}
}
\begin{tabproofsplitwide}
  \proofpartwideindR[Reflexivität]{%
    \begin{aligned}[t]
      &\CommutativeMonoidAlg{M}{\star}{e}\dsep
        x\in\FormalDifferenceCarrier{M,\star,e}\vdash{}\\[-2pt]
      &x\FormalDifferenceRel{M,\star,e}x
    \end{aligned}
  }{%
    \CommutativeMonoidAlg{M}{\star}{e}\dsep
    x\in\FormalDifferenceCarrier{M,\star,e}
    \vdash x\FormalDifferenceRel{M,\star,e}x
  }[GroupCompletionPairEquivalenceReflexive]
    \proofstepwidestar[1]{\CommutativeMonoidAlg{M}{\star}{e}}{\rA}
    \proofstepwidestar[2]{x\in\FormalDifferenceCarrier{M,\star,e}}{\rA}
    \proofstepwidestar[1]{\FormalDifferenceCarrier{M,\star,e}=M\times M}{%
      \FormulaRefAuto{GroupCompletionPairCarrier}[def]{1}}
    \proofstepwidestar[1,2]{x\in M\times M}{\rIE{3,2}}
    \proofstepwidestar[1,2]{\exists a\in M\,\exists b\in M\,x=(a,b)}{%
      \FormulaRefAuto{x\in A\times B\eqvdash
        \exists a\in A\,\exists b\in B\,x=(a,b)}{4}}
    \proofstepwidestar[6]{a,b\in M\land x=(a,b)}{\rA}
    \proofstepwidestar[1,6]{%
      (a,b)\FormalDifferenceRel{M,\star,e}(a,b)}{%
      \FormulaRefAuto{GroupCompletionPairRelationReflexive}{1,6}}
    \proofstepwidestar[1,6]{x\FormalDifferenceRel{M,\star,e}x}{\rIE{6,7}}
    \proofstepwidestar[1,2]{x\FormalDifferenceRel{M,\star,e}x}{%
      \rEE{5,6,8}}
  \closeproofpartwideindR

  \proofpartwideindR[Symmetrie]{%
    \begin{aligned}[t]
      &\CommutativeMonoidAlg{M}{\star}{e}\dsep
        x\FormalDifferenceRel{M,\star,e}y\vdash{}\\[-2pt]
      &y\FormalDifferenceRel{M,\star,e}x
    \end{aligned}
  }{%
    \CommutativeMonoidAlg{M}{\star}{e}\dsep
    x\FormalDifferenceRel{M,\star,e}y
    \vdash y\FormalDifferenceRel{M,\star,e}x
  }[GroupCompletionPairEquivalenceSymmetric]
    \proofstepwidestar[1]{\CommutativeMonoidAlg{M}{\star}{e}}{\rA}
    \proofstepwidestar[2]{x\FormalDifferenceRel{M,\star,e}y}{\rA}
    \proofstepwidestar[3]{x\in\FormalDifferenceCarrier{M,\star,e}}{\rA}
    \proofstepwidestar[4]{y\in\FormalDifferenceCarrier{M,\star,e}}{\rA}
    \proofstepwidestar[1,3]{\exists a\in M\,\exists b\in M\,x=(a,b)}{%
      \FormulaRefAuto{GroupCompletionPairRepresentative}{1,3}}
    \proofstepwidestar[1,4]{\exists c\in M\,\exists d\in M\,y=(c,d)}{%
      \FormulaRefAuto{GroupCompletionPairRepresentative}{1,4}}
    \proofstepwidestar[7]{a,b\in M\land x=(a,b)}{\rA}
    \proofstepwidestar[8]{c,d\in M\land y=(c,d)}{\rA}
    \proofstepwidestar[7,8]{a,b,c,d\in M}{%
      \rAI{\rAEa{7},\rAEa{8}}}
    \proofstepwidestar[2,7,8]{%
      (a,b)\FormalDifferenceRel{M,\star,e}(c,d)}{%
      \rIE{\rAEb{7},\rAEb{8},2}}
    \proofstepwidestar[1,2,7,8]{%
      (c,d)\FormalDifferenceRel{M,\star,e}(a,b)}{%
      \FormulaRefAuto{GroupCompletionPairRelationSymmetric}{1,9,10}}
    \proofstepwidestar[1,2,7,8]{y\FormalDifferenceRel{M,\star,e}x}{%
      \rIE{\rAEb{8},\rAEb{7},11}}
    \proofstepwidestar[1,2,3,4,7]{y\FormalDifferenceRel{M,\star,e}x}{%
      \rEE{6,8,12}}
    \proofstepwidestar[1,2,3,4]{y\FormalDifferenceRel{M,\star,e}x}{%
      \rEE{5,7,13}}
  \closeproofpartwideindR

  \proofpartwideindR[Transitivität]{%
    \begin{aligned}[t]
      &\CommutativeMonoidAlg{M}{\star}{e}\dsep
        \mathsf{K\ddot{u}rz}(M,\star)\dsep
        x\FormalDifferenceRel{M,\star,e}y\dsep
        y\FormalDifferenceRel{M,\star,e}z\vdash{}\\[-2pt]
      &x\FormalDifferenceRel{M,\star,e}z
    \end{aligned}
  }{%
    \CommutativeMonoidAlg{M}{\star}{e}\dsep
    \mathsf{K\ddot{u}rz}(M,\star)\dsep
    x\FormalDifferenceRel{M,\star,e}y\dsep
    y\FormalDifferenceRel{M,\star,e}z
    \vdash x\FormalDifferenceRel{M,\star,e}z
  }[GroupCompletionPairEquivalenceTransitive]
    \proofstepwidestar[1]{\CommutativeMonoidAlg{M}{\star}{e}}{\rA}
    \proofstepwidestar[2]{\mathsf{K\ddot{u}rz}(M,\star)}{\rA}
    \proofstepwidestar[3]{x\FormalDifferenceRel{M,\star,e}y}{\rA}
    \proofstepwidestar[4]{y\FormalDifferenceRel{M,\star,e}z}{\rA}
    \proofstepwidestar[5]{x\in\FormalDifferenceCarrier{M,\star,e}}{\rA}
    \proofstepwidestar[6]{y\in\FormalDifferenceCarrier{M,\star,e}}{\rA}
    \proofstepwidestar[7]{z\in\FormalDifferenceCarrier{M,\star,e}}{\rA}
    \proofstepwidestar[1,5]{\exists a\in M\,\exists b\in M\,x=(a,b)}{%
      \FormulaRefAuto{GroupCompletionPairRepresentative}{1,5}}
    \proofstepwidestar[1,6]{\exists c\in M\,\exists d\in M\,y=(c,d)}{%
      \FormulaRefAuto{GroupCompletionPairRepresentative}{1,6}}
    \proofstepwidestar[1,7]{\exists f\in M\,\exists g\in M\,z=(f,g)}{%
      \FormulaRefAuto{GroupCompletionPairRepresentative}{1,7}}
    \proofstepwidestar[11]{a,b\in M\land x=(a,b)}{\rA}
    \proofstepwidestar[12]{c,d\in M\land y=(c,d)}{\rA}
    \proofstepwidestar[13]{f,g\in M\land z=(f,g)}{\rA}
    \proofstepwidestar[11,12,13]{a,b,c,d,f,g\in M}{%
      \rAI{\rAEa{11},\rAI{\rAEa{12},\rAEa{13}}}}
    \proofstepwidestar[3,11,12]{%
      (a,b)\FormalDifferenceRel{M,\star,e}(c,d)}{%
      \rIE{\rAEb{11},\rAEb{12},3}}
    \proofstepwidestar[4,12,13]{%
      (c,d)\FormalDifferenceRel{M,\star,e}(f,g)}{%
      \rIE{\rAEb{12},\rAEb{13},4}}
    \proofstepwidestar[1,2,3,4,11,12,13]{%
      (a,b)\FormalDifferenceRel{M,\star,e}(f,g)}{%
      \FormulaRefAuto{GroupCompletionPairRelationTransitive}{1,2,14,15,16}}
    \proofstepwidestar[1,2,3,4,11,12,13]{%
      x\FormalDifferenceRel{M,\star,e}z}{%
      \rIE{\rAEb{11},\rAEb{13},17}}
    \proofstepwidestar[1,2,3,4,5,6,7,11,12]{%
      x\FormalDifferenceRel{M,\star,e}z}{\rEE{10,13,18}}
    \proofstepwidestar[1,2,3,4,5,6,7,11]{%
      x\FormalDifferenceRel{M,\star,e}z}{\rEE{9,12,19}}
    \proofstepwidestar[1,2,3,4,5,6,7]{%
      x\FormalDifferenceRel{M,\star,e}z}{\rEE{8,11,20}}
  \closeproofpartwideindR

  \proofpartwide{Schluss}
    \proofstepwidestar[1]{\CommutativeMonoidAlg{M}{\star}{e}}{\rA}
    \proofstepwidestar[2]{\mathsf{K\ddot{u}rz}(M,\star)}{\rA}
    \proofstepwidestar[1]{%
      x\in\FormalDifferenceCarrier{M,\star,e}\rightarrow
      x\FormalDifferenceRel{M,\star,e}x}{%
      \FormulaRefAuto{GroupCompletionPairEquivalenceReflexive}{1}}
    \proofstepwidestar[1]{%
      x\FormalDifferenceRel{M,\star,e}y\rightarrow
      y\FormalDifferenceRel{M,\star,e}x}{%
      \FormulaRefAuto{GroupCompletionPairEquivalenceSymmetric}{1}}
    \proofstepwidestar[3]{x\FormalDifferenceRel{M,\star,e}y}{\rA}
    \proofstepwidestar[4]{y\FormalDifferenceRel{M,\star,e}z}{\rA}
    \proofstepwidestar[1,2,3,4]{x\FormalDifferenceRel{M,\star,e}z}{%
      \FormulaRefAuto{GroupCompletionPairEquivalenceTransitive}{1,2,3,4}}
    \proofstepwidestar[1,2]{%
      \EqRel{\FormalDifferenceCarrier{M,\star,e},
        \FormalDifferenceRel{M,\star,e}}}{%
      \FormulaRefAuto{EqRelDef}{3,4,7}}
  \closeproofpartwide
\end{tabproofsplitwide}

\section{Quotient und Differenzenklassen}

\FormulaDefDeltaK[Träger der formalen Differenzen]{%
  \begin{aligned}[t]
    &\CommutativeMonoidAlg{M}{\star}{e}\dsep
      \mathsf{K\ddot{u}rz}(M,\star)\vdash{}\\[-2pt]
    &\GroupCompletionCarrier{M,\star,e}\coloneqq
      \QuotientSet{\FormalDifferenceCarrier{M,\star,e}}
        {\FormalDifferenceRel{M,\star,e}}
  \end{aligned}
}{GroupCompletionCarrierDef}{%
  \DeltaRow{Mengen}{M\dsep e\dsep\GroupCompletionCarrier{M,\star,e}}
  \DeltaRow{Funktionen}{\star}
  \DeltaRow{Relationsoperatoren}{\FormalDifferenceRel{M,\star,e}}
  \DeltaRow{Neue Symbole}{\GroupCompletionCarrier{M,\star,e}}
}

\FormulaDefDeltaK[Klasse eines formalen Differenzenpaares]{%
  \begin{aligned}[t]
    &\CommutativeMonoidAlg{M}{\star}{e}\dsep a,b\in M\vdash{}\\[-2pt]
    &\FormalDifferenceClass{M,\star,e}{a}{b}\coloneqq
      \EqClass{(a,b)}{\FormalDifferenceRel{M,\star,e}}
        {\FormalDifferenceCarrier{M,\star,e}}
  \end{aligned}
}{GroupCompletionClassDef}{%
  \DeltaRow{Mengen}{M\dsep e\dsep a\dsep b}
  \DeltaRow{Funktionen}{\star}
  \DeltaRow{Relationsoperatoren}{\FormalDifferenceRel{M,\star,e}}
  \DeltaRow{Neue Symbole}{\FormalDifferenceClass{M,\star,e}{a}{b}}
}

\FormulaThmDeltaK[Differenzenklassen liegen im Quotienten]{%
  \begin{aligned}[t]
    &\CommutativeMonoidAlg{M}{\star}{e}\dsep
      \mathsf{K\ddot{u}rz}(M,\star)\dsep a,b\in M\vdash{}\\[-2pt]
    &\FormalDifferenceClass{M,\star,e}{a}{b}
      \in\GroupCompletionCarrier{M,\star,e}
  \end{aligned}
}{GroupCompletionClassInCarrier}{%
  \DeltaRow{Mengen}{M\dsep e\dsep a\dsep b}
  \DeltaRow{Funktionen}{\star}
}
\begin{tabproof}
  \proofstep{1}{\CommutativeMonoidAlg{M}{\star}{e}}{\rA}
  \proofstep{2}{\mathsf{K\ddot{u}rz}(M,\star)}{\rA}
  \proofstep{3}{a,b\in M}{\rA}
  \proofstep{1}{\FormalDifferenceCarrier{M,\star,e}=M\times M}{%
    \FormulaRefAuto{GroupCompletionPairCarrier}[def]{1}}
  \proofstep{1,3}{(a,b)\in\FormalDifferenceCarrier{M,\star,e}}{%
    \rIE{4,\FormulaRefAuto{(a,b)\in A\times B\eqvdash
      a\in A\land b\in B}{3}}}
  \proofstep{1,2}{%
    \EqRel{\FormalDifferenceCarrier{M,\star,e},
      \FormalDifferenceRel{M,\star,e}}}{%
    \FormulaRefAuto{GroupCompletionPairEquivalence}{1,2}}
  \proofstep{1,2,3}{%
    \EqClass{(a,b)}{\FormalDifferenceRel{M,\star,e}}
      {\FormalDifferenceCarrier{M,\star,e}}
    \in\QuotientSet{\FormalDifferenceCarrier{M,\star,e}}
      {\FormalDifferenceRel{M,\star,e}}}{%
    \FormulaRefAuto{EqClassInQuotient}{6,5}}
  \proofstep{1,3}{%
    \FormalDifferenceClass{M,\star,e}{a}{b}
    =\EqClass{(a,b)}{\FormalDifferenceRel{M,\star,e}}
      {\FormalDifferenceCarrier{M,\star,e}}}{%
    \FormulaRefAuto{GroupCompletionClassDef}[def]{1,3}}
  \proofstep{1,2}{%
    \GroupCompletionCarrier{M,\star,e}
    =\QuotientSet{\FormalDifferenceCarrier{M,\star,e}}
      {\FormalDifferenceRel{M,\star,e}}}{%
    \FormulaRefAuto{GroupCompletionCarrierDef}[def]{1,2}}
  \proofstep{1,2,3}{%
    \FormalDifferenceClass{M,\star,e}{a}{b}
      \in\GroupCompletionCarrier{M,\star,e}}{\rIE{8,9,7}}
\end{tabproof}

\FormulaThmDeltaK[Gleichheit formaler Differenzenklassen]{%
  \begin{aligned}[t]
    &\CommutativeMonoidAlg{M}{\star}{e}\dsep
      \mathsf{K\ddot{u}rz}(M,\star)\dsep a,b,c,d\in M\vdash{}\\[-2pt]
    &\FormalDifferenceClass{M,\star,e}{a}{b}
      =\FormalDifferenceClass{M,\star,e}{c}{d}
      \leftrightarrow a\star d=b\star c
  \end{aligned}
}{GroupCompletionClassEquality}{%
  \DeltaRow{Mengen}{M\dsep e\dsep a\dsep b\dsep c\dsep d}
  \DeltaRow{Funktionen}{\star}
  \DeltaRow{Relationsoperatoren}{\FormalDifferenceRel{M,\star,e}}
}
\begin{tabproofsplitwide}
  \proofpartwideindR[Klassenkriterium]{%
    \begin{aligned}[t]
      &\CommutativeMonoidAlg{M}{\star}{e}\dsep
        \mathsf{K\ddot{u}rz}(M,\star)\dsep a,b,c,d\in M\vdash{}\\[-2pt]
      &\FormalDifferenceClass{M,\star,e}{a}{b}
        =\FormalDifferenceClass{M,\star,e}{c}{d}
        \leftrightarrow(a,b)\FormalDifferenceRel{M,\star,e}(c,d)
    \end{aligned}
  }{%
    \CommutativeMonoidAlg{M}{\star}{e}\dsep
    \mathsf{K\ddot{u}rz}(M,\star)\dsep a,b,c,d\in M\vdash
    \FormalDifferenceClass{M,\star,e}{a}{b}
      =\FormalDifferenceClass{M,\star,e}{c}{d}
      \leftrightarrow(a,b)\FormalDifferenceRel{M,\star,e}(c,d)
  }[GroupCompletionClassEqualityViaRelation]
    \proofstepwidestar[1]{\CommutativeMonoidAlg{M}{\star}{e}}{\rA}
    \proofstepwidestar[2]{\mathsf{K\ddot{u}rz}(M,\star)}{\rA}
    \proofstepwidestar[3]{a,b,c,d\in M}{\rA}
    \proofstepwidestar[1,2]{%
      \EqRel{\FormalDifferenceCarrier{M,\star,e},
        \FormalDifferenceRel{M,\star,e}}}{%
      \FormulaRefAuto{GroupCompletionPairEquivalence}{1,2}}
    \proofstepwidestar[1]{%
      \FormalDifferenceCarrier{M,\star,e}=M\times M}{%
      \FormulaRefAuto{GroupCompletionPairCarrier}[def]{1}}
    \proofstepwidestar[1,3]{%
      (a,b)\in\FormalDifferenceCarrier{M,\star,e}}{%
      \rIE{5,\FormulaRefAuto{(a,b)\in A\times B
        \eqvdash a\in A\land b\in B}{3}}}
    \proofstepwidestar[1,3]{%
      (c,d)\in\FormalDifferenceCarrier{M,\star,e}}{%
      \rIE{5,\FormulaRefAuto{(a,b)\in A\times B
        \eqvdash a\in A\land b\in B}{3}}}
    \proofstepwidestar[1,2,3]{%
      \begin{aligned}[t]
        &\EqClass{(a,b)}{\FormalDifferenceRel{M,\star,e}}
          {\FormalDifferenceCarrier{M,\star,e}}
          =\EqClass{(c,d)}{\FormalDifferenceRel{M,\star,e}}
          {\FormalDifferenceCarrier{M,\star,e}}\\[-2pt]
        &\qquad\leftrightarrow
          (a,b)\FormalDifferenceRel{M,\star,e}(c,d)
      \end{aligned}}{%
      \FormulaRefAuto{EqClassEqualIffRel}{4,6,7}}
    \proofstepwidestar[1,3]{%
      \FormalDifferenceClass{M,\star,e}{a}{b}
      =\EqClass{(a,b)}{\FormalDifferenceRel{M,\star,e}}
        {\FormalDifferenceCarrier{M,\star,e}}}{%
      \FormulaRefAuto{GroupCompletionClassDef}[def]{1,3}}
    \proofstepwidestar[1,3]{%
      \FormalDifferenceClass{M,\star,e}{c}{d}
      =\EqClass{(c,d)}{\FormalDifferenceRel{M,\star,e}}
        {\FormalDifferenceCarrier{M,\star,e}}}{%
      \FormulaRefAuto{GroupCompletionClassDef}[def]{1,3}}
    \proofstepwidestar[1,2,3]{%
      \begin{aligned}[t]
        &\FormalDifferenceClass{M,\star,e}{a}{b}
          =\FormalDifferenceClass{M,\star,e}{c}{d}\\[-2pt]
        &\qquad\leftrightarrow
          (a,b)\FormalDifferenceRel{M,\star,e}(c,d)
      \end{aligned}}{%
      \rIE{8,9,10}}
  \closeproofpartwideindR

  \proofpartwide{Schluss}
    \proofstepwidestar[1]{\CommutativeMonoidAlg{M}{\star}{e}}{\rA}
    \proofstepwidestar[2]{\mathsf{K\ddot{u}rz}(M,\star)}{\rA}
    \proofstepwidestar[3]{a,b,c,d\in M}{\rA}
    \proofstepwidestar[1,2,3]{%
      \FormalDifferenceClass{M,\star,e}{a}{b}
      =\FormalDifferenceClass{M,\star,e}{c}{d}
      \leftrightarrow(a,b)\FormalDifferenceRel{M,\star,e}(c,d)}{%
      \FormulaRefAuto{GroupCompletionClassEqualityViaRelation}{1,2,3}}
    \proofstepwidestar[1,3]{%
      (a,b)\FormalDifferenceRel{M,\star,e}(c,d)
      \leftrightarrow a\star d=b\star c}{%
      \FormulaRefAuto{GroupCompletionPairRelationChar}{1,3}}
    \proofstepwidestar[1,2,3]{%
      \FormalDifferenceClass{M,\star,e}{a}{b}
      =\FormalDifferenceClass{M,\star,e}{c}{d}
      \leftrightarrow a\star d=b\star c}{%
      \FormulaRefAuto{P\leftrightarrow Q\dsep
        Q\leftrightarrow R\vdash P\leftrightarrow R}{4,5}}
  \closeproofpartwide
\end{tabproofsplitwide}

\FormulaThmDeltaK[Jede formale Differenz besitzt ein Paar]{%
  \begin{aligned}[t]
    &\CommutativeMonoidAlg{M}{\star}{e}\dsep
      \mathsf{K\ddot{u}rz}(M,\star)\dsep
      z\in\GroupCompletionCarrier{M,\star,e}\vdash{}\\[-2pt]
    &\exists a\in M\,\exists b\in M\,
      z=\FormalDifferenceClass{M,\star,e}{a}{b}
  \end{aligned}
}{GroupCompletionRepresentative}{%
  \DeltaRow{Mengen}{M\dsep e\dsep z\dsep x\dsep a\dsep b}
  \DeltaRow{Funktionen}{\star}
}
\begin{tabproofsplitwide}
  \proofpartwideindR[Quotientenrepräsentant]{%
    \begin{aligned}[t]
      &\CommutativeMonoidAlg{M}{\star}{e}\dsep
        \mathsf{K\ddot{u}rz}(M,\star)\dsep
        z\in\GroupCompletionCarrier{M,\star,e}\vdash{}\\[-2pt]
      &\exists x\in\FormalDifferenceCarrier{M,\star,e}\,
        z=\EqClass{x}{\FormalDifferenceRel{M,\star,e}}
          {\FormalDifferenceCarrier{M,\star,e}}
    \end{aligned}
  }{%
    \CommutativeMonoidAlg{M}{\star}{e}\dsep
    \mathsf{K\ddot{u}rz}(M,\star)\dsep
    z\in\GroupCompletionCarrier{M,\star,e}\vdash
    \exists x\in\FormalDifferenceCarrier{M,\star,e}\,
      z=\EqClass{x}{\FormalDifferenceRel{M,\star,e}}
        {\FormalDifferenceCarrier{M,\star,e}}
  }[GroupCompletionQuotientRepresentative]
    \proofstepwidestar[1]{\CommutativeMonoidAlg{M}{\star}{e}}{\rA}
    \proofstepwidestar[2]{\mathsf{K\ddot{u}rz}(M,\star)}{\rA}
    \proofstepwidestar[3]{z\in\GroupCompletionCarrier{M,\star,e}}{\rA}
    \proofstepwidestar[1,2]{%
      \GroupCompletionCarrier{M,\star,e}
      =\QuotientSet{\FormalDifferenceCarrier{M,\star,e}}
        {\FormalDifferenceRel{M,\star,e}}}{%
      \FormulaRefAuto{GroupCompletionCarrierDef}[def]{1,2}}
    \proofstepwidestar[1,2,3]{%
      z\in\QuotientSet{\FormalDifferenceCarrier{M,\star,e}}
        {\FormalDifferenceRel{M,\star,e}}}{\rIE{4,3}}
    \proofstepwidestar[1,2]{%
      \EqRel{\FormalDifferenceCarrier{M,\star,e},
        \FormalDifferenceRel{M,\star,e}}}{%
      \FormulaRefAuto{GroupCompletionPairEquivalence}{1,2}}
    \proofstepwidestar[1,2,3]{%
      \exists x\in\FormalDifferenceCarrier{M,\star,e}\,
      z=\EqClass{x}{\FormalDifferenceRel{M,\star,e}}
        {\FormalDifferenceCarrier{M,\star,e}}}{%
      \FormulaRefAuto{QuotientRepresentative}{6,5}}
  \closeproofpartwideindR

  \proofpartwide{Schluss}
    \proofstepwidestar[1]{\CommutativeMonoidAlg{M}{\star}{e}}{\rA}
    \proofstepwidestar[2]{\mathsf{K\ddot{u}rz}(M,\star)}{\rA}
    \proofstepwidestar[3]{z\in\GroupCompletionCarrier{M,\star,e}}{\rA}
    \proofstepwidestar[1,2,3]{%
      \exists x\in\FormalDifferenceCarrier{M,\star,e}\,
      z=\EqClass{x}{\FormalDifferenceRel{M,\star,e}}
        {\FormalDifferenceCarrier{M,\star,e}}}{%
      \FormulaRefAuto{GroupCompletionQuotientRepresentative}{1,2,3}}
    \proofstepwidestar[5]{%
      x\in\FormalDifferenceCarrier{M,\star,e}\land
      z=\EqClass{x}{\FormalDifferenceRel{M,\star,e}}
        {\FormalDifferenceCarrier{M,\star,e}}}{\rA}
    \proofstepwidestar[1,5]{x\in M\times M}{%
      \rIE{\FormulaRefAuto{GroupCompletionPairCarrier}[def]{1},\rAEa{5}}}
    \proofstepwidestar[1,5]{\exists a\in M\,\exists b\in M\,x=(a,b)}{%
      \FormulaRefAuto{x\in A\times B\eqvdash
        \exists a\in A\,\exists b\in B\,x=(a,b)}{6}}
    \proofstepwidestar[8]{a,b\in M\land x=(a,b)}{\rA}
    \proofstepwidestar[5,8]{%
      z=\EqClass{(a,b)}{\FormalDifferenceRel{M,\star,e}}
        {\FormalDifferenceCarrier{M,\star,e}}}{%
      \rIE{\rAEb{5},\rAEb{8}}}
    \proofstepwidestar[1,8]{%
      \FormalDifferenceClass{M,\star,e}{a}{b}
      =\EqClass{(a,b)}{\FormalDifferenceRel{M,\star,e}}
        {\FormalDifferenceCarrier{M,\star,e}}}{%
      \FormulaRefAuto{GroupCompletionClassDef}[def]{1,8}}
    \proofstepwidestar[1,5,8]{%
      z=\FormalDifferenceClass{M,\star,e}{a}{b}}{%
      \FormulaRefAuto{a=b\dsep c=b\vdash a=c}{9,10}}
    \proofstepwidestar[1,5]{%
      \exists a\in M\,\exists b\in M\,
      z=\FormalDifferenceClass{M,\star,e}{a}{b}}{\rEE{7,8,11}}
    \proofstepwidestar[1,2,3]{%
      \exists a\in M\,\exists b\in M\,
      z=\FormalDifferenceClass{M,\star,e}{a}{b}}{\rEE{4,5,12}}
  \closeproofpartwide
\end{tabproofsplitwide}

\section{Operation auf formalen Differenzen}

\FormulaThmDeltaK[Repräsentantenunabhängigkeit der Differenzenoperation]{%
  \begin{aligned}[t]
    &\CommutativeMonoidAlg{M}{\star}{e}\dsep
      \mathsf{K\ddot{u}rz}(M,\star)\dsep
      a,b,a',b',c,d,c',d'\in M\dsep{}\\[-2pt]
    &\FormalDifferenceClass{M,\star,e}{a}{b}
      =\FormalDifferenceClass{M,\star,e}{a'}{b'}\dsep
      \FormalDifferenceClass{M,\star,e}{c}{d}
      =\FormalDifferenceClass{M,\star,e}{c'}{d'}\vdash{}\\[-2pt]
    &\FormalDifferenceClass{M,\star,e}{a\star c}{b\star d}
      =\FormalDifferenceClass{M,\star,e}{a'\star c'}{b'\star d'}
  \end{aligned}
}{GroupCompletionOperationCompatible}{%
  \DeltaRow{Mengen}{M\dsep e\dsep a\dsep b\dsep a'\dsep b'\dsep
    c\dsep d\dsep c'\dsep d'}
  \DeltaRow{Funktionen}{\star}
}
\begin{tabproofwide}
  \proofstepwidestar[1]{\CommutativeMonoidAlg{M}{\star}{e}}{\rA}
  \proofstepwidestar[2]{\mathsf{K\ddot{u}rz}(M,\star)}{\rA}
  \proofstepwidestar[3]{a,b,a',b',c,d,c',d'\in M}{\rA}
  \proofstepwidestar[4]{%
    \FormalDifferenceClass{M,\star,e}{a}{b}
      =\FormalDifferenceClass{M,\star,e}{a'}{b'}}{\rA}
  \proofstepwidestar[5]{%
    \FormalDifferenceClass{M,\star,e}{c}{d}
      =\FormalDifferenceClass{M,\star,e}{c'}{d'}}{\rA}
  \proofstepwidestar[1,2,3,4]{a\star b'=b\star a'}{%
    \FormulaRefAuto{GroupCompletionClassEquality}{1,2,3,4}}
  \proofstepwidestar[1,2,3,5]{c\star d'=d\star c'}{%
    \FormulaRefAuto{GroupCompletionClassEquality}{1,2,3,5}}
  \proofstepwidestar[1]{\CommutativeSemigroupAlg{M}{\star}}{%
    \FormulaRefAuto{CommutativeMonoidIsCommutativeSemigroup}{1}}
  \proofstepwide[1,2,3,4,5]{%
    (a\star c)\star(b'\star d')}{=}{%
    (a\star b')\star(c\star d')}{%
    \FormulaRefAuto{CommutativeSemigroupFourTermExchange}{8,3}}
  \proofstepwide[1,2,3,4,5]{}{=}{%
    (b\star a')\star(d\star c')}{\rIE{6,7}}
  \proofstepwide[1,2,3,4,5]{}{=}{%
    (b\star d)\star(a'\star c')}{%
    \FormulaRefAuto{CommutativeSemigroupFourTermExchange}{8,3}}
  \proofstepwidestar[1,2,3,4,5]{%
    (a\star c)\star(b'\star d')
      =(b\star d)\star(a'\star c')}{\rChain{9,11}}
  \proofstepwidestar[1,3]{a\star c,b\star d,a'\star c',b'\star d'\in M}{%
    \FormulaRefAuto{SemigroupClosureAxiom}[ax]{%
      \FormulaRefAuto{CommutativeSemigroupBaseAxiom}[ax]{8},3}}
  \proofstepwidestar[1,2,3,4,5]{%
    \FormalDifferenceClass{M,\star,e}{a\star c}{b\star d}
      =\FormalDifferenceClass{M,\star,e}{a'\star c'}{b'\star d'}}{%
    \FormulaRefAuto{GroupCompletionClassEquality}{1,2,13,12}}
\end{tabproofwide}

\FormulaThmDeltaK[Quotientenabstieg der Differenzenoperation]{%
  \begin{aligned}[t]
    &\CommutativeMonoidAlg{M}{\star}{e}\dsep
      \mathsf{K\ddot{u}rz}(M,\star)\dsep
      z,w\in\GroupCompletionCarrier{M,\star,e}\vdash{}\\[-2pt]
    &\exists! u\in\GroupCompletionCarrier{M,\star,e}\,
      \exists p,q,r,s\in M\,\Bigl(
        z=\FormalDifferenceClass{M,\star,e}{p}{q}\land{}\\[-2pt]
    &\qquad w=\FormalDifferenceClass{M,\star,e}{r}{s}\land
        u=\FormalDifferenceClass{M,\star,e}{p\star r}{q\star s}
      \Bigr)
  \end{aligned}
}{GroupCompletionOperationQuotientDescent}{%
  \DeltaRow{Mengen}{M\dsep e\dsep z\dsep w\dsep u\dsep v\dsep
    a\dsep b\dsep c\dsep d\dsep a'\dsep b'\dsep c'\dsep d'\dsep
    p\dsep q\dsep r\dsep s}
  \DeltaRow{Funktionen}{\star}
}
\begin{tabproofsplitwide}
  \proofpartwideindR[Ergebniskandidat fester Repräsentanten]{%
    \begin{aligned}[t]
      &\CommutativeMonoidAlg{M}{\star}{e}\dsep
        \mathsf{K\ddot{u}rz}(M,\star)\dsep a,b,c,d\in M\dsep{}\\[-2pt]
      &z=\FormalDifferenceClass{M,\star,e}{a}{b}\dsep
        w=\FormalDifferenceClass{M,\star,e}{c}{d}\vdash{}\\[-2pt]
      &\FormalDifferenceClass{M,\star,e}{a\star c}{b\star d}
        \in\GroupCompletionCarrier{M,\star,e}
    \end{aligned}
  }{%
    \CommutativeMonoidAlg{M}{\star}{e}\dsep
    \mathsf{K\ddot{u}rz}(M,\star)\dsep a,b,c,d\in M\dsep
    z=\FormalDifferenceClass{M,\star,e}{a}{b}\dsep
    w=\FormalDifferenceClass{M,\star,e}{c}{d}
    \vdash
    \FormalDifferenceClass{M,\star,e}{a\star c}{b\star d}
      \in\GroupCompletionCarrier{M,\star,e}
  }[GroupCompletionOperationDescentCandidate]
    \proofstepwidestar[1]{\CommutativeMonoidAlg{M}{\star}{e}}{\rA}
    \proofstepwidestar[2]{\mathsf{K\ddot{u}rz}(M,\star)}{\rA}
    \proofstepwidestar[3]{a,b,c,d\in M}{\rA}
    \proofstepwidestar[1]{\CommutativeSemigroupAlg{M}{\star}}{%
      \FormulaRefAuto{CommutativeMonoidIsCommutativeSemigroup}{1}}
    \proofstepwidestar[1,3]{a\star c,b\star d\in M}{%
      \FormulaRefAuto{SemigroupClosureAxiom}[ax]{%
        \FormulaRefAuto{CommutativeSemigroupBaseAxiom}[ax]{4},3}}
    \proofstepwidestar[1,2,3]{%
      \FormalDifferenceClass{M,\star,e}{a\star c}{b\star d}
      \in\GroupCompletionCarrier{M,\star,e}}{%
      \FormulaRefAuto{GroupCompletionClassInCarrier}{1,2,5}}
  \closeproofpartwideindR

  \proofpartwideindR[Existenz durch Repräsentantenwahl]{%
    \begin{aligned}[t]
      &\CommutativeMonoidAlg{M}{\star}{e}\dsep
        \mathsf{K\ddot{u}rz}(M,\star)\dsep
        z,w\in\GroupCompletionCarrier{M,\star,e}\vdash{}\\[-2pt]
      &\exists u\in\GroupCompletionCarrier{M,\star,e}\,
        \exists p,q,r,s\in M\,\Bigl(
          z=\FormalDifferenceClass{M,\star,e}{p}{q}\land{}\\[-2pt]
      &\qquad w=\FormalDifferenceClass{M,\star,e}{r}{s}\land
          u=\FormalDifferenceClass{M,\star,e}{p\star r}{q\star s}
        \Bigr)
    \end{aligned}
  }{%
    \CommutativeMonoidAlg{M}{\star}{e}\dsep
    \mathsf{K\ddot{u}rz}(M,\star)\dsep
    z,w\in\GroupCompletionCarrier{M,\star,e}\vdash
    \exists u\in\GroupCompletionCarrier{M,\star,e}\,
    \exists p,q,r,s\in M\,
    \bigl(z=\FormalDifferenceClass{M,\star,e}{p}{q}\land
      w=\FormalDifferenceClass{M,\star,e}{r}{s}\land
      u=\FormalDifferenceClass{M,\star,e}{p\star r}{q\star s}\bigr)
  }[GroupCompletionOperationDescentExistence]
    \proofstepwidestar[1]{\CommutativeMonoidAlg{M}{\star}{e}}{\rA}
    \proofstepwidestar[2]{\mathsf{K\ddot{u}rz}(M,\star)}{\rA}
    \proofstepwidestar[3]{z,w\in\GroupCompletionCarrier{M,\star,e}}{\rA}
    \proofstepwidestar[1,2,3]{\exists a,b\in M\,
      z=\FormalDifferenceClass{M,\star,e}{a}{b}}{%
      \FormulaRefAuto{GroupCompletionRepresentative}{1,2,\rAEa{3}}}
    \proofstepwidestar[1,2,3]{\exists c,d\in M\,
      w=\FormalDifferenceClass{M,\star,e}{c}{d}}{%
      \FormulaRefAuto{GroupCompletionRepresentative}{1,2,\rAEb{3}}}
    \proofstepwidestar[6]{%
      \begin{aligned}[t]
        &a,b,c,d\in M\land
          z=\FormalDifferenceClass{M,\star,e}{a}{b}\land{}\\[-2pt]
        &w=\FormalDifferenceClass{M,\star,e}{c}{d}
      \end{aligned}}{\rA}
    \proofstepwidestar[1,2,6]{%
      \FormalDifferenceClass{M,\star,e}{a\star c}{b\star d}
      \in\GroupCompletionCarrier{M,\star,e}}{%
      \FormulaRefAuto{GroupCompletionOperationDescentCandidate}{1,2,6}}
    \proofstepwidestar[1,2,6]{%
      \begin{aligned}[t]
        &\exists u\in\GroupCompletionCarrier{M,\star,e}\,
          \exists p,q,r,s\in M\,\Bigl(
          z=\FormalDifferenceClass{M,\star,e}{p}{q}\land{}\\[-2pt]
        &\qquad w=\FormalDifferenceClass{M,\star,e}{r}{s}\land
          u=\FormalDifferenceClass{M,\star,e}{p\star r}{q\star s}
          \Bigr)
      \end{aligned}}{%
      \rEI{\rAI{7,\rAI{6,\rII}}}}
    \proofstepwidestar[1,2,3]{%
      \begin{aligned}[t]
        &\exists u\in\GroupCompletionCarrier{M,\star,e}\,
          \exists p,q,r,s\in M\,\Bigl(
          z=\FormalDifferenceClass{M,\star,e}{p}{q}\land{}\\[-2pt]
        &\qquad w=\FormalDifferenceClass{M,\star,e}{r}{s}\land
          u=\FormalDifferenceClass{M,\star,e}{p\star r}{q\star s}
          \Bigr)
      \end{aligned}}{%
      \rEE{4,5,6,8}}
  \closeproofpartwideindR

  \proofpartwideindR[Eindeutigkeit fester Repräsentanten]{%
    \begin{aligned}[t]
      &\CommutativeMonoidAlg{M}{\star}{e}\dsep
        \mathsf{K\ddot{u}rz}(M,\star)\dsep
        a,b,c,d,a',b',c',d'\in M\dsep{}\\[-2pt]
      &z=\FormalDifferenceClass{M,\star,e}{a}{b}\land
        w=\FormalDifferenceClass{M,\star,e}{c}{d}\land
        u=\FormalDifferenceClass{M,\star,e}{a\star c}{b\star d}\dsep{}\\[-2pt]
      &z=\FormalDifferenceClass{M,\star,e}{a'}{b'}\land
        w=\FormalDifferenceClass{M,\star,e}{c'}{d'}\land
        v=\FormalDifferenceClass{M,\star,e}{a'\star c'}{b'\star d'}
        \vdash u=v
    \end{aligned}
  }{%
    \CommutativeMonoidAlg{M}{\star}{e}\dsep
    \mathsf{K\ddot{u}rz}(M,\star)\dsep
    a,b,c,d,a',b',c',d'\in M\dsep
    z=\FormalDifferenceClass{M,\star,e}{a}{b}\land
    w=\FormalDifferenceClass{M,\star,e}{c}{d}\land
    u=\FormalDifferenceClass{M,\star,e}{a\star c}{b\star d}\dsep
    z=\FormalDifferenceClass{M,\star,e}{a'}{b'}\land
    w=\FormalDifferenceClass{M,\star,e}{c'}{d'}\land
    v=\FormalDifferenceClass{M,\star,e}{a'\star c'}{b'\star d'}
    \vdash u=v
  }[GroupCompletionOperationDescentFixedUniqueness]
    \proofstepwidestar[1]{\CommutativeMonoidAlg{M}{\star}{e}}{\rA}
    \proofstepwidestar[2]{\mathsf{K\ddot{u}rz}(M,\star)}{\rA}
    \proofstepwidestar[3]{a,b,c,d,a',b',c',d'\in M}{\rA}
    \proofstepwidestar[4]{%
      \begin{aligned}[t]
        &z=\FormalDifferenceClass{M,\star,e}{a}{b}\land
          w=\FormalDifferenceClass{M,\star,e}{c}{d}\land{}\\[-2pt]
        &u=\FormalDifferenceClass{M,\star,e}{a\star c}{b\star d}
      \end{aligned}}{\rA}
    \proofstepwidestar[5]{%
      \begin{aligned}[t]
        &z=\FormalDifferenceClass{M,\star,e}{a'}{b'}\land
          w=\FormalDifferenceClass{M,\star,e}{c'}{d'}\land{}\\[-2pt]
        &v=\FormalDifferenceClass{M,\star,e}{a'\star c'}{b'\star d'}
      \end{aligned}}{\rA}
    \proofstepwidestar[4,5]{%
      \begin{aligned}[t]
        &\FormalDifferenceClass{M,\star,e}{a}{b}={}\\[-2pt]
        &\quad\FormalDifferenceClass{M,\star,e}{a'}{b'}
      \end{aligned}}{%
      \rIE{\rAEa{4},\rAEa{5}}}
    \proofstepwidestar[4,5]{%
      \begin{aligned}[t]
        &\FormalDifferenceClass{M,\star,e}{c}{d}={}\\[-2pt]
        &\quad\FormalDifferenceClass{M,\star,e}{c'}{d'}
      \end{aligned}}{%
      \rIE{\rAEa{\rAEb{4}},\rAEa{\rAEb{5}}}}
    \proofstepwidestar[1,2,3,4,5]{%
      \begin{aligned}[t]
        &\FormalDifferenceClass{M,\star,e}{a\star c}{b\star d}={}\\[-2pt]
        &\quad\FormalDifferenceClass{M,\star,e}{a'\star c'}{b'\star d'}
      \end{aligned}}{%
      \FormulaRefAuto{GroupCompletionOperationCompatible}{1,2,3,6,7}}
    \proofstepwidestar[1,2,3,4,5]{u=v}{%
      \rIE{\rAEb{\rAEb{4}},\rAEb{\rAEb{5}},8}}
  \closeproofpartwideindR

  \proofpartwideindR[Eindeutigkeit beliebiger Ergebniszeugen]{%
    \begin{aligned}[t]
      &\CommutativeMonoidAlg{M}{\star}{e}\dsep
        \mathsf{K\ddot{u}rz}(M,\star)\dsep{}\\[-2pt]
      &u\in\GroupCompletionCarrier{M,\star,e}\land
        \exists a,b,c,d\in M\,\Bigl(
        z=\FormalDifferenceClass{M,\star,e}{a}{b}\land{}\\[-2pt]
      &\qquad w=\FormalDifferenceClass{M,\star,e}{c}{d}\land
        u=\FormalDifferenceClass{M,\star,e}{a\star c}{b\star d}\Bigr)
        \dsep{}\\[-2pt]
      &v\in\GroupCompletionCarrier{M,\star,e}\land
        \exists a',b',c',d'\in M\,\Bigl(
        z=\FormalDifferenceClass{M,\star,e}{a'}{b'}\land{}\\[-2pt]
      &\qquad w=\FormalDifferenceClass{M,\star,e}{c'}{d'}\land
        v=\FormalDifferenceClass{M,\star,e}{a'\star c'}{b'\star d'}\Bigr)
        \vdash u=v
    \end{aligned}
  }{%
    \begin{aligned}[t]
      &\CommutativeMonoidAlg{M}{\star}{e}\dsep
        \mathsf{K\ddot{u}rz}(M,\star)\dsep{}\\[-2pt]
      &u\in\GroupCompletionCarrier{M,\star,e}\land
        \exists a,b,c,d\in M\,\Bigl(
        z=\FormalDifferenceClass{M,\star,e}{a}{b}\land{}\\[-2pt]
      &\qquad w=\FormalDifferenceClass{M,\star,e}{c}{d}\land
        u=\FormalDifferenceClass{M,\star,e}{a\star c}{b\star d}\Bigr)
        \dsep{}\\[-2pt]
      &v\in\GroupCompletionCarrier{M,\star,e}\land
        \exists a',b',c',d'\in M\,\Bigl(
        z=\FormalDifferenceClass{M,\star,e}{a'}{b'}\land{}\\[-2pt]
      &\qquad w=\FormalDifferenceClass{M,\star,e}{c'}{d'}\land
        v=\FormalDifferenceClass{M,\star,e}{a'\star c'}{b'\star d'}\Bigr)
        \vdash u=v
    \end{aligned}
  }[GroupCompletionOperationDescentUniqueness]
    \proofstepwidestar[1]{\CommutativeMonoidAlg{M}{\star}{e}}{\rA}
    \proofstepwidestar[2]{\mathsf{K\ddot{u}rz}(M,\star)}{\rA}
    \proofstepwidestar[3]{%
      \begin{aligned}[t]
        &u\in\GroupCompletionCarrier{M,\star,e}\land
          \exists a,b,c,d\in M\,\Bigl(
          z=\FormalDifferenceClass{M,\star,e}{a}{b}\land{}\\[-2pt]
        &\qquad w=\FormalDifferenceClass{M,\star,e}{c}{d}\land{}\\[-2pt]
        &\qquad
          u=\FormalDifferenceClass{M,\star,e}{a\star c}{b\star d}\Bigr)
      \end{aligned}}{\rA}
    \proofstepwidestar[4]{%
      \begin{aligned}[t]
        &v\in\GroupCompletionCarrier{M,\star,e}\land
          \exists a',b',c',d'\in M\,\Bigl(
          z=\FormalDifferenceClass{M,\star,e}{a'}{b'}\land{}\\[-2pt]
        &\qquad w=\FormalDifferenceClass{M,\star,e}{c'}{d'}\land{}\\[-2pt]
        &\qquad
          v=\FormalDifferenceClass{M,\star,e}{a'\star c'}{b'\star d'}\Bigr)
      \end{aligned}}{\rA}
    \proofstepwidestar[5]{%
      \begin{aligned}[t]
        &a,b,c,d\in M\land{}\\[-2pt]
        &z=\FormalDifferenceClass{M,\star,e}{a}{b}\land
          w=\FormalDifferenceClass{M,\star,e}{c}{d}\land{}\\[-2pt]
        &\quad
          u=\FormalDifferenceClass{M,\star,e}{a\star c}{b\star d}
      \end{aligned}}{\rA}
    \proofstepwidestar[6]{%
      \begin{aligned}[t]
        &a',b',c',d'\in M\land{}\\[-2pt]
        &z=\FormalDifferenceClass{M,\star,e}{a'}{b'}\land
          w=\FormalDifferenceClass{M,\star,e}{c'}{d'}\land{}\\[-2pt]
        &\quad
          v=\FormalDifferenceClass{M,\star,e}{a'\star c'}{b'\star d'}
      \end{aligned}}{\rA}
    \proofstepwidestar[1,2,5,6]{u=v}{%
      \BandXXIXProofRefStack{%
        \FormulaRefAuto{GroupCompletionOperationDescentFixedUniqueness}\\
        \left(1,2,5,6\right)}}
    \proofstepwidestar[1,2,3,4]{u=v}{%
      \rEE{\rAEb{3},\rAEb{4},5,6,7}}
  \closeproofpartwideindR

  \proofpartwide{Abschluss des Abstiegs}
    \proofstepwidestar[1]{\CommutativeMonoidAlg{M}{\star}{e}}{\rA}
    \proofstepwidestar[2]{\mathsf{K\ddot{u}rz}(M,\star)}{\rA}
    \proofstepwidestar[3]{z,w\in\GroupCompletionCarrier{M,\star,e}}{\rA}
    \proofstepwidestar[1,2,3]{%
      \begin{aligned}[t]
        &\exists u\in\GroupCompletionCarrier{M,\star,e}\,
          \exists p,q,r,s\in M\,\Bigl(
          z=\FormalDifferenceClass{M,\star,e}{p}{q}\land{}\\[-2pt]
        &\qquad w=\FormalDifferenceClass{M,\star,e}{r}{s}\land
          u=\FormalDifferenceClass{M,\star,e}{p\star r}{q\star s}\Bigr)
      \end{aligned}}{%
      \FormulaRefAuto{GroupCompletionOperationDescentExistence}{1,2,3}}
    \proofstepwidestar[5]{%
      \begin{aligned}[t]
        &u\in\GroupCompletionCarrier{M,\star,e}\land
          \exists p,q,r,s\in M\,\Bigl(
          z=\FormalDifferenceClass{M,\star,e}{p}{q}\land{}\\[-2pt]
        &\qquad w=\FormalDifferenceClass{M,\star,e}{r}{s}\land
          u=\FormalDifferenceClass{M,\star,e}{p\star r}{q\star s}\Bigr)
      \end{aligned}}{\rA}
    \proofstepwidestar[6]{%
      \begin{aligned}[t]
        &v\in\GroupCompletionCarrier{M,\star,e}\land
          \exists p,q,r,s\in M\,\Bigl(
          z=\FormalDifferenceClass{M,\star,e}{p}{q}\land{}\\[-2pt]
        &\qquad w=\FormalDifferenceClass{M,\star,e}{r}{s}\land
          v=\FormalDifferenceClass{M,\star,e}{p\star r}{q\star s}\Bigr)
      \end{aligned}}{\rA}
    \proofstepwidestar[1,2,5,6]{u=v}{%
      \FormulaRefAuto{GroupCompletionOperationDescentUniqueness}{1,2,5,6}}
    \proofstepwidestar[1,2,3]{%
      \begin{aligned}[t]
        &\exists!u\in\GroupCompletionCarrier{M,\star,e}\,
          \exists p,q,r,s\in M\,\Bigl(
          z=\FormalDifferenceClass{M,\star,e}{p}{q}\land{}\\[-2pt]
        &\qquad w=\FormalDifferenceClass{M,\star,e}{r}{s}\land
          u=\FormalDifferenceClass{M,\star,e}{p\star r}{q\star s}\Bigr)
      \end{aligned}}{%
      \UEI{4,5,6,7}}
  \closeproofpartwide
\end{tabproofsplitwide}

Der eindeutige Existenzsatz ist der Quotientenabstieg der
Differenzenoperation; die folgende Klassenregel bezeichnet dieses eindeutig
bestimmte Ergebnis.

\FormulaDefDeltaK[Operation auf formalen Differenzen]{%
  \begin{aligned}[t]
    &\CommutativeMonoidAlg{M}{\star}{e}\dsep
      \mathsf{K\ddot{u}rz}(M,\star)\dsep a,b,c,d\in M\vdash{}\\[-2pt]
    &\FormalDifferenceClass{M,\star,e}{a}{b}
      \GroupCompletionOp{M,\star,e}
      \FormalDifferenceClass{M,\star,e}{c}{d}
      \coloneqq
      \FormalDifferenceClass{M,\star,e}{a\star c}{b\star d}
  \end{aligned}
}{GroupCompletionOperationDef}{%
  \DeltaRow{Mengen}{M\dsep e\dsep a\dsep b\dsep c\dsep d}
  \DeltaRow{Funktionen}{\star\dsep\GroupCompletionOp{M,\star,e}}
  \DeltaRow{Neue Symbole}{\GroupCompletionOp{M,\star,e}}
}

\FormulaDefDeltaK[Neutrales Element der formalen Differenzen]{%
  \begin{aligned}[t]
    &\CommutativeMonoidAlg{M}{\star}{e}\dsep
      \mathsf{K\ddot{u}rz}(M,\star)\vdash{}\\[-2pt]
    &\GroupCompletionZero{M,\star,e}\coloneqq
      \FormalDifferenceClass{M,\star,e}{e}{e}
  \end{aligned}
}{GroupCompletionZeroDef}{%
  \DeltaRow{Mengen}{M\dsep e\dsep\GroupCompletionZero{M,\star,e}}
  \DeltaRow{Funktionen}{\star}
  \DeltaRow{Neue Symbole}{\GroupCompletionZero{M,\star,e}}
}

\FormulaThmDeltaK[Das neutrale Element liegt im Differenzenquotienten]{%
  \begin{aligned}[t]
    &\CommutativeMonoidAlg{M}{\star}{e}\dsep
      \mathsf{K\ddot{u}rz}(M,\star)\vdash{}\\[-2pt]
    &\GroupCompletionZero{M,\star,e}
      \in\GroupCompletionCarrier{M,\star,e}
  \end{aligned}
}{GroupCompletionZeroCarrier}{%
  \DeltaRow{Mengen}{M\dsep e}
  \DeltaRow{Funktionen}{\star}
}
\begin{tabproof}
  \proofstep{1}{\CommutativeMonoidAlg{M}{\star}{e}}{\rA}
  \proofstep{2}{\mathsf{K\ddot{u}rz}(M,\star)}{\rA}
  \proofstep{1}{\MonoidAlg{M}{\star}{e}}{%
    \FormulaRefAuto{CommutativeMonoidBaseMonoidAxiom}[ax]{1}}
  \proofstep{1}{e\in M}{%
    \FormulaRefAuto{MonoidIdentityCarrierAxiom}[ax]{3}}
  \proofstep{1,2}{\FormalDifferenceClass{M,\star,e}{e}{e}
    \in\GroupCompletionCarrier{M,\star,e}}{%
    \FormulaRefAuto{GroupCompletionClassInCarrier}{1,2,4,4}}
  \proofstep{1,2}{\GroupCompletionZero{M,\star,e}
    =\FormalDifferenceClass{M,\star,e}{e}{e}}{%
    \FormulaRefAuto{GroupCompletionZeroDef}[def]{1,2}}
  \proofstep{1,2}{\GroupCompletionZero{M,\star,e}
    \in\GroupCompletionCarrier{M,\star,e}}{\rIE{6,5}}
\end{tabproof}

\FormulaThmDeltaK[Abgeschlossenheit der Differenzenoperation]{%
  \begin{aligned}[t]
    &\CommutativeMonoidAlg{M}{\star}{e}\dsep
      \mathsf{K\ddot{u}rz}(M,\star)\dsep
      z,w\in\GroupCompletionCarrier{M,\star,e}\vdash{}\\[-2pt]
    &z\GroupCompletionOp{M,\star,e}w
      \in\GroupCompletionCarrier{M,\star,e}
  \end{aligned}
}{GroupCompletionOperationClosure}{%
  \DeltaRow{Mengen}{M\dsep e\dsep z\dsep w\dsep a\dsep b\dsep c\dsep d}
  \DeltaRow{Funktionen}{\star\dsep\GroupCompletionOp{M,\star,e}}
}
\begin{tabproofsplitwide}
  \proofpartwideindR[Abgeschlossenheit auf Repräsentanten]{%
    \begin{aligned}[t]
      &\CommutativeMonoidAlg{M}{\star}{e}\dsep
        \mathsf{K\ddot{u}rz}(M,\star)\dsep a,b,c,d\in M\dsep{}\\[-2pt]
      &z=\FormalDifferenceClass{M,\star,e}{a}{b}\dsep
        w=\FormalDifferenceClass{M,\star,e}{c}{d}\vdash{}\\[-2pt]
      &z\GroupCompletionOp{M,\star,e}w
        \in\GroupCompletionCarrier{M,\star,e}
    \end{aligned}
  }{%
    \CommutativeMonoidAlg{M}{\star}{e}\dsep
    \mathsf{K\ddot{u}rz}(M,\star)\dsep a,b,c,d\in M\dsep
    z=\FormalDifferenceClass{M,\star,e}{a}{b}\dsep
    w=\FormalDifferenceClass{M,\star,e}{c}{d}
    \vdash z\GroupCompletionOp{M,\star,e}w
      \in\GroupCompletionCarrier{M,\star,e}
  }[GroupCompletionOperationClosureOnRepresentatives]
    \proofstepwidestar[1]{\CommutativeMonoidAlg{M}{\star}{e}}{\rA}
    \proofstepwidestar[2]{\mathsf{K\ddot{u}rz}(M,\star)}{\rA}
    \proofstepwidestar[3]{a,b,c,d\in M}{\rA}
    \proofstepwidestar[4]{z=\FormalDifferenceClass{M,\star,e}{a}{b}}{\rA}
    \proofstepwidestar[5]{w=\FormalDifferenceClass{M,\star,e}{c}{d}}{\rA}
    \proofstepwidestar[1,2,3,4,5]{%
      z\GroupCompletionOp{M,\star,e}w
      =\FormalDifferenceClass{M,\star,e}{a\star c}{b\star d}}{%
      \FormulaRefAuto{GroupCompletionOperationDef}[def]{1,2,3,4,5}}
    \proofstepwidestar[1]{\CommutativeSemigroupAlg{M}{\star}}{%
      \FormulaRefAuto{CommutativeMonoidIsCommutativeSemigroup}{1}}
    \proofstepwidestar[1,3]{a\star c,b\star d\in M}{%
      \FormulaRefAuto{SemigroupClosureAxiom}[ax]{%
        \FormulaRefAuto{CommutativeSemigroupBaseAxiom}[ax]{7},3}}
    \proofstepwidestar[1,2,3]{%
      \FormalDifferenceClass{M,\star,e}{a\star c}{b\star d}
      \in\GroupCompletionCarrier{M,\star,e}}{%
      \FormulaRefAuto{GroupCompletionClassInCarrier}{1,2,8}}
    \proofstepwidestar[1,2,3,4,5]{%
      z\GroupCompletionOp{M,\star,e}w
      \in\GroupCompletionCarrier{M,\star,e}}{\rIE{6,9}}
  \closeproofpartwideindR

  \proofpartwide{Elimination der Repräsentanten}
    \proofstepwidestar[1]{\CommutativeMonoidAlg{M}{\star}{e}}{\rA}
    \proofstepwidestar[2]{\mathsf{K\ddot{u}rz}(M,\star)}{\rA}
    \proofstepwidestar[3]{z,w\in\GroupCompletionCarrier{M,\star,e}}{\rA}
    \proofstepwidestar[1,2,3]{\exists a,b\in M\,
      z=\FormalDifferenceClass{M,\star,e}{a}{b}}{%
      \FormulaRefAuto{GroupCompletionRepresentative}{1,2,\rAEa{3}}}
    \proofstepwidestar[1,2,3]{\exists c,d\in M\,
      w=\FormalDifferenceClass{M,\star,e}{c}{d}}{%
      \FormulaRefAuto{GroupCompletionRepresentative}{1,2,\rAEb{3}}}
    \proofstepwidestar[6]{a,b,c,d\in M\land
      z=\FormalDifferenceClass{M,\star,e}{a}{b}\land
      w=\FormalDifferenceClass{M,\star,e}{c}{d}}{\rA}
    \proofstepwidestar[1,2,6]{%
      z\GroupCompletionOp{M,\star,e}w
      \in\GroupCompletionCarrier{M,\star,e}}{%
      \FormulaRefAuto{GroupCompletionOperationClosureOnRepresentatives}{1,2,6}}
    \proofstepwidestar[1,2,3]{%
      z\GroupCompletionOp{M,\star,e}w
      \in\GroupCompletionCarrier{M,\star,e}}{\rEE{4,5,6,7}}
  \closeproofpartwide
\end{tabproofsplitwide}

\FormulaThmDeltaK[Assoziativität der Differenzenoperation]{%
  \begin{aligned}[t]
    &\CommutativeMonoidAlg{M}{\star}{e}\dsep
      \mathsf{K\ddot{u}rz}(M,\star)\dsep
      x,y,z\in\GroupCompletionCarrier{M,\star,e}\vdash{}\\[-2pt]
    &(x\GroupCompletionOp{M,\star,e}y)
      \GroupCompletionOp{M,\star,e}z
      =x\GroupCompletionOp{M,\star,e}
        (y\GroupCompletionOp{M,\star,e}z)
  \end{aligned}
}{GroupCompletionOperationAssociative}{%
  \DeltaRow{Mengen}{M\dsep e\dsep x\dsep y\dsep z\dsep
    a\dsep b\dsep c\dsep d\dsep f\dsep g}
  \DeltaRow{Funktionen}{\star\dsep\GroupCompletionOp{M,\star,e}}
}
\begin{tabproofsplitwide}
  \proofpartwideindR[Rechnung auf Differenzenklassen]{%
    \begin{aligned}[t]
      &\CommutativeMonoidAlg{M}{\star}{e}\dsep
        \mathsf{K\ddot{u}rz}(M,\star)\dsep a,b,c,d,f,g\in M\vdash{}\\[-2pt]
      &(\FormalDifferenceClass{M,\star,e}{a}{b}
        \GroupCompletionOp{M,\star,e}\FormalDifferenceClass{M,\star,e}{c}{d})
        \GroupCompletionOp{M,\star,e}\FormalDifferenceClass{M,\star,e}{f}{g}\\[-2pt]
      &\quad=
        \FormalDifferenceClass{M,\star,e}{a}{b}
        \GroupCompletionOp{M,\star,e}
        (\FormalDifferenceClass{M,\star,e}{c}{d}
        \GroupCompletionOp{M,\star,e}\FormalDifferenceClass{M,\star,e}{f}{g})
    \end{aligned}
  }{%
    \begin{aligned}[t]
      &\CommutativeMonoidAlg{M}{\star}{e}\dsep
        \mathsf{K\ddot{u}rz}(M,\star)\dsep a,b,c,d,f,g\in M\vdash{}\\[-2pt]
      &(\FormalDifferenceClass{M,\star,e}{a}{b}
        \GroupCompletionOp{M,\star,e}\FormalDifferenceClass{M,\star,e}{c}{d})
        \GroupCompletionOp{M,\star,e}\FormalDifferenceClass{M,\star,e}{f}{g}\\[-2pt]
      &\qquad=\FormalDifferenceClass{M,\star,e}{a}{b}
        \GroupCompletionOp{M,\star,e}
        (\FormalDifferenceClass{M,\star,e}{c}{d}
        \GroupCompletionOp{M,\star,e}\FormalDifferenceClass{M,\star,e}{f}{g})
    \end{aligned}
  }[GroupCompletionOperationAssociativeOnClasses]
    \proofstepwidestar[1]{\CommutativeMonoidAlg{M}{\star}{e}}{\rA}
    \proofstepwidestar[2]{\mathsf{K\ddot{u}rz}(M,\star)}{\rA}
    \proofstepwidestar[3]{a,b,c,d,f,g\in M}{\rA}
    \proofstepwidestar[1,2,3]{%
      \begin{aligned}[t]
        &(\FormalDifferenceClass{M,\star,e}{a}{b}
          \GroupCompletionOp{M,\star,e}\FormalDifferenceClass{M,\star,e}{c}{d})
          \GroupCompletionOp{M,\star,e}\FormalDifferenceClass{M,\star,e}{f}{g}\\[-2pt]
        &\qquad=
          \FormalDifferenceClass{M,\star,e}{(a\star c)\star f}{(b\star d)\star g}
      \end{aligned}}{%
      \FormulaRefAuto{GroupCompletionOperationDef}[def]{1,2,3}}
    \proofstepwidestar[1,2,3]{%
      \begin{aligned}[t]
        &\FormalDifferenceClass{M,\star,e}{a}{b}\GroupCompletionOp{M,\star,e}
          (\FormalDifferenceClass{M,\star,e}{c}{d}
          \GroupCompletionOp{M,\star,e}\FormalDifferenceClass{M,\star,e}{f}{g})\\[-2pt]
        &\qquad=\FormalDifferenceClass{M,\star,e}
          {a\star(c\star f)}{b\star(d\star g)}
      \end{aligned}}{%
      \FormulaRefAuto{GroupCompletionOperationDef}[def]{1,2,3}}
    \proofstepwidestar[1]{\MonoidAlg{M}{\star}{e}}{%
      \FormulaRefAuto{CommutativeMonoidBaseMonoidAxiom}[ax]{1}}
    \proofstepwidestar[1]{\SemigroupAlg{M}{\star}}{%
      \FormulaRefAuto{MonoidBaseSemigroupAxiom}[ax]{6}}
    \proofstepwidestar[1,3]{(a\star c)\star f=a\star(c\star f)}{%
      \FormulaRefAuto{SemigroupAssociativityAxiom}[ax]{7,3}}
    \proofstepwidestar[1,3]{(b\star d)\star g=b\star(d\star g)}{%
      \FormulaRefAuto{SemigroupAssociativityAxiom}[ax]{7,3}}
    \proofstepwidestar[1,2,3]{%
      \begin{aligned}[t]
        &\FormalDifferenceClass{M,\star,e}{(a\star c)\star f}{(b\star d)\star g}={}\\[-2pt]
        &\quad\FormalDifferenceClass{M,\star,e}
          {a\star(c\star f)}{b\star(d\star g)}
      \end{aligned}}{\rIE{8,9}}
    \proofstepwidestar[1,2,3]{%
      \begin{aligned}[t]
        &(\FormalDifferenceClass{M,\star,e}{a}{b}
          \GroupCompletionOp{M,\star,e}\FormalDifferenceClass{M,\star,e}{c}{d})
          \GroupCompletionOp{M,\star,e}\FormalDifferenceClass{M,\star,e}{f}{g}\\[-2pt]
        &\qquad=\FormalDifferenceClass{M,\star,e}{a}{b}
          \GroupCompletionOp{M,\star,e}
          (\FormalDifferenceClass{M,\star,e}{c}{d}
          \GroupCompletionOp{M,\star,e}\FormalDifferenceClass{M,\star,e}{f}{g})
      \end{aligned}}{%
      \rIE{4,5,10}}
  \closeproofpartwideindR

  \proofpartwide{Elimination der Repräsentanten}
    \proofstepwidestar[1]{\CommutativeMonoidAlg{M}{\star}{e}}{\rA}
    \proofstepwidestar[2]{\mathsf{K\ddot{u}rz}(M,\star)}{\rA}
    \proofstepwidestar[3]{x,y,z\in\GroupCompletionCarrier{M,\star,e}}{\rA}
    \proofstepwidestar[1,2,3]{%
      \begin{aligned}[t]
        &\exists a,b\in M\,{}\\[-2pt]
        &\quad x=\FormalDifferenceClass{M,\star,e}{a}{b}
      \end{aligned}}{%
      \FormulaRefAuto{GroupCompletionRepresentative}{1,2,\rAEa{3}}}
    \proofstepwidestar[1,2,3]{%
      \begin{aligned}[t]
        &\exists c,d\in M\,{}\\[-2pt]
        &\quad y=\FormalDifferenceClass{M,\star,e}{c}{d}
      \end{aligned}}{%
      \FormulaRefAuto{GroupCompletionRepresentative}{1,2,\rAEa{\rAEb{3}}}}
    \proofstepwidestar[1,2,3]{%
      \begin{aligned}[t]
        &\exists f,g\in M\,{}\\[-2pt]
        &\quad z=\FormalDifferenceClass{M,\star,e}{f}{g}
      \end{aligned}}{%
      \FormulaRefAuto{GroupCompletionRepresentative}{1,2,\rAEb{\rAEb{3}}}}
    \proofstepwidestar[7]{%
      \begin{aligned}[t]
        &a,b,c,d,f,g\in M\land
          x=\FormalDifferenceClass{M,\star,e}{a}{b}\land{}\\[-2pt]
        &y=\FormalDifferenceClass{M,\star,e}{c}{d}\land
          z=\FormalDifferenceClass{M,\star,e}{f}{g}
      \end{aligned}}{\rA}
    \proofstepwidestar[1,2,7]{%
      \begin{aligned}[t]
        &(x\GroupCompletionOp{M,\star,e}y)\GroupCompletionOp{M,\star,e}z={}\\[-2pt]
        &\quad x\GroupCompletionOp{M,\star,e}
          (y\GroupCompletionOp{M,\star,e}z)
      \end{aligned}}{%
      \rIE{7,\FormulaRefAuto{GroupCompletionOperationAssociativeOnClasses}{1,2,7}}}
    \proofstepwidestar[1,2,3]{%
      \begin{aligned}[t]
        &(x\GroupCompletionOp{M,\star,e}y)\GroupCompletionOp{M,\star,e}z={}\\[-2pt]
        &\quad x\GroupCompletionOp{M,\star,e}
          (y\GroupCompletionOp{M,\star,e}z)
      \end{aligned}}{%
      \rEE{4,5,6,7,8}}
  \closeproofpartwide
\end{tabproofsplitwide}

\FormulaThmDeltaK[Kommutativität der Differenzenoperation]{%
  \begin{aligned}[t]
    &\CommutativeMonoidAlg{M}{\star}{e}\dsep
      \mathsf{K\ddot{u}rz}(M,\star)\dsep
      z,w\in\GroupCompletionCarrier{M,\star,e}\vdash{}\\[-2pt]
    &z\GroupCompletionOp{M,\star,e}w
      =w\GroupCompletionOp{M,\star,e}z
  \end{aligned}
}{GroupCompletionOperationCommutative}{%
  \DeltaRow{Mengen}{M\dsep e\dsep z\dsep w\dsep a\dsep b\dsep c\dsep d}
  \DeltaRow{Funktionen}{\star\dsep\GroupCompletionOp{M,\star,e}}
}
\begin{tabproofsplitwide}
  \proofpartwideindR[Vertauschung auf Differenzenklassen]{%
    \begin{aligned}[t]
      &\CommutativeMonoidAlg{M}{\star}{e}\dsep
        \mathsf{K\ddot{u}rz}(M,\star)\dsep a,b,c,d\in M\vdash{}\\[-2pt]
      &\FormalDifferenceClass{M,\star,e}{a}{b}\GroupCompletionOp{M,\star,e}\FormalDifferenceClass{M,\star,e}{c}{d}
      =\FormalDifferenceClass{M,\star,e}{c}{d}\GroupCompletionOp{M,\star,e}\FormalDifferenceClass{M,\star,e}{a}{b}
    \end{aligned}
  }{%
    \CommutativeMonoidAlg{M}{\star}{e}\dsep
    \mathsf{K\ddot{u}rz}(M,\star)\dsep a,b,c,d\in M\vdash
    \FormalDifferenceClass{M,\star,e}{a}{b}\GroupCompletionOp{M,\star,e}\FormalDifferenceClass{M,\star,e}{c}{d}
    =\FormalDifferenceClass{M,\star,e}{c}{d}\GroupCompletionOp{M,\star,e}\FormalDifferenceClass{M,\star,e}{a}{b}
  }[GroupCompletionOperationCommutativeOnClasses]
    \proofstepwidestar[1]{\CommutativeMonoidAlg{M}{\star}{e}}{\rA}
    \proofstepwidestar[2]{\mathsf{K\ddot{u}rz}(M,\star)}{\rA}
    \proofstepwidestar[3]{a,b,c,d\in M}{\rA}
    \proofstepwidestar[1,2,3]{%
      \FormalDifferenceClass{M,\star,e}{a}{b}\GroupCompletionOp{M,\star,e}\FormalDifferenceClass{M,\star,e}{c}{d}
      =\FormalDifferenceClass{M,\star,e}{a\star c}{b\star d}}{%
      \FormulaRefAuto{GroupCompletionOperationDef}[def]{1,2,3}}
    \proofstepwidestar[1,2,3]{%
      \FormalDifferenceClass{M,\star,e}{c}{d}\GroupCompletionOp{M,\star,e}\FormalDifferenceClass{M,\star,e}{a}{b}
      =\FormalDifferenceClass{M,\star,e}{c\star a}{d\star b}}{%
      \FormulaRefAuto{GroupCompletionOperationDef}[def]{1,2,3}}
    \proofstepwidestar[1,3]{a\star c=c\star a}{%
      \FormulaRefAuto{CommutativeMonoidCommutativityAxiom}[ax]{1,3}}
    \proofstepwidestar[1,3]{b\star d=d\star b}{%
      \FormulaRefAuto{CommutativeMonoidCommutativityAxiom}[ax]{1,3}}
    \proofstepwidestar[1,2,3]{%
      \FormalDifferenceClass{M,\star,e}{a}{b}\GroupCompletionOp{M,\star,e}\FormalDifferenceClass{M,\star,e}{c}{d}
      =\FormalDifferenceClass{M,\star,e}{c}{d}\GroupCompletionOp{M,\star,e}\FormalDifferenceClass{M,\star,e}{a}{b}}{%
      \rIE{4,5,6,7}}
  \closeproofpartwideindR

  \proofpartwide{Elimination der Repräsentanten}
    \proofstepwidestar[1]{\CommutativeMonoidAlg{M}{\star}{e}}{\rA}
    \proofstepwidestar[2]{\mathsf{K\ddot{u}rz}(M,\star)}{\rA}
    \proofstepwidestar[3]{z,w\in\GroupCompletionCarrier{M,\star,e}}{\rA}
    \proofstepwidestar[1,2,3]{\exists a,b\in M\,z=\FormalDifferenceClass{M,\star,e}{a}{b}}{%
      \FormulaRefAuto{GroupCompletionRepresentative}{1,2,\rAEa{3}}}
    \proofstepwidestar[1,2,3]{\exists c,d\in M\,w=\FormalDifferenceClass{M,\star,e}{c}{d}}{%
      \FormulaRefAuto{GroupCompletionRepresentative}{1,2,\rAEb{3}}}
    \proofstepwidestar[6]{a,b,c,d\in M\land z=\FormalDifferenceClass{M,\star,e}{a}{b}\land
      w=\FormalDifferenceClass{M,\star,e}{c}{d}}{\rA}
    \proofstepwidestar[1,2,6]{%
      z\GroupCompletionOp{M,\star,e}w=w\GroupCompletionOp{M,\star,e}z}{%
      \rIE{6,\FormulaRefAuto{GroupCompletionOperationCommutativeOnClasses}{1,2,6}}}
    \proofstepwidestar[1,2,3]{%
      z\GroupCompletionOp{M,\star,e}w=w\GroupCompletionOp{M,\star,e}z}{%
      \rEE{4,5,6,7}}
  \closeproofpartwide
\end{tabproofsplitwide}

\FormulaThmDeltaK[Neutralitätsgesetze der Differenzenoperation]{%
  \begin{aligned}[t]
    &\CommutativeMonoidAlg{M}{\star}{e}\dsep
      \mathsf{K\ddot{u}rz}(M,\star)\dsep
      z\in\GroupCompletionCarrier{M,\star,e}\vdash{}\\[-2pt]
    &z\GroupCompletionOp{M,\star,e}\GroupCompletionZero{M,\star,e}=z
      \land
      \GroupCompletionZero{M,\star,e}\GroupCompletionOp{M,\star,e}z=z
  \end{aligned}
}{GroupCompletionOperationIdentity}{%
  \DeltaRow{Mengen}{M\dsep e\dsep z\dsep a\dsep b}
  \DeltaRow{Funktionen}{\star\dsep\GroupCompletionOp{M,\star,e}}
}
\begin{tabproofsplitwide}
  \proofpartwideindR[Rechtsneutralität auf einer Klasse]{%
    \begin{aligned}[t]
      &\CommutativeMonoidAlg{M}{\star}{e}\dsep
        \mathsf{K\ddot{u}rz}(M,\star)\dsep a,b\in M\vdash{}\\[-2pt]
      &\FormalDifferenceClass{M,\star,e}{a}{b}\GroupCompletionOp{M,\star,e}
        \GroupCompletionZero{M,\star,e}=\FormalDifferenceClass{M,\star,e}{a}{b}
    \end{aligned}
  }{%
    \CommutativeMonoidAlg{M}{\star}{e}\dsep
    \mathsf{K\ddot{u}rz}(M,\star)\dsep a,b\in M\vdash
    \FormalDifferenceClass{M,\star,e}{a}{b}\GroupCompletionOp{M,\star,e}
      \GroupCompletionZero{M,\star,e}=\FormalDifferenceClass{M,\star,e}{a}{b}
  }[GroupCompletionOperationRightIdentityOnClass]
    \proofstepwidestar[1]{\CommutativeMonoidAlg{M}{\star}{e}}{\rA}
    \proofstepwidestar[2]{\mathsf{K\ddot{u}rz}(M,\star)}{\rA}
    \proofstepwidestar[3]{a,b\in M}{\rA}
    \proofstepwidestar[1,2]{\GroupCompletionZero{M,\star,e}=\FormalDifferenceClass{M,\star,e}{e}{e}}{%
      \FormulaRefAuto{GroupCompletionZeroDef}[def]{1,2}}
    \proofstepwidestar[1,2,3]{%
      \FormalDifferenceClass{M,\star,e}{a}{b}\GroupCompletionOp{M,\star,e}
        \GroupCompletionZero{M,\star,e}
      =\FormalDifferenceClass{M,\star,e}{a\star e}{b\star e}}{%
      \FormulaRefAuto{GroupCompletionOperationDef}[def]{1,2,3,4}}
    \proofstepwidestar[1]{\MonoidAlg{M}{\star}{e}}{%
      \FormulaRefAuto{CommutativeMonoidBaseMonoidAxiom}[ax]{1}}
    \proofstepwidestar[1,3]{a\star e=a}{%
      \FormulaRefAuto{MonoidRightIdentityAxiom}[ax]{6,3}}
    \proofstepwidestar[1,3]{b\star e=b}{%
      \FormulaRefAuto{MonoidRightIdentityAxiom}[ax]{6,3}}
    \proofstepwidestar[1,2,3]{%
      \FormalDifferenceClass{M,\star,e}{a}{b}\GroupCompletionOp{M,\star,e}
        \GroupCompletionZero{M,\star,e}=\FormalDifferenceClass{M,\star,e}{a}{b}}{\rIE{5,7,8}}
  \closeproofpartwideindR

  \proofpartwide{Elimination des Repräsentanten}
    \proofstepwidestar[1]{\CommutativeMonoidAlg{M}{\star}{e}}{\rA}
    \proofstepwidestar[2]{\mathsf{K\ddot{u}rz}(M,\star)}{\rA}
    \proofstepwidestar[3]{z\in\GroupCompletionCarrier{M,\star,e}}{\rA}
    \proofstepwidestar[1,2,3]{\exists a,b\in M\,z=\FormalDifferenceClass{M,\star,e}{a}{b}}{%
      \FormulaRefAuto{GroupCompletionRepresentative}{1,2,3}}
    \proofstepwidestar[5]{a,b\in M\land z=\FormalDifferenceClass{M,\star,e}{a}{b}}{\rA}
    \proofstepwidestar[1,2,5]{%
      \begin{aligned}[t]
        &z\GroupCompletionOp{M,\star,e}\GroupCompletionZero{M,\star,e}={}\\[-2pt]
        &\quad z
      \end{aligned}}{%
      \rIE{5,\FormulaRefAuto{GroupCompletionOperationRightIdentityOnClass}{1,2,5}}}
    \proofstepwidestar[1,2,3]{%
      \begin{aligned}[t]
        &z\GroupCompletionOp{M,\star,e}\GroupCompletionZero{M,\star,e}={}\\[-2pt]
        &\quad z
      \end{aligned}}{%
      \rEE{4,5,6}}
    \proofstepwidestar[1,2,3]{%
      \begin{aligned}[t]
        &\GroupCompletionZero{M,\star,e}\GroupCompletionOp{M,\star,e}z={}\\[-2pt]
        &\quad z
      \end{aligned}}{%
      \rIE{\FormulaRefAuto{GroupCompletionOperationCommutative}{%
        1,2,\FormulaRefAuto{GroupCompletionZeroCarrier}{1,2},3},7}}
    \proofstepwidestar[1,2,3]{%
      \begin{aligned}[t]
        &z\GroupCompletionOp{M,\star,e}\GroupCompletionZero{M,\star,e}=z
          \land{}\\[-2pt]
        &\GroupCompletionZero{M,\star,e}\GroupCompletionOp{M,\star,e}z=z
      \end{aligned}}{%
      \rAI{7,8}}
  \closeproofpartwide
\end{tabproofsplitwide}

\FormulaThmDeltaK[Die formalen Differenzen bilden ein kommutatives Monoid]{%
  \begin{aligned}[t]
    &\CommutativeMonoidAlg{M}{\star}{e}\dsep
      \mathsf{K\ddot{u}rz}(M,\star)\vdash{}\\[-2pt]
    &\CommutativeMonoidAlg{\GroupCompletionCarrier{M,\star,e}}
      {\GroupCompletionOp{M,\star,e}}{\GroupCompletionZero{M,\star,e}}
  \end{aligned}
}{GroupCompletionCommutativeMonoid}{%
  \DeltaRow{Mengen}{M\dsep e\dsep z\dsep x\dsep y}
  \DeltaRow{Funktionen}{\star\dsep\GroupCompletionOp{M,\star,e}}
}
\begin{tabproofwide}
  \proofstepwidestar[1]{\CommutativeMonoidAlg{M}{\star}{e}}{\rA}
  \proofstepwidestar[2]{\mathsf{K\ddot{u}rz}(M,\star)}{\rA}
  \proofstepwidestar[1,2]{%
    \SemigroupAlg{\GroupCompletionCarrier{M,\star,e}}
      {\GroupCompletionOp{M,\star,e}}}{%
    \BandXXIXProofRefStack{%
      \FormulaRefAuto{SemigroupDef}[def]\\
      \left(
        \FormulaRefAuto{GroupCompletionOperationClosure}{1,2},
        \FormulaRefAuto{GroupCompletionOperationAssociative}{1,2}
      \right)}}
  \proofstepwidestar[1,2]{%
    \GroupCompletionZero{M,\star,e}
      \in\GroupCompletionCarrier{M,\star,e}}{%
    \FormulaRefAuto{GroupCompletionZeroCarrier}{1,2}}
  \proofstepwidestar[5]{z\in\GroupCompletionCarrier{M,\star,e}}{\rA}
  \proofstepwidestar[1,2,5]{%
    \begin{aligned}[t]
      &z\GroupCompletionOp{M,\star,e}\GroupCompletionZero{M,\star,e}=z
        \land{}\\[-2pt]
      &\GroupCompletionZero{M,\star,e}\GroupCompletionOp{M,\star,e}z=z
    \end{aligned}}{%
    \FormulaRefAuto{GroupCompletionOperationIdentity}{1,2,5}}
  \proofstepwidestar[1,2,5]{%
    \GroupCompletionZero{M,\star,e}\GroupCompletionOp{M,\star,e}z=z}{%
    \rAEb{6}}
  \proofstepwidestar[1,2,5]{%
    z\GroupCompletionOp{M,\star,e}\GroupCompletionZero{M,\star,e}=z}{%
    \rAEa{6}}
  \proofstepwidestar[1,2]{%
    \MonoidAlg{\GroupCompletionCarrier{M,\star,e}}
      {\GroupCompletionOp{M,\star,e}}{\GroupCompletionZero{M,\star,e}}}{%
    \BandXXIXProofRefStack{%
      \FormulaRefAuto{MonoidDef}[def]\\
      \left(3,4,7,8\right)}}
  \proofstepwidestar[10]{x\in\GroupCompletionCarrier{M,\star,e}}{\rA}
  \proofstepwidestar[11]{y\in\GroupCompletionCarrier{M,\star,e}}{\rA}
  \proofstepwidestar[1,2,10,11]{%
    x\GroupCompletionOp{M,\star,e}y
      =y\GroupCompletionOp{M,\star,e}x}{%
    \FormulaRefAuto{GroupCompletionOperationCommutative}{1,2,10,11}}
  \proofstepwidestar[1,2]{%
    \CommutativeMonoidAlg{\GroupCompletionCarrier{M,\star,e}}
      {\GroupCompletionOp{M,\star,e}}{\GroupCompletionZero{M,\star,e}}}{%
    \BandXXIXProofRefStack{%
      \FormulaRefAuto{CommutativeMonoidDef}[def]\\
      \left(9,12\right)}}
\end{tabproofwide}

\newpage
\section{Kanonische Einbettung}

\FormulaDefDeltaK[Kanonische Einbettung]{%
  \CommutativeMonoidAlg{M}{\star}{e}\dsep
  \mathsf{K\ddot{u}rz}(M,\star)\vdash
  \begin{aligned}[t]
    &\GroupCompletionEmbed{M,\star,e}\colon
      M\to\GroupCompletionCarrier{M,\star,e},\\[-2pt]
    &\forall a\in M\,
      \bigl(
        \GroupCompletionEmbed{M,\star,e}(a)\coloneqq
        \FormalDifferenceClass{M,\star,e}{a}{e}
      \bigr)
  \end{aligned}
}{GroupCompletionEmbeddingDef}{%
  \DeltaRow{Mengen}{M\dsep e\dsep a}
  \DeltaRow{Funktionen}{\star\dsep\GroupCompletionEmbed{M,\star,e}}
  \DeltaRow{Neue Symbole}{\GroupCompletionEmbed{M,\star,e}}
}
\begin{tabproof}
  \proofstep{1}{\CommutativeMonoidAlg{M}{\star}{e}}{\rA}
  \proofstep{2}{\mathsf{K\ddot{u}rz}(M,\star)}{\rA}
  \proofstep{3}{a\in M}{\rA}
  \proofstep{1}{\MonoidAlg{M}{\star}{e}}{%
    \FormulaRefAuto{CommutativeMonoidBaseMonoidAxiom}[ax]{1}}
  \proofstep{1}{e\in M}{%
    \FormulaRefAuto{MonoidIdentityCarrierAxiom}[ax]{4}}
  \proofstep{1,2,3}{\FormalDifferenceClass{M,\star,e}{a}{e}
    \in\GroupCompletionCarrier{M,\star,e}}{%
    \FormulaRefAuto{GroupCompletionClassInCarrier}{1,2,3,5}}
  \proofstep{1,2}{%
    \forall a\in M\,
      \FormalDifferenceClass{M,\star,e}{a}{e}
        \in\GroupCompletionCarrier{M,\star,e}}{%
    \rUI{\rRI{3,6}}}
\end{tabproof}

\FormulaThmDeltaK[Typisierung der kanonischen Einbettung]{%
  \begin{aligned}[t]
    &\CommutativeMonoidAlg{M}{\star}{e}\dsep
      \mathsf{K\ddot{u}rz}(M,\star)\vdash{}\\[-2pt]
    &\GroupCompletionEmbed{M,\star,e}\colon
      M\to\GroupCompletionCarrier{M,\star,e}
  \end{aligned}
}{GroupCompletionEmbeddingFunction}{%
  \DeltaRow{Mengen}{M\dsep e\dsep a}
  \DeltaRow{Funktionen}{\star\dsep\GroupCompletionEmbed{M,\star,e}}
}
\begin{tabproof}
  \proofstep{1}{\CommutativeMonoidAlg{M}{\star}{e}}{\rA}
  \proofstep{2}{\mathsf{K\ddot{u}rz}(M,\star)}{\rA}
  \proofstep{1,2}{\GroupCompletionEmbed{M,\star,e}\colon
    M\to\GroupCompletionCarrier{M,\star,e}}{%
    \FormulaRefAuto{GroupCompletionEmbeddingDef}[def]{1,2}}
\end{tabproof}

\FormulaThmDeltaK[Operationserhaltung der kanonischen Einbettung]{%
  \begin{aligned}[t]
    &\CommutativeMonoidAlg{M}{\star}{e}\dsep
      \mathsf{K\ddot{u}rz}(M,\star)\dsep a,b\in M\vdash{}\\[-2pt]
    &\GroupCompletionEmbed{M,\star,e}(a\star b)
      =\GroupCompletionEmbed{M,\star,e}(a)
        \GroupCompletionOp{M,\star,e}
        \GroupCompletionEmbed{M,\star,e}(b)
  \end{aligned}
}{GroupCompletionEmbeddingOperation}{%
  \DeltaRow{Mengen}{M\dsep e\dsep a\dsep b}
  \DeltaRow{Funktionen}{\star\dsep\GroupCompletionOp{M,\star,e}
    \dsep\GroupCompletionEmbed{M,\star,e}}
}
\begin{tabproofwide}
  \proofstepwidestar[1]{\CommutativeMonoidAlg{M}{\star}{e}}{\rA}
  \proofstepwidestar[2]{\mathsf{K\ddot{u}rz}(M,\star)}{\rA}
  \proofstepwidestar[3]{a,b\in M}{\rA}
  \proofstepwidestar[1]{\CommutativeSemigroupAlg{M}{\star}}{%
    \FormulaRefAuto{CommutativeMonoidIsCommutativeSemigroup}{1}}
  \proofstepwidestar[1,3]{a\star b\in M}{%
    \FormulaRefAuto{SemigroupClosureAxiom}[ax]{%
      \FormulaRefAuto{CommutativeSemigroupBaseAxiom}[ax]{4},3}}
  \proofstepwidestar[1]{\MonoidAlg{M}{\star}{e}}{%
    \FormulaRefAuto{CommutativeMonoidBaseMonoidAxiom}[ax]{1}}
  \proofstepwidestar[1]{e\in M}{%
    \FormulaRefAuto{MonoidIdentityCarrierAxiom}[ax]{6}}
  \proofstepwidestar[1,2]{%
    \forall x\in M\,
      \GroupCompletionEmbed{M,\star,e}(x)
        =\FormalDifferenceClass{M,\star,e}{x}{e}}{%
    \FormulaRefAuto{GroupCompletionEmbeddingDef}[def]{1,2}}
  \proofstepwidestar[1,2,3]{%
    \begin{aligned}[t]
      &\GroupCompletionEmbed{M,\star,e}(a\star b)={}\\[-2pt]
      &\quad\FormalDifferenceClass{M,\star,e}{a\star b}{e}
    \end{aligned}}{%
    \rRE{\rUE{8},5}}
  \proofstepwidestar[1,2,3]{%
    \GroupCompletionEmbed{M,\star,e}(a)
      =\FormalDifferenceClass{M,\star,e}{a}{e}}{%
    \rRE{\rUE{8},\rAEa{3}}}
  \proofstepwidestar[1,2,3]{%
    \GroupCompletionEmbed{M,\star,e}(b)
      =\FormalDifferenceClass{M,\star,e}{b}{e}}{%
    \rRE{\rUE{8},\rAEb{3}}}
  \proofstepwidestarbreakkeep[1,2,3]{%
    \begin{aligned}[t]
      &\GroupCompletionEmbed{M,\star,e}(a)
        \GroupCompletionOp{M,\star,e}
        \GroupCompletionEmbed{M,\star,e}(b)={}\\[-2pt]
      &\quad\FormalDifferenceClass{M,\star,e}{a\star b}{e\star e}
    \end{aligned}}{%
    \rIE{10,11,\FormulaRefAuto{GroupCompletionOperationDef}[def]{1,2,3}}}
  \proofstepwidestar[1]{e\star e=e}{%
    \FormulaRefAuto{MonoidRightIdentityAxiom}[ax]{6,7}}
  \proofstepwidestar[1,2,3]{%
    \begin{aligned}[t]
      &\FormalDifferenceClass{M,\star,e}{a\star b}{e\star e}={}\\[-2pt]
      &\quad\FormalDifferenceClass{M,\star,e}{a\star b}{e}
    \end{aligned}}{\rIE{13}}
  \proofstepwidestar[1,2,3]{%
    \begin{aligned}[t]
      &\GroupCompletionEmbed{M,\star,e}(a\star b)={}\\[-2pt]
      &\quad\GroupCompletionEmbed{M,\star,e}(a)
        \GroupCompletionOp{M,\star,e}
        \GroupCompletionEmbed{M,\star,e}(b)
    \end{aligned}}{\rIE{9,12,14}}
\end{tabproofwide}

\FormulaThmDeltaK[Neutralenerhaltung der kanonischen Einbettung]{%
  \begin{aligned}[t]
    &\CommutativeMonoidAlg{M}{\star}{e}\dsep
      \mathsf{K\ddot{u}rz}(M,\star)\vdash{}\\[-2pt]
    &\GroupCompletionEmbed{M,\star,e}(e)
      =\GroupCompletionZero{M,\star,e}
  \end{aligned}
}{GroupCompletionEmbeddingIdentity}{%
  \DeltaRow{Mengen}{M\dsep e}
  \DeltaRow{Funktionen}{\star\dsep\GroupCompletionEmbed{M,\star,e}}
}
\begin{tabproof}
  \proofstep{1}{\CommutativeMonoidAlg{M}{\star}{e}}{\rA}
  \proofstep{2}{\mathsf{K\ddot{u}rz}(M,\star)}{\rA}
  \proofstep{1}{\MonoidAlg{M}{\star}{e}}{%
    \FormulaRefAuto{CommutativeMonoidBaseMonoidAxiom}[ax]{1}}
  \proofstep{1}{e\in M}{\FormulaRefAuto{MonoidIdentityCarrierAxiom}[ax]{3}}
  \proofstep{1,2}{%
    \forall a\in M\,
      \GroupCompletionEmbed{M,\star,e}(a)
        =\FormalDifferenceClass{M,\star,e}{a}{e}}{%
    \FormulaRefAuto{GroupCompletionEmbeddingDef}[def]{1,2}}
  \proofstep{1,2}{\GroupCompletionEmbed{M,\star,e}(e)
    =\FormalDifferenceClass{M,\star,e}{e}{e}}{%
    \rRE{\rUE{5},4}}
  \proofstep{1,2}{\GroupCompletionZero{M,\star,e}
    =\FormalDifferenceClass{M,\star,e}{e}{e}}{%
    \FormulaRefAuto{GroupCompletionZeroDef}[def]{1,2}}
  \proofstep{1,2}{\GroupCompletionEmbed{M,\star,e}(e)
    =\GroupCompletionZero{M,\star,e}}{\rIE{6,7}}
\end{tabproof}

\FormulaThmDeltaK[Die kanonische Einbettung ist ein Monoidhomomorphismus]{%
  \begin{aligned}[t]
    &\CommutativeMonoidAlg{M}{\star}{e}\dsep
      \mathsf{K\ddot{u}rz}(M,\star)\vdash{}\\[-2pt]
    &\MonoidHom{\GroupCompletionEmbed{M,\star,e}}{M,\star,e}
      {\GroupCompletionCarrier{M,\star,e},
       \GroupCompletionOp{M,\star,e},\GroupCompletionZero{M,\star,e}}
  \end{aligned}
}{GroupCompletionEmbeddingMonoidHom}{%
  \DeltaRow{Mengen}{M\dsep e\dsep a\dsep b}
  \DeltaRow{Funktionen}{\star\dsep\GroupCompletionOp{M,\star,e}
    \dsep\GroupCompletionEmbed{M,\star,e}}
}
\begin{tabproofwide}
  \proofstepwidestar[1]{\CommutativeMonoidAlg{M}{\star}{e}}{\rA}
  \proofstepwidestar[2]{\mathsf{K\ddot{u}rz}(M,\star)}{\rA}
  \proofstepwidestar[1]{\MonoidAlg{M}{\star}{e}}{%
    \FormulaRefAuto{CommutativeMonoidBaseMonoidAxiom}[ax]{1}}
  \proofstepwidestar[1,2]{%
    \CommutativeMonoidAlg{\GroupCompletionCarrier{M,\star,e}}
      {\GroupCompletionOp{M,\star,e}}{\GroupCompletionZero{M,\star,e}}}{%
    \FormulaRefAuto{GroupCompletionCommutativeMonoid}{1,2}}
  \proofstepwidestar[1,2]{%
    \MonoidAlg{\GroupCompletionCarrier{M,\star,e}}
      {\GroupCompletionOp{M,\star,e}}{\GroupCompletionZero{M,\star,e}}}{%
    \FormulaRefAuto{CommutativeMonoidBaseMonoidAxiom}[ax]{4}}
  \proofstepwidestar[1]{\SemigroupAlg{M}{\star}}{%
    \FormulaRefAuto{MonoidBaseSemigroupAxiom}[ax]{3}}
  \proofstepwidestar[1,2]{%
    \SemigroupAlg{\GroupCompletionCarrier{M,\star,e}}
      {\GroupCompletionOp{M,\star,e}}}{%
    \FormulaRefAuto{MonoidBaseSemigroupAxiom}[ax]{5}}
  \proofstepwidestar[1,2]{%
    \GroupCompletionEmbed{M,\star,e}\colon
      M\to\GroupCompletionCarrier{M,\star,e}}{%
    \FormulaRefAuto{GroupCompletionEmbeddingFunction}{1,2}}
  \proofstepwidestar[9]{a\in M}{\rA}
  \proofstepwidestar[10]{b\in M}{\rA}
  \proofstepwidestar[1,2,9,10]{%
    \GroupCompletionEmbed{M,\star,e}(a\star b)
      =\GroupCompletionEmbed{M,\star,e}(a)
        \GroupCompletionOp{M,\star,e}
        \GroupCompletionEmbed{M,\star,e}(b)}{%
    \FormulaRefAuto{GroupCompletionEmbeddingOperation}{1,2,9,10}}
  \proofstepwidestar[1,2]{%
    \SemigroupHom{\GroupCompletionEmbed{M,\star,e}}{M,\star}
      {\GroupCompletionCarrier{M,\star,e},
       \GroupCompletionOp{M,\star,e}}}{%
    \BandXXIXProofRefStack{%
      \FormulaRefAuto{SemigroupHomDef}[def]\\
      \left(6,7,8,11\right)}}
  \proofstepwidestar[1,2]{\GroupCompletionEmbed{M,\star,e}(e)
    =\GroupCompletionZero{M,\star,e}}{%
    \FormulaRefAuto{GroupCompletionEmbeddingIdentity}{1,2}}
  \proofstepwidestarbreakkeep[1,2]{%
    \MonoidHom{\GroupCompletionEmbed{M,\star,e}}{M,\star,e}
      {\GroupCompletionCarrier{M,\star,e},
       \GroupCompletionOp{M,\star,e},\GroupCompletionZero{M,\star,e}}}{%
    \BandXXIXProofRefStack{%
      \FormulaRefAuto{MonoidHomDef}[def]\\
      \left(3,5,12,13\right)}}
\end{tabproofwide}

\FormulaThmDeltaK[Injektivität der kanonischen Einbettung]{%
  \begin{aligned}[t]
    &\CommutativeMonoidAlg{M}{\star}{e}\dsep
      \mathsf{K\ddot{u}rz}(M,\star)\vdash{}\\[-2pt]
    &\GroupCompletionEmbed{M,\star,e}\colon
      M\inj\GroupCompletionCarrier{M,\star,e}
  \end{aligned}
}{GroupCompletionEmbeddingInjective}{%
  \DeltaRow{Mengen}{M\dsep e\dsep a\dsep b}
  \DeltaRow{Funktionen}{\star\dsep\GroupCompletionEmbed{M,\star,e}}
}
\begin{tabproofsplitwide}
  \proofpartwideindR[Gleichheitsreflexion]{%
    \begin{aligned}[t]
      &\CommutativeMonoidAlg{M}{\star}{e}\dsep
        \mathsf{K\ddot{u}rz}(M,\star)\dsep a,b\in M\dsep{}\\[-2pt]
      &\GroupCompletionEmbed{M,\star,e}(a)
        =\GroupCompletionEmbed{M,\star,e}(b)\vdash a=b
    \end{aligned}
  }{%
    \CommutativeMonoidAlg{M}{\star}{e}\dsep
    \mathsf{K\ddot{u}rz}(M,\star)\dsep a,b\in M\dsep
    \GroupCompletionEmbed{M,\star,e}(a)
      =\GroupCompletionEmbed{M,\star,e}(b)\vdash a=b
  }[GroupCompletionEmbeddingReflectsEquality]
    \proofstepwidestar[1]{\CommutativeMonoidAlg{M}{\star}{e}}{\rA}
    \proofstepwidestar[2]{\mathsf{K\ddot{u}rz}(M,\star)}{\rA}
    \proofstepwidestar[3]{a,b\in M}{\rA}
    \proofstepwidestar[4]{\GroupCompletionEmbed{M,\star,e}(a)
      =\GroupCompletionEmbed{M,\star,e}(b)}{\rA}
    \proofstepwidestar[1,2]{%
      \forall x\in M\,
        \GroupCompletionEmbed{M,\star,e}(x)
          =\FormalDifferenceClass{M,\star,e}{x}{e}}{%
      \FormulaRefAuto{GroupCompletionEmbeddingDef}[def]{1,2}}
    \proofstepwidestar[1,2,3]{%
      \GroupCompletionEmbed{M,\star,e}(a)
        =\FormalDifferenceClass{M,\star,e}{a}{e}}{%
      \rRE{\rUE{5},\rAEa{3}}}
    \proofstepwidestar[1,2,3]{%
      \GroupCompletionEmbed{M,\star,e}(b)
        =\FormalDifferenceClass{M,\star,e}{b}{e}}{%
      \rRE{\rUE{5},\rAEb{3}}}
    \proofstepwidestar[1,2,3,4]{%
      \FormalDifferenceClass{M,\star,e}{a}{e}
        =\FormalDifferenceClass{M,\star,e}{b}{e}}{%
      \rIE{6,4,7}}
    \proofstepwidestar[1]{\MonoidAlg{M}{\star}{e}}{%
      \FormulaRefAuto{CommutativeMonoidBaseMonoidAxiom}[ax]{1}}
    \proofstepwidestar[1]{e\in M}{%
      \FormulaRefAuto{MonoidIdentityCarrierAxiom}[ax]{9}}
    \proofstepwidestar[1,2,3,4]{a\star e=e\star b}{%
      \FormulaRefAuto{GroupCompletionClassEquality}{1,2,3,10,8}}
    \proofstepwidestar[1,3]{a\star e=a}{%
      \FormulaRefAuto{MonoidRightIdentityAxiom}[ax]{9,3}}
    \proofstepwidestar[1,3]{e\star b=b}{%
      \FormulaRefAuto{MonoidLeftIdentityAxiom}[ax]{9,3}}
    \proofstepwidestar[1,2,3,4]{a=b}{\rIE{12,13,11}}
  \closeproofpartwideindR

  \newpage
  \proofpartwide{Schluss}
    \proofstepwidestar[1]{\CommutativeMonoidAlg{M}{\star}{e}}{\rA}
    \proofstepwidestar[2]{\mathsf{K\ddot{u}rz}(M,\star)}{\rA}
    \proofstepwidestar[1,2]{\GroupCompletionEmbed{M,\star,e}\colon
      M\to\GroupCompletionCarrier{M,\star,e}}{%
      \FormulaRefAuto{GroupCompletionEmbeddingFunction}{1,2}}
    \proofstepwidestar[4]{a\in M}{\rA}
    \proofstepwidestar[5]{b\in M}{\rA}
    \proofstepwidestar[6]{\GroupCompletionEmbed{M,\star,e}(a)
      =\GroupCompletionEmbed{M,\star,e}(b)}{\hspace{0.8em}\rA}
    \proofstepwidestar[1,2,4,5,6]{a=b}{%
      \FormulaRefAuto{GroupCompletionEmbeddingReflectsEquality}{1,2,4,5,6}}
    \proofstepwidestarbreakkeep[1,2]{\GroupCompletionEmbed{M,\star,e}\colon
      M\inj\GroupCompletionCarrier{M,\star,e}}{%
      \FormulaRefAuto{Injektive Funktion}{3,\rUI{\rRI{4,5,6,7}}}}
  \closeproofpartwide
\end{tabproofsplitwide}

\section{Inverse Klassen und der Einbettungssatz}

\FormulaThmDeltaK[Vertauschte formale Differenzen sind inverse Klassen]{%
  \begin{aligned}[t]
    &\CommutativeMonoidAlg{M}{\star}{e}\dsep
      \mathsf{K\ddot{u}rz}(M,\star)\dsep a,b\in M\vdash{}\\[-2pt]
    &\FormalDifferenceClass{M,\star,e}{a}{b}
       \GroupCompletionOp{M,\star,e}
       \FormalDifferenceClass{M,\star,e}{b}{a}
       =\GroupCompletionZero{M,\star,e}\land{}\\[-2pt]
    &\FormalDifferenceClass{M,\star,e}{b}{a}
       \GroupCompletionOp{M,\star,e}
       \FormalDifferenceClass{M,\star,e}{a}{b}
       =\GroupCompletionZero{M,\star,e}
  \end{aligned}
}{GroupCompletionClassInverse}{%
  \DeltaRow{Mengen}{M\dsep e\dsep a\dsep b}
  \DeltaRow{Funktionen}{\star\dsep\GroupCompletionOp{M,\star,e}}
}
\begin{tabproofwide}
  \proofstepwidestar[1]{\CommutativeMonoidAlg{M}{\star}{e}}{\rA}
  \proofstepwidestar[2]{\mathsf{K\ddot{u}rz}(M,\star)}{\rA}
  \proofstepwidestar[3]{a,b\in M}{\rA}
  \proofstepwidestar[1]{\MonoidAlg{M}{\star}{e}}{%
    \FormulaRefAuto{CommutativeMonoidBaseMonoidAxiom}[ax]{1}}
  \proofstepwidestar[1]{\SemigroupAlg{M}{\star}}{%
    \FormulaRefAuto{MonoidBaseSemigroupAxiom}[ax]{4}}
  \proofstepwidestar[1,3]{a\star b,b\star a\in M}{%
    \FormulaRefAuto{SemigroupClosureAxiom}[ax]{5,3}}
  \proofstepwidestar[1]{e\in M}{%
    \FormulaRefAuto{MonoidIdentityCarrierAxiom}[ax]{4}}
  \proofstepwide[1,3]{(a\star b)\star e}{=}{a\star b}{%
    \FormulaRefAuto{MonoidRightIdentityAxiom}[ax]{4,6}}
  \proofstepwide[1,3]{}{=}{b\star a}{%
    \FormulaRefAuto{CommutativeMonoidCommutativityAxiom}[ax]{1,3}}
  \proofstepwide[1,3]{}{=}{(b\star a)\star e}{%
    \FormulaRefAuto{a=b\vdash b=a}{%
      \FormulaRefAuto{MonoidRightIdentityAxiom}[ax]{4,6}}}
  \proofstepwidestar[1,2,3]{%
    \begin{aligned}[t]
      &\FormalDifferenceClass{M,\star,e}{a\star b}{b\star a}={}\\[-2pt]
      &\quad\FormalDifferenceClass{M,\star,e}{e}{e}
    \end{aligned}}{%
    \FormulaRefAuto{GroupCompletionClassEquality}{1,2,6,7,\rChain{8,10}}}
  \proofstepwidestar[1,2]{\GroupCompletionZero{M,\star,e}
    =\FormalDifferenceClass{M,\star,e}{e}{e}}{%
    \FormulaRefAuto{GroupCompletionZeroDef}[def]{1,2}}
  \proofstepwidestar[1,2,3]{%
    \begin{aligned}[t]
      &\FormalDifferenceClass{M,\star,e}{a}{b}
        \GroupCompletionOp{M,\star,e}
        \FormalDifferenceClass{M,\star,e}{b}{a}={}\\[-2pt]
      &\quad\FormalDifferenceClass{M,\star,e}{a\star b}{b\star a}
    \end{aligned}}{%
    \FormulaRefAuto{GroupCompletionOperationDef}[def]{1,2,3}}
  \proofstepwidestar[1,2,3]{%
    \begin{aligned}[t]
      &\FormalDifferenceClass{M,\star,e}{a}{b}
        \GroupCompletionOp{M,\star,e}
        \FormalDifferenceClass{M,\star,e}{b}{a}={}\\[-2pt]
      &\quad\GroupCompletionZero{M,\star,e}
    \end{aligned}}{\rIE{13,11,12}}
  \proofstepwidestar[1,2,3]{%
    \begin{aligned}[t]
      &\FormalDifferenceClass{M,\star,e}{b}{a}
        \GroupCompletionOp{M,\star,e}
        \FormalDifferenceClass{M,\star,e}{a}{b}={}\\[-2pt]
      &\quad\FormalDifferenceClass{M,\star,e}{a}{b}
        \GroupCompletionOp{M,\star,e}
        \FormalDifferenceClass{M,\star,e}{b}{a}
    \end{aligned}}{%
    \BandXXIXProofRefStack{%
      \FormulaRefAuto{GroupCompletionOperationCommutative}\\
      \left(
        1,2,\FormulaRefAuto{GroupCompletionClassInCarrier}{1,2,3},
        \FormulaRefAuto{GroupCompletionClassInCarrier}{1,2,3}
      \right)}}
  \proofstepwidestar[1,2,3]{%
    \begin{aligned}[t]
      &\FormalDifferenceClass{M,\star,e}{b}{a}
        \GroupCompletionOp{M,\star,e}
        \FormalDifferenceClass{M,\star,e}{a}{b}={}\\[-2pt]
      &\quad\GroupCompletionZero{M,\star,e}
    \end{aligned}}{\rIE{15,14}}
  \proofstepwidestar[1,2,3]{\begin{aligned}[t]
    &\FormalDifferenceClass{M,\star,e}{a}{b}
      \GroupCompletionOp{M,\star,e}
      \FormalDifferenceClass{M,\star,e}{b}{a}={}\\[-2pt]
    &\quad\GroupCompletionZero{M,\star,e}\land{}\\[-2pt]
    &\FormalDifferenceClass{M,\star,e}{b}{a}
      \GroupCompletionOp{M,\star,e}
      \FormalDifferenceClass{M,\star,e}{a}{b}={}\\[-2pt]
    &\quad\GroupCompletionZero{M,\star,e}
  \end{aligned}}{\rAI{14,16}}
\end{tabproofwide}

\FormulaThmDeltaK[Jede Differenzenklasse besitzt ein beidseitiges Inverses]{%
  \begin{aligned}[t]
    &\CommutativeMonoidAlg{M}{\star}{e}\dsep
      \mathsf{K\ddot{u}rz}(M,\star)\dsep
      z\in\GroupCompletionCarrier{M,\star,e}\vdash{}\\[-2pt]
    &\exists w\in\GroupCompletionCarrier{M,\star,e}\,
      \bigl(z\GroupCompletionOp{M,\star,e}w=\GroupCompletionZero{M,\star,e}
      \land w\GroupCompletionOp{M,\star,e}z=\GroupCompletionZero{M,\star,e}\bigr)
  \end{aligned}
}{GroupCompletionInverseExistence}{%
  \DeltaRow{Mengen}{M\dsep e\dsep z\dsep w\dsep a\dsep b}
  \DeltaRow{Funktionen}{\star\dsep\GroupCompletionOp{M,\star,e}}
}
\begin{tabproofwide}
  \proofstepwidestar[1]{\CommutativeMonoidAlg{M}{\star}{e}}{\rA}
  \proofstepwidestar[2]{\mathsf{K\ddot{u}rz}(M,\star)}{\rA}
  \proofstepwidestar[3]{z\in\GroupCompletionCarrier{M,\star,e}}{\rA}
  \proofstepwidestar[1,2,3]{\exists a,b\in M\,
    z=\FormalDifferenceClass{M,\star,e}{a}{b}}{%
    \FormulaRefAuto{GroupCompletionRepresentative}{1,2,3}}
  \proofstepwidestar[5]{a,b\in M\land
    z=\FormalDifferenceClass{M,\star,e}{a}{b}}{\rA}
  \proofstepwidestar[1,2,5]{\FormalDifferenceClass{M,\star,e}{b}{a}
    \in\GroupCompletionCarrier{M,\star,e}}{%
    \FormulaRefAuto{GroupCompletionClassInCarrier}{1,2,5}}
  \proofstepwidestar[1,2,5]{\begin{aligned}[t]
    &z\GroupCompletionOp{M,\star,e}\FormalDifferenceClass{M,\star,e}{b}{a}={}\\[-2pt]
    &\quad\GroupCompletionZero{M,\star,e}\land{}\\[-2pt]
    &\FormalDifferenceClass{M,\star,e}{b}{a}
      \GroupCompletionOp{M,\star,e}z={}\\[-2pt]
    &\quad\GroupCompletionZero{M,\star,e}
  \end{aligned}}{%
    \rIE{5,\FormulaRefAuto{GroupCompletionClassInverse}{1,2,5}}}
  \proofstepwidestar[1,2,5]{%
    \begin{aligned}[t]
      &\exists w\in\GroupCompletionCarrier{M,\star,e}\,\Bigl(
        z\GroupCompletionOp{M,\star,e}w
        =\GroupCompletionZero{M,\star,e}\land{}\\[-2pt]
      &\qquad w\GroupCompletionOp{M,\star,e}z
        =\GroupCompletionZero{M,\star,e}\Bigr)
    \end{aligned}}{%
    \rEI{\rAI{6,7}}}
  \proofstepwidestar[1,2,3]{%
    \begin{aligned}[t]
      &\exists w\in\GroupCompletionCarrier{M,\star,e}\,\Bigl(
        z\GroupCompletionOp{M,\star,e}w
        =\GroupCompletionZero{M,\star,e}\land{}\\[-2pt]
      &\qquad w\GroupCompletionOp{M,\star,e}z
        =\GroupCompletionZero{M,\star,e}\Bigr)
    \end{aligned}}{%
    \rEE{4,5,8}}
\end{tabproofwide}

\newpage
\FormulaThmDeltaK[Die formalen Differenzen bilden eine abelsche Gruppe]{%
  \begin{aligned}[t]
    &\CommutativeMonoidAlg{M}{\star}{e}\dsep
      \mathsf{K\ddot{u}rz}(M,\star)\vdash{}\\[-2pt]
    &\AbelianGroupAlg{\GroupCompletionCarrier{M,\star,e}}
      {\GroupCompletionOp{M,\star,e}}{\GroupCompletionZero{M,\star,e}}
  \end{aligned}
}{GroupCompletionAbelianGroup}{%
  \DeltaRow{Mengen}{M\dsep e\dsep z\dsep w}
  \DeltaRow{Funktionen}{\star\dsep\GroupCompletionOp{M,\star,e}}
}
\begin{tabproofwide}
  \proofstepwidestar[1]{\CommutativeMonoidAlg{M}{\star}{e}}{\rA}
  \proofstepwidestar[2]{\mathsf{K\ddot{u}rz}(M,\star)}{\rA}
  \proofstepwidestar[1,2]{%
    \CommutativeMonoidAlg{\GroupCompletionCarrier{M,\star,e}}
      {\GroupCompletionOp{M,\star,e}}{\GroupCompletionZero{M,\star,e}}}{%
    \FormulaRefAuto{GroupCompletionCommutativeMonoid}{1,2}}
  \proofstepwidestar[1,2]{%
    \MonoidAlg{\GroupCompletionCarrier{M,\star,e}}
      {\GroupCompletionOp{M,\star,e}}{\GroupCompletionZero{M,\star,e}}}{%
    \FormulaRefAuto{CommutativeMonoidBaseMonoidAxiom}[ax]{3}}
  \proofstepwidestar[5]{z\in\GroupCompletionCarrier{M,\star,e}}{\rA}
  \proofstepwidestar[1,2,5]{%
    \begin{aligned}[t]
      &\exists w\in\GroupCompletionCarrier{M,\star,e}\,\Bigl(
        z\GroupCompletionOp{M,\star,e}w
        =\GroupCompletionZero{M,\star,e}\land{}\\[-2pt]
      &\qquad w\GroupCompletionOp{M,\star,e}z
        =\GroupCompletionZero{M,\star,e}\Bigr)
    \end{aligned}}{%
    \FormulaRefAuto{GroupCompletionInverseExistence}{1,2,5}}
  \proofstepwidestar[1,2]{%
    \GroupAlg{\GroupCompletionCarrier{M,\star,e}}
      {\GroupCompletionOp{M,\star,e}}{\GroupCompletionZero{M,\star,e}}}{%
    \FormulaRefAuto{GroupDef}[def]{4,6}}
  \proofstepwidestar[8]{w\in\GroupCompletionCarrier{M,\star,e}}{\rA}
  \proofstepwidestar[1,2,5,8]{%
    \begin{aligned}[t]
      &z\GroupCompletionOp{M,\star,e}w={}\\[-2pt]
      &\quad w\GroupCompletionOp{M,\star,e}z
    \end{aligned}}{%
    \FormulaRefAuto{GroupCompletionOperationCommutative}{1,2,5,8}}
  \proofstepwidestar[1,2]{%
    \AbelianGroupAlg{\GroupCompletionCarrier{M,\star,e}}
      {\GroupCompletionOp{M,\star,e}}{\GroupCompletionZero{M,\star,e}}}{%
    \FormulaRefAuto{AbelianGroupDef}[def]{7,9}}
\end{tabproofwide}

\FormulaThmDeltaK[Einbettungssatz für kommutative kürzbare Monoide]{%
  \begin{aligned}[t]
    &\CommutativeMonoidAlg{M}{\star}{e}\dsep
      \mathsf{K\ddot{u}rz}(M,\star)\vdash{}\\[-2pt]
    &\AbelianGroupAlg{\GroupCompletionCarrier{M,\star,e}}
      {\GroupCompletionOp{M,\star,e}}{\GroupCompletionZero{M,\star,e}}
      \land{}\\[-2pt]
    &\CommutativeMonoidHom{\GroupCompletionEmbed{M,\star,e}}{M,\star,e}
      {\GroupCompletionCarrier{M,\star,e},
       \GroupCompletionOp{M,\star,e},\GroupCompletionZero{M,\star,e}}
      \land{}\\[-2pt]
    &\GroupCompletionEmbed{M,\star,e}\colon
      M\inj\GroupCompletionCarrier{M,\star,e}
  \end{aligned}
}{GroupCompletionEmbeddingTheorem}{%
  \DeltaRow{Mengen}{M\dsep e}
  \DeltaRow{Funktionen}{\star\dsep\GroupCompletionOp{M,\star,e}
    \dsep\GroupCompletionEmbed{M,\star,e}}
}
\begin{tabproofwide}
  \proofstepwidestar[1]{\CommutativeMonoidAlg{M}{\star}{e}}{\rA}
  \proofstepwidestar[2]{\mathsf{K\ddot{u}rz}(M,\star)}{\rA}
  \proofstepwidestar[1,2]{%
    \AbelianGroupAlg{\GroupCompletionCarrier{M,\star,e}}
      {\GroupCompletionOp{M,\star,e}}{\GroupCompletionZero{M,\star,e}}}{%
    \FormulaRefAuto{GroupCompletionAbelianGroup}{1,2}}
  \proofstepwidestar[1,2]{%
    \CommutativeMonoidAlg{\GroupCompletionCarrier{M,\star,e}}
      {\GroupCompletionOp{M,\star,e}}{\GroupCompletionZero{M,\star,e}}}{%
    \FormulaRefAuto{AbelianGroupIsCommutativeMonoid}{3}}
  \proofstepwidestar[1,2]{%
    \MonoidHom{\GroupCompletionEmbed{M,\star,e}}{M,\star,e}
      {\GroupCompletionCarrier{M,\star,e},
       \GroupCompletionOp{M,\star,e},\GroupCompletionZero{M,\star,e}}}{%
    \FormulaRefAuto{GroupCompletionEmbeddingMonoidHom}{1,2}}
  \proofstepwidestar[1,2]{%
    \CommutativeMonoidHom{\GroupCompletionEmbed{M,\star,e}}{M,\star,e}
      {\GroupCompletionCarrier{M,\star,e},
       \GroupCompletionOp{M,\star,e},\GroupCompletionZero{M,\star,e}}}{%
    \FormulaRefAuto{CommutativeMonoidHomDef}[def]{1,4,5}}
  \proofstepwidestar[1,2]{\GroupCompletionEmbed{M,\star,e}\colon
    M\inj\GroupCompletionCarrier{M,\star,e}}{%
    \FormulaRefAuto{GroupCompletionEmbeddingInjective}{1,2}}
  \proofstepwidestar[1,2]{\begin{aligned}[t]
    &\AbelianGroupAlg{\GroupCompletionCarrier{M,\star,e}}
      {\GroupCompletionOp{M,\star,e}}{\GroupCompletionZero{M,\star,e}}
      \land{}\\[-2pt]
    &\CommutativeMonoidHom{\GroupCompletionEmbed{M,\star,e}}{M,\star,e}
      {\GroupCompletionCarrier{M,\star,e},
       \GroupCompletionOp{M,\star,e},\GroupCompletionZero{M,\star,e}}
      \land{}\\[-2pt]
    &\GroupCompletionEmbed{M,\star,e}\colon
      M\inj\GroupCompletionCarrier{M,\star,e}
  \end{aligned}}{\rAI{3,\rAI{6,7}}}
\end{tabproofwide}

\begin{remark}
Die Kürzbarkeit wird in dieser Konstruktion beim Nachweis der Transitivität
der Relation formaler Differenzen verwendet. Die Gleichheitsreflexion der
kanonischen Einbettung folgt anschließend bereits aus den Neutralitätsgesetzen.
Für \((M,\star,e)=(\mathbb N,+,0)\) ergeben die Definitionen dieselbe
Kreuzsummenrelation und dieselbe Klassenoperation wie die Konstruktion der
ganzen Zahlen in Band~17.
\end{remark}

\chapter{Ringe mit Eins}

\section{Ringaxiome}

\FormulaAxiomDeltaK[Linksdistributivität]{%
  a,b,c\in R\vdash
  a\cdot(b+c)=a\cdot b+a\cdot c%
}{RingLeftDistributivityLaw}{%
  \DeltaRow{Mengen}{R\dsep a\dsep b\dsep c}
  \DeltaRow{Funktionen}{+\dsep\cdot}
}

\FormulaAxiomDeltaK[Rechtsdistributivität]{%
  a,b,c\in R\vdash
  (a+b)\cdot c=a\cdot c+b\cdot c%
}{RingRightDistributivityLaw}{%
  \DeltaRow{Mengen}{R\dsep a\dsep b\dsep c}
  \DeltaRow{Funktionen}{+\dsep\cdot}
}

\FormulaDefDeltaK[Begriff des Rings mit Eins]{%
  \RingAlg{R}{+}{\cdot}{0}{1}%
}{RingDef}{%
  \DeltaRow{Mengen}{R\dsep 0\dsep 1\dsep a\dsep b\dsep c}
  \DeltaRow{Funktionen}{+\dsep\cdot}
  \DeltaRow{\textbf{Axiome}}{}
  \DeltaRow{Additive abelsche Gruppe}
    {\AbelianGroupAlg{R}{+}{0}}
    [\FormulaRefAuto{AbelianGroupDef}[def]]
  \DeltaRow{Multiplikatives Monoid}
    {\MonoidAlg{R}{\cdot}{1}}
    [\FormulaRefAuto{MonoidDef}[def]]
  \DeltaRow{Linksdistributivität}{%
    \begin{aligned}[t]
      &a,b,c\in R\vdash{}\\[-2pt]
      &a\cdot(b+c)=a\cdot b+a\cdot c
    \end{aligned}}
    [\FormulaRefAuto{RingLeftDistributivityLaw}[ax]]
  \DeltaRow{Rechtsdistributivität}{%
    \begin{aligned}[t]
      &a,b,c\in R\vdash{}\\[-2pt]
      &(a+b)\cdot c=a\cdot c+b\cdot c
    \end{aligned}}
    [\FormulaRefAuto{RingRightDistributivityLaw}[ax]]
  \DeltaRow{\textbf{Neues Symbol}}{}
  \DeltaPrem{Ring mit Eins}{\RingAlg{R}{+}{\cdot}{0}{1}}
}

\FormulaAxiomDeltaK[Additiver Anteil eines Rings]{%
  \RingAlg{R}{+}{\cdot}{0}{1}
  \vdash\AbelianGroupAlg{R}{+}{0}%
}{RingAdditiveAbelianGroupAxiom}{%
  \DeltaRow{Mengen}{R\dsep 0\dsep 1}
  \DeltaRow{Funktionen}{+\dsep\cdot}
}

\FormulaAxiomDeltaK[Multiplikativer Anteil eines Rings]{%
  \RingAlg{R}{+}{\cdot}{0}{1}
  \vdash\MonoidAlg{R}{\cdot}{1}%
}{RingMultiplicativeMonoidAxiom}{%
  \DeltaRow{Mengen}{R\dsep 0\dsep 1}
  \DeltaRow{Funktionen}{+\dsep\cdot}
}

\FormulaAxiomDeltaK[Linkes Distributivgesetz eines Rings]{%
  \RingAlg{R}{+}{\cdot}{0}{1}\dsep a,b,c\in R
  \vdash a\cdot(b+c)=a\cdot b+a\cdot c%
}{RingLeftDistributivityAxiom}{%
  \DeltaRow{Mengen}{R\dsep 0\dsep 1\dsep a\dsep b\dsep c}
  \DeltaRow{Funktionen}{+\dsep\cdot}
}

\FormulaAxiomDeltaK[Rechtes Distributivgesetz eines Rings]{%
  \RingAlg{R}{+}{\cdot}{0}{1}\dsep a,b,c\in R
  \vdash (a+b)\cdot c=a\cdot c+b\cdot c%
}{RingRightDistributivityAxiom}{%
  \DeltaRow{Mengen}{R\dsep 0\dsep 1\dsep a\dsep b\dsep c}
  \DeltaRow{Funktionen}{+\dsep\cdot}
}

\section{Kommutative Ringe}

\FormulaAxiomDeltaK[Kommutativität der Ringmultiplikation]{%
  \RingAlg{R}{+}{\cdot}{0}{1}\dsep a,b\in R
  \vdash a\cdot b=b\cdot a%
}{CommutativeRingMultiplicationLaw}{%
  \DeltaRow{Mengen}{R\dsep 0\dsep 1\dsep a\dsep b}
  \DeltaRow{Funktionen}{+\dsep\cdot}
}

\FormulaDefDeltaK[Begriff des kommutativen Rings mit Eins]{%
  \CommutativeRingAlg{R}{+}{\cdot}{0}{1}%
}{CommutativeRingDef}{%
  \DeltaRow{Mengen}{R\dsep 0\dsep 1\dsep a\dsep b}
  \DeltaRow{Funktionen}{+\dsep\cdot}
  \DeltaRow{\textbf{Axiome}}{}
  \DeltaRow{Ring mit Eins}{\RingAlg{R}{+}{\cdot}{0}{1}}
    [\FormulaRefAuto{RingDef}[def]]
  \DeltaRow{Kommutative Multiplikation}{%
    a,b\in R\vdash a\cdot b=b\cdot a}
    [\FormulaRefAuto{CommutativeRingMultiplicationLaw}[ax]]
  \DeltaRow{\textbf{Neues Symbol}}{}
  \DeltaPrem{Kommutativer Ring mit Eins}
    {\CommutativeRingAlg{R}{+}{\cdot}{0}{1}}
}

\FormulaAxiomDeltaK[Ringanteil eines kommutativen Rings]{%
  \CommutativeRingAlg{R}{+}{\cdot}{0}{1}
  \vdash\RingAlg{R}{+}{\cdot}{0}{1}%
}{CommutativeRingBaseRingAxiom}{%
  \DeltaRow{Mengen}{R\dsep 0\dsep 1}
  \DeltaRow{Funktionen}{+\dsep\cdot}
}

\FormulaAxiomDeltaK[Multiplikationskommutativität im kommutativen Ring]{%
  \CommutativeRingAlg{R}{+}{\cdot}{0}{1}\dsep a,b\in R
  \vdash a\cdot b=b\cdot a%
}{CommutativeRingMultiplicationAxiom}{%
  \DeltaRow{Mengen}{R\dsep 0\dsep 1\dsep a\dsep b}
  \DeltaRow{Funktionen}{+\dsep\cdot}
}

\section{Abgeleitete Trägergesetze}

\FormulaThmDeltaK[Additive Gruppe eines Rings]{%
  \RingAlg{R}{+}{\cdot}{0}{1}\vdash\GroupAlg{R}{+}{0}%
}{RingAdditiveGroup}{%
  \DeltaRow{Mengen}{R\dsep 0\dsep 1}
  \DeltaRow{Funktionen}{+\dsep\cdot}
}
\begin{tabproof}
  \proofstep{1}{\RingAlg{R}{+}{\cdot}{0}{1}}{\rA}
  \proofstep{1}{\AbelianGroupAlg{R}{+}{0}}{%
    \FormulaRefAuto{RingAdditiveAbelianGroupAxiom}[ax]{1}}
  \proofstep{1}{\GroupAlg{R}{+}{0}}{%
    \FormulaRefAuto{AbelianGroupBaseGroupAxiom}[ax]{2}}
\end{tabproof}

\FormulaThmDeltaK[Multiplikative Halbgruppe eines Rings]{%
  \RingAlg{R}{+}{\cdot}{0}{1}\vdash\SemigroupAlg{R}{\cdot}%
}{RingMultiplicativeSemigroup}{%
  \DeltaRow{Mengen}{R\dsep 0\dsep 1}
  \DeltaRow{Funktionen}{+\dsep\cdot}
}
\begin{tabproof}
  \proofstep{1}{\RingAlg{R}{+}{\cdot}{0}{1}}{\rA}
  \proofstep{1}{\MonoidAlg{R}{\cdot}{1}}{%
    \FormulaRefAuto{RingMultiplicativeMonoidAxiom}[ax]{1}}
  \proofstep{1}{\SemigroupAlg{R}{\cdot}}{%
    \FormulaRefAuto{MonoidBaseSemigroupAxiom}[ax]{2}}
\end{tabproof}

\FormulaThmDeltaK[Abgeschlossenheit der Ringmultiplikation]{%
  \RingAlg{R}{+}{\cdot}{0}{1}\dsep a,b\in R
  \vdash a\cdot b\in R%
}{RingMultiplicationClosure}{%
  \DeltaRow{Mengen}{R\dsep 0\dsep 1\dsep a\dsep b}
  \DeltaRow{Funktionen}{+\dsep\cdot}
}
\begin{tabproof}
  \proofstep{1}{\RingAlg{R}{+}{\cdot}{0}{1}}{\rA}
  \proofstep{2}{a\in R}{\rA}
  \proofstep{3}{b\in R}{\rA}
  \proofstep{1}{\SemigroupAlg{R}{\cdot}}{%
    \FormulaRefAuto{RingMultiplicativeSemigroup}[thm]{1}}
  \proofstep{1,2,3}{a\cdot b\in R}{%
    \FormulaRefAuto{SemigroupClosureAxiom}[ax]{4,2,3}}
\end{tabproof}

\section{Additives Inverses und Vorzeichen}

\FormulaDefDeltaK[Additives Inverses in einem Ring]{%
  \begin{aligned}[t]
    &\RingAlg{R}{+}{\cdot}{0}{1}\dsep a\in R\vdash{}\\
    &-a\coloneqq
      \iota b\,
      \bigl(b\in R\land a+b=0\land b+a=0\bigr)
  \end{aligned}
}{RingNegativeElement}{%
  \DeltaRow{Mengen}{R\dsep 0\dsep 1\dsep a\dsep b}
  \DeltaRow{Funktionen}{+\dsep\cdot}
  \DeltaRow{Träger}{-a\in R}
  \DeltaRow{Rechtsinverse}{a+(-a)=0}
  \DeltaRow{Linksinverse}{(-a)+a=0}
  \DeltaRow{Eindeutigkeit}{\FormulaRefAuto{GroupInverseUnique}[thm]}
  \DeltaRow{\textbf{Neues Symbol}}{}
  \DeltaPrem{Additives Inverses}{-a}
}

\FormulaThmDeltaK[Additives Inverses liegt im Ringträger]{%
  \RingAlg{R}{+}{\cdot}{0}{1}\dsep a\in R\vdash -a\in R%
}{RingNegativeCarrier}{%
  \DeltaRow{Mengen}{R\dsep 0\dsep 1\dsep a}
  \DeltaRow{Funktionen}{+\dsep\cdot}
}
\begin{tabproof}
  \proofstep{1}{\RingAlg{R}{+}{\cdot}{0}{1}}{\rA}
  \proofstep{2}{a\in R}{\rA}
  \proofstep{1,2}{-a\in R}{%
    \FormulaRefAuto{RingNegativeElement}[def]{1,2}}
\end{tabproof}

\FormulaThmDeltaK[Rechte additive Inversengleichung]{%
  \RingAlg{R}{+}{\cdot}{0}{1}\dsep a\in R
  \vdash a+(-a)=0%
}{RingNegativeRight}{%
  \DeltaRow{Mengen}{R\dsep 0\dsep 1\dsep a}
  \DeltaRow{Funktionen}{+\dsep\cdot}
}
\begin{tabproof}
  \proofstep{1}{\RingAlg{R}{+}{\cdot}{0}{1}}{\rA}
  \proofstep{2}{a\in R}{\rA}
  \proofstep{1,2}{a+(-a)=0}{%
    \FormulaRefAuto{RingNegativeElement}[def]{1,2}}
\end{tabproof}

\FormulaThmDeltaK[Linke additive Inversengleichung]{%
  \RingAlg{R}{+}{\cdot}{0}{1}\dsep a\in R
  \vdash (-a)+a=0%
}{RingNegativeLeft}{%
  \DeltaRow{Mengen}{R\dsep 0\dsep 1\dsep a}
  \DeltaRow{Funktionen}{+\dsep\cdot}
}
\begin{tabproof}
  \proofstep{1}{\RingAlg{R}{+}{\cdot}{0}{1}}{\rA}
  \proofstep{2}{a\in R}{\rA}
  \proofstep{1,2}{(-a)+a=0}{%
    \FormulaRefAuto{RingNegativeElement}[def]{1,2}}
\end{tabproof}

\FormulaThmDeltaK[Doppelte Negation]{%
  \RingAlg{R}{+}{\cdot}{0}{1}\dsep a\in R
  \vdash -(-a)=a%
}{RingDoubleNegative}{%
  \DeltaRow{Mengen}{R\dsep 0\dsep 1\dsep a}
  \DeltaRow{Funktionen}{+\dsep\cdot}
}
\begin{tabproofwide}
  \proofstepwidestar[1]{\RingAlg{R}{+}{\cdot}{0}{1}}{\rA}
  \proofstepwidestar[2]{a\in R}{\rA}
  \proofstepwidestar[1]{\GroupAlg{R}{+}{0}}{%
    \FormulaRefAuto{RingAdditiveGroup}[thm]{1}}
  \proofstepwidestar[1,2]{-a\in R}{%
    \FormulaRefAuto{RingNegativeCarrier}[thm]{1,2}}
  \proofstepwidestar[1,4]{-(-a)\in R}{%
    \FormulaRefAuto{RingNegativeCarrier}[thm]{1,4}}
  \proofstepwidestar[1,4]{-(-a)+(-a)=0}{%
    \FormulaRefAuto{RingNegativeLeft}[thm]{1,4}}
  \proofstepwidestar[1,2]{(-a)+a=0}{%
    \FormulaRefAuto{RingNegativeLeft}[thm]{1,2}}
  \proofstepwidestar[1,2,3,4,5,6,7]{-(-a)=a}{%
    \FormulaRefAuto{GroupInverseUnique}[thm]{3,4,5,2,6,7}}
\end{tabproofwide}

\section{Nullabsorption}

\FormulaThmDeltaK[Linke Nullabsorption]{%
  \RingAlg{R}{+}{\cdot}{0}{1}\dsep a\in R
  \vdash 0\cdot a=0%
}{RingLeftZeroAbsorption}{%
  \DeltaRow{Mengen}{R\dsep 0\dsep 1\dsep a}
  \DeltaRow{Funktionen}{+\dsep\cdot}
}
\begin{tabproofsplitwide}
  \proofpartwideindR[Verdopplung]{%
    \RingAlg{R}{+}{\cdot}{0}{1}\dsep a\in R
    \vdash 0\cdot a=0\cdot a+0\cdot a
  }{%
    \RingAlg{R}{+}{\cdot}{0}{1}\dsep a\in R
    \vdash 0\cdot a=0\cdot a+0\cdot a%
  }[RingLeftZeroDoubling]
    \proofstepwidestar[1]{\RingAlg{R}{+}{\cdot}{0}{1}}{\rA}
    \proofstepwidestar[2]{a\in R}{\rA}
    \proofstepwidestar[1]{\GroupAlg{R}{+}{0}}{%
      \FormulaRefAuto{RingAdditiveGroup}[thm]{1}}
    \proofstepwidestar[1]{0\in R}{%
      \FormulaRefAuto{GroupIdentityCarrier}[thm]{3}}
    \proofstepwide[3,4]{0+0}{=}{0}{%
      \FormulaRefAuto{GroupRightIdentity}[thm]{3,4}}
    \proofstepwide[5]{0\cdot a}{=}{(0+0)\cdot a}{%
      \rIE{\FormulaRefAuto{a=b\vdash b=a}{5}}}
    \proofstepwide[1,2,4]{(0+0)\cdot a}{=}{%
      0\cdot a+0\cdot a}{%
      \FormulaRefAuto{RingRightDistributivityAxiom}[ax]{1,4,4,2}}
    \proofstepwidestar[1,2]{0\cdot a=0\cdot a+0\cdot a}{%
      \rChain{6,7}}
  \closeproofpartwideindR

  \proofpartwide{Kürzung}
    \proofstepwidestar[1]{\RingAlg{R}{+}{\cdot}{0}{1}}{\rA}
    \proofstepwidestar[2]{a\in R}{\rA}
    \proofstepwidestar[1]{\GroupAlg{R}{+}{0}}{%
      \FormulaRefAuto{RingAdditiveGroup}[thm]{1}}
    \proofstepwidestar[1,2]{0\cdot a\in R}{%
      \FormulaRefAuto{RingMultiplicationClosure}[thm]{%
        1,\FormulaRefAuto{GroupIdentityCarrier}[thm]{3},2}}
    \proofstepwide[3,4]{0\cdot a+0}{=}{0\cdot a}{%
      \FormulaRefAuto{GroupRightIdentity}[thm]{3,4}}
    \proofstepwide[1,2]{0\cdot a}{=}{0\cdot a+0\cdot a}{%
      \FormulaRefAuto{RingLeftZeroDoubling}[thm]{1,2}}
    \proofstepwidestar[1,2,3,4,5,6]{%
      0\cdot a+0=0\cdot a+0\cdot a}{\rChain{5,6}}
    \proofstepwidestar[1,2,3,4,7]{0=0\cdot a}{%
      \FormulaRefAuto{GroupLeftCancellation}[thm]{3,4,
        \FormulaRefAuto{GroupIdentityCarrier}[thm]{3},4,7}}
    \proofstepwidestar[1,2]{0\cdot a=0}{%
      \FormulaRefAuto{a=b\vdash b=a}{8}}
  \closeproofpartwide
\end{tabproofsplitwide}

\FormulaThmDeltaK[Rechte Nullabsorption]{%
  \RingAlg{R}{+}{\cdot}{0}{1}\dsep a\in R
  \vdash a\cdot0=0%
}{RingRightZeroAbsorption}{%
  \DeltaRow{Mengen}{R\dsep 0\dsep 1\dsep a}
  \DeltaRow{Funktionen}{+\dsep\cdot}
}
\begin{tabproofsplitwide}
  \proofpartwideindR[Verdopplung]{%
    \RingAlg{R}{+}{\cdot}{0}{1}\dsep a\in R
    \vdash a\cdot0=a\cdot0+a\cdot0
  }{%
    \RingAlg{R}{+}{\cdot}{0}{1}\dsep a\in R
    \vdash a\cdot0=a\cdot0+a\cdot0%
  }[RingRightZeroDoubling]
    \proofstepwidestar[1]{\RingAlg{R}{+}{\cdot}{0}{1}}{\rA}
    \proofstepwidestar[2]{a\in R}{\rA}
    \proofstepwidestar[1]{\GroupAlg{R}{+}{0}}{%
      \FormulaRefAuto{RingAdditiveGroup}[thm]{1}}
    \proofstepwidestar[1]{0\in R}{%
      \FormulaRefAuto{GroupIdentityCarrier}[thm]{3}}
    \proofstepwide[3,4]{0+0}{=}{0}{%
      \FormulaRefAuto{GroupRightIdentity}[thm]{3,4}}
    \proofstepwide[5]{a\cdot0}{=}{a\cdot(0+0)}{%
      \rIE{\FormulaRefAuto{a=b\vdash b=a}{5}}}
    \proofstepwide[1,2,4]{a\cdot(0+0)}{=}{%
      a\cdot0+a\cdot0}{%
      \FormulaRefAuto{RingLeftDistributivityAxiom}[ax]{1,2,4,4}}
    \proofstepwidestar[1,2]{a\cdot0=a\cdot0+a\cdot0}{%
      \rChain{6,7}}
  \closeproofpartwideindR

  \proofpartwide{Kürzung}
    \proofstepwidestar[1]{\RingAlg{R}{+}{\cdot}{0}{1}}{\rA}
    \proofstepwidestar[2]{a\in R}{\rA}
    \proofstepwidestar[1]{\GroupAlg{R}{+}{0}}{%
      \FormulaRefAuto{RingAdditiveGroup}[thm]{1}}
    \proofstepwidestar[1,2]{a\cdot0\in R}{%
      \FormulaRefAuto{RingMultiplicationClosure}[thm]{%
        1,2,\FormulaRefAuto{GroupIdentityCarrier}[thm]{3}}}
    \proofstepwide[3,4]{a\cdot0+0}{=}{a\cdot0}{%
      \FormulaRefAuto{GroupRightIdentity}[thm]{3,4}}
    \proofstepwide[1,2]{a\cdot0}{=}{a\cdot0+a\cdot0}{%
      \FormulaRefAuto{RingRightZeroDoubling}[thm]{1,2}}
    \proofstepwidestar[1,2,3,4,5,6]{%
      a\cdot0+0=a\cdot0+a\cdot0}{\rChain{5,6}}
    \proofstepwidestar[1,2,3,4,7]{0=a\cdot0}{%
      \FormulaRefAuto{GroupLeftCancellation}[thm]{3,4,
        \FormulaRefAuto{GroupIdentityCarrier}[thm]{3},4,7}}
    \proofstepwidestar[1,2]{a\cdot0=0}{%
      \FormulaRefAuto{a=b\vdash b=a}{8}}
  \closeproofpartwide
\end{tabproofsplitwide}

\section{Vorzeichenregeln}

\FormulaThmDeltaK[Negatives Vorzeichen im linken Faktor]{%
  \RingAlg{R}{+}{\cdot}{0}{1}\dsep a,b\in R
  \vdash (-a)\cdot b=-(a\cdot b)%
}{RingLeftNegativeProduct}{%
  \DeltaRow{Mengen}{R\dsep 0\dsep 1\dsep a\dsep b}
  \DeltaRow{Funktionen}{+\dsep\cdot}
}
\begin{tabproofsplitwide}
  \proofpartwideindR[Rechtsinverse Eigenschaft]{%
    \begin{aligned}[t]
      &\RingAlg{R}{+}{\cdot}{0}{1}\dsep a,b\in R\vdash{}\\
      &a\cdot b+(-a)\cdot b=0
    \end{aligned}
  }{%
    \RingAlg{R}{+}{\cdot}{0}{1}\dsep a,b\in R
    \vdash a\cdot b+(-a)\cdot b=0%
  }[RingLeftNegativeRightInverse]
    \proofstepwidestar[1]{\RingAlg{R}{+}{\cdot}{0}{1}}{\rA}
    \proofstepwidestar[2]{a\in R}{\rA}
    \proofstepwidestar[3]{b\in R}{\rA}
    \proofstepwidestar[1,2]{-a\in R}{%
      \FormulaRefAuto{RingNegativeCarrier}[thm]{1,2}}
    \proofstepwide[1,2,3,4]{a\cdot b+(-a)\cdot b}{=}{%
      (a+(-a))\cdot b}{%
      \FormulaRefAuto{a=b\vdash b=a}{%
        \FormulaRefAuto{RingRightDistributivityAxiom}[ax]{1,2,4,3}}}
    \proofstepwide[1,2,3]{(a+(-a))\cdot b}{=}{0\cdot b}{%
      \FormulaRefAuto{RingNegativeRight}[thm]{1,2}}
    \proofstepwide[1,3]{0\cdot b}{=}{0}{%
      \FormulaRefAuto{RingLeftZeroAbsorption}[thm]{1,3}}
    \proofstepwidestar[1,2,3]{a\cdot b+(-a)\cdot b=0}{%
      \rChain{5,6,7}}
  \closeproofpartwideindR

  \proofpartwideindR[Linksinverse Eigenschaft]{%
    \begin{aligned}[t]
      &\RingAlg{R}{+}{\cdot}{0}{1}\dsep a,b\in R\vdash{}\\
      &(-a)\cdot b+a\cdot b=0
    \end{aligned}
  }{%
    \RingAlg{R}{+}{\cdot}{0}{1}\dsep a,b\in R
    \vdash (-a)\cdot b+a\cdot b=0%
  }[RingLeftNegativeLeftInverse]
    \proofstepwidestar[1]{\RingAlg{R}{+}{\cdot}{0}{1}}{\rA}
    \proofstepwidestar[2]{a\in R}{\rA}
    \proofstepwidestar[3]{b\in R}{\rA}
    \proofstepwidestar[1,2]{-a\in R}{%
      \FormulaRefAuto{RingNegativeCarrier}[thm]{1,2}}
    \proofstepwide[1,2,3,4]{(-a)\cdot b+a\cdot b}{=}{%
      ((-a)+a)\cdot b}{%
      \FormulaRefAuto{a=b\vdash b=a}{%
        \FormulaRefAuto{RingRightDistributivityAxiom}[ax]{1,4,2,3}}}
    \proofstepwide[1,2,3]{((-a)+a)\cdot b}{=}{0\cdot b}{%
      \FormulaRefAuto{RingNegativeLeft}[thm]{1,2}}
    \proofstepwide[1,3]{0\cdot b}{=}{0}{%
      \FormulaRefAuto{RingLeftZeroAbsorption}[thm]{1,3}}
    \proofstepwidestar[1,2,3]{(-a)\cdot b+a\cdot b=0}{%
      \rChain{5,6,7}}
  \closeproofpartwideindR

  \proofpartwide{Eindeutigkeit}
    \proofstepwidestar[1]{\RingAlg{R}{+}{\cdot}{0}{1}}{\rA}
    \proofstepwidestar[2]{a\in R}{\rA}
    \proofstepwidestar[3]{b\in R}{\rA}
    \proofstepwidestar[1]{\GroupAlg{R}{+}{0}}{%
      \FormulaRefAuto{RingAdditiveGroup}[thm]{1}}
    \proofstepwidestar[1,2,3]{a\cdot b\in R}{%
      \FormulaRefAuto{RingMultiplicationClosure}[thm]{1,2,3}}
    \proofstepwidestar[1,2,3]{(-a)\cdot b\in R}{%
      \FormulaRefAuto{RingMultiplicationClosure}[thm]{%
        1,\FormulaRefAuto{RingNegativeCarrier}[thm]{1,2},3}}
    \proofstepwidestar[1,5]{-(a\cdot b)\in R}{%
      \FormulaRefAuto{RingNegativeCarrier}[thm]{1,5}}
    \proofstepwidestar[1,2,3]{(-a)\cdot b+a\cdot b=0}{%
      \FormulaRefAuto{RingLeftNegativeLeftInverse}[thm]{1,2,3}}
    \proofstepwidestar[1,5]{a\cdot b+(-(a\cdot b))=0}{%
      \FormulaRefAuto{RingNegativeRight}[thm]{1,5}}
    \proofstepwidestar[1,2,3,4,5,6,7,8,9]{%
      (-a)\cdot b=-(a\cdot b)}{%
      \FormulaRefAuto{GroupInverseUnique}[thm]{4,5,6,7,8,9}}
  \closeproofpartwide
\end{tabproofsplitwide}

\FormulaThmDeltaK[Negatives Vorzeichen im rechten Faktor]{%
  \RingAlg{R}{+}{\cdot}{0}{1}\dsep a,b\in R
  \vdash a\cdot(-b)=-(a\cdot b)%
}{RingRightNegativeProduct}{%
  \DeltaRow{Mengen}{R\dsep 0\dsep 1\dsep a\dsep b}
  \DeltaRow{Funktionen}{+\dsep\cdot}
}
\begin{tabproofsplitwide}
  \proofpartwideindR[Rechtsinverse Eigenschaft]{%
    \begin{aligned}[t]
      &\RingAlg{R}{+}{\cdot}{0}{1}\dsep a,b\in R\vdash{}\\
      &a\cdot b+a\cdot(-b)=0
    \end{aligned}
  }{%
    \RingAlg{R}{+}{\cdot}{0}{1}\dsep a,b\in R
    \vdash a\cdot b+a\cdot(-b)=0%
  }[RingRightNegativeRightInverse]
    \proofstepwidestar[1]{\RingAlg{R}{+}{\cdot}{0}{1}}{\rA}
    \proofstepwidestar[2]{a\in R}{\rA}
    \proofstepwidestar[3]{b\in R}{\rA}
    \proofstepwidestar[1,3]{-b\in R}{%
      \FormulaRefAuto{RingNegativeCarrier}[thm]{1,3}}
    \proofstepwide[1,2,3,4]{a\cdot b+a\cdot(-b)}{=}{%
      a\cdot(b+(-b))}{%
      \FormulaRefAuto{a=b\vdash b=a}{%
        \FormulaRefAuto{RingLeftDistributivityAxiom}[ax]{1,2,3,4}}}
    \proofstepwide[1,2,3]{a\cdot(b+(-b))}{=}{a\cdot0}{%
      \FormulaRefAuto{RingNegativeRight}[thm]{1,3}}
    \proofstepwide[1,2]{a\cdot0}{=}{0}{%
      \FormulaRefAuto{RingRightZeroAbsorption}[thm]{1,2}}
    \proofstepwidestar[1,2,3]{a\cdot b+a\cdot(-b)=0}{%
      \rChain{5,6,7}}
  \closeproofpartwideindR

  \proofpartwideindR[Linksinverse Eigenschaft]{%
    \begin{aligned}[t]
      &\RingAlg{R}{+}{\cdot}{0}{1}\dsep a,b\in R\vdash{}\\
      &a\cdot(-b)+a\cdot b=0
    \end{aligned}
  }{%
    \RingAlg{R}{+}{\cdot}{0}{1}\dsep a,b\in R
    \vdash a\cdot(-b)+a\cdot b=0%
  }[RingRightNegativeLeftInverse]
    \proofstepwidestar[1]{\RingAlg{R}{+}{\cdot}{0}{1}}{\rA}
    \proofstepwidestar[2]{a\in R}{\rA}
    \proofstepwidestar[3]{b\in R}{\rA}
    \proofstepwidestar[1,3]{-b\in R}{%
      \FormulaRefAuto{RingNegativeCarrier}[thm]{1,3}}
    \proofstepwide[1,2,3,4]{a\cdot(-b)+a\cdot b}{=}{%
      a\cdot((-b)+b)}{%
      \FormulaRefAuto{a=b\vdash b=a}{%
        \FormulaRefAuto{RingLeftDistributivityAxiom}[ax]{1,2,4,3}}}
    \proofstepwide[1,2,3]{a\cdot((-b)+b)}{=}{a\cdot0}{%
      \FormulaRefAuto{RingNegativeLeft}[thm]{1,3}}
    \proofstepwide[1,2]{a\cdot0}{=}{0}{%
      \FormulaRefAuto{RingRightZeroAbsorption}[thm]{1,2}}
    \proofstepwidestar[1,2,3]{a\cdot(-b)+a\cdot b=0}{%
      \rChain{5,6,7}}
  \closeproofpartwideindR

  \proofpartwide{Eindeutigkeit}
    \proofstepwidestar[1]{\RingAlg{R}{+}{\cdot}{0}{1}}{\rA}
    \proofstepwidestar[2]{a\in R}{\rA}
    \proofstepwidestar[3]{b\in R}{\rA}
    \proofstepwidestar[1]{\GroupAlg{R}{+}{0}}{%
      \FormulaRefAuto{RingAdditiveGroup}[thm]{1}}
    \proofstepwidestar[1,2,3]{a\cdot b\in R}{%
      \FormulaRefAuto{RingMultiplicationClosure}[thm]{1,2,3}}
    \proofstepwidestar[1,2,3]{a\cdot(-b)\in R}{%
      \FormulaRefAuto{RingMultiplicationClosure}[thm]{%
        1,2,\FormulaRefAuto{RingNegativeCarrier}[thm]{1,3}}}
    \proofstepwidestar[1,5]{-(a\cdot b)\in R}{%
      \FormulaRefAuto{RingNegativeCarrier}[thm]{1,5}}
    \proofstepwidestar[1,2,3]{a\cdot(-b)+a\cdot b=0}{%
      \FormulaRefAuto{RingRightNegativeLeftInverse}[thm]{1,2,3}}
    \proofstepwidestar[1,5]{a\cdot b+(-(a\cdot b))=0}{%
      \FormulaRefAuto{RingNegativeRight}[thm]{1,5}}
    \proofstepwidestar[1,2,3,4,5,6,7,8,9]{%
      a\cdot(-b)=-(a\cdot b)}{%
      \FormulaRefAuto{GroupInverseUnique}[thm]{4,5,6,7,8,9}}
  \closeproofpartwide
\end{tabproofsplitwide}

\FormulaThmDeltaK[Produkt zweier negativer Faktoren]{%
  \RingAlg{R}{+}{\cdot}{0}{1}\dsep a,b\in R
  \vdash (-a)\cdot(-b)=a\cdot b%
}{RingNegativeTimesNegative}{%
  \DeltaRow{Mengen}{R\dsep 0\dsep 1\dsep a\dsep b}
  \DeltaRow{Funktionen}{+\dsep\cdot}
}
\begin{tabproofwide}
  \proofstepwidestar[1]{\RingAlg{R}{+}{\cdot}{0}{1}}{\rA}
  \proofstepwidestar[2]{a\in R}{\rA}
  \proofstepwidestar[3]{b\in R}{\rA}
  \proofstepwide[1,2,3]{(-a)\cdot(-b)}{=}{-(a\cdot(-b))}{%
    \FormulaRefAuto{RingLeftNegativeProduct}[thm]{%
      1,2,\FormulaRefAuto{RingNegativeCarrier}[thm]{1,3}}}
  \proofstepwide[1,2,3]{-(a\cdot(-b))}{=}{-(-(a\cdot b))}{%
    \FormulaRefAuto{RingRightNegativeProduct}[thm]{1,2,3}}
  \proofstepwide[1,2,3]{-(-(a\cdot b))}{=}{a\cdot b}{%
    \FormulaRefAuto{RingDoubleNegative}[thm]{%
      1,\FormulaRefAuto{RingMultiplicationClosure}[thm]{1,2,3}}}
  \proofstepwidestar[1,2,3]{(-a)\cdot(-b)=a\cdot b}{%
    \rChain{4,5,6}}
\end{tabproofwide}

\chapter{Homomorphismen von Ringen}

\section{Der einem Ring zugrunde liegende Halbring}

Die Nullabsorption wurde für Ringe aus den übrigen Ringaxiomen hergeleitet.
Damit ist jeder Ring mit Eins insbesondere ein Halbring mit Eins. Dieses
Resultat erlaubt es, die kanonische Abbildung aus \(\mathbb N\) ohne eine
zweite Rekursionskonstruktion zu verwenden.

\FormulaThmDeltaK[Jeder Ring mit Eins ist ein Halbring mit Eins]{%
  \RingAlg{R}{+}{\cdot}{0}{1}
  \vdash\SemiringAlg{R}{+}{\cdot}{0}{1}%
}{RingUnderlyingSemiring}{%
  \DeltaRow{Mengen}{R\dsep0\dsep1\dsep a\dsep b\dsep c}
  \DeltaRow{Funktionen}{+\dsep\cdot}
}
\begin{tabproofwide}
  \proofstepwidestar[1]{\RingAlg{R}{+}{\cdot}{0}{1}}{\rA}
  \proofstepwidestar[1]{\AbelianGroupAlg{R}{+}{0}}{%
    \FormulaRefAuto{RingAdditiveAbelianGroupAxiom}[ax]{1}}
  \proofstepwidestar[1]{\CommutativeMonoidAlg{R}{+}{0}}{%
    \FormulaRefAuto{AbelianGroupIsCommutativeMonoid}[thm]{2}}
  \proofstepwidestar[1]{\MonoidAlg{R}{\cdot}{1}}{%
    \FormulaRefAuto{RingMultiplicativeMonoidAxiom}[ax]{1}}
  \proofstepwidestar[5]{a\in R}{\rA}
  \proofstepwidestar[6]{b\in R}{\rA}
  \proofstepwidestar[7]{c\in R}{\rA}
  \proofstepwidestar[1,5,6,7]{%
    a\cdot(b+c)=a\cdot b+a\cdot c}{%
    \FormulaRefAuto{RingLeftDistributivityAxiom}[ax]{1,5,6,7}}
  \proofstepwidestar[1,5,6,7]{%
    (a+b)\cdot c=a\cdot c+b\cdot c}{%
    \FormulaRefAuto{RingRightDistributivityAxiom}[ax]{1,5,6,7}}
  \proofstepwidestar[1,5]{0\cdot a=0}{%
    \FormulaRefAuto{RingLeftZeroAbsorption}[thm]{1,5}}
  \proofstepwidestar[1,5]{a\cdot0=0}{%
    \FormulaRefAuto{RingRightZeroAbsorption}[thm]{1,5}}
  \proofstepwidestar[1]{\SemiringAlg{R}{+}{\cdot}{0}{1}}{%
    \FormulaRefAuto{SemiringDef}[def]{3,4,8,9,10,11}}
\end{tabproofwide}

\section{Ringhomomorphismen}

Da alle Ringe dieses Bandes ein ausgezeichnetes Einselement besitzen, meint
\emph{Ringhomomorphismus} im Folgenden stets einen einserhaltenden
Ringhomomorphismus.

\FormulaDefDeltaK[Begriff des Ringhomomorphismus]{%
  \RingHom{f}{R,\oplus,\odot,0_R,1_R}
    {S,\boxplus,\boxtimes,0_S,1_S}%
}{RingHomDef}{%
  \DeltaRow{Mengen}{R\dsep S\dsep0_R\dsep1_R\dsep0_S\dsep1_S}
  \DeltaRow{Funktionen}{f\dsep\oplus\dsep\odot\dsep\boxplus\dsep\boxtimes}
  \DeltaRow{\textbf{Axiome}}{}
  \DeltaPrem{Quellring}{%
    \RingAlg{R}{\oplus}{\odot}{0_R}{1_R}}
  \DeltaPrem{Zielring}{%
    \RingAlg{S}{\boxplus}{\boxtimes}{0_S}{1_S}}
  \DeltaPrem{Additiver Anteil}{%
    \GroupHom{f}{R,\oplus,0_R}{S,\boxplus,0_S}}
  \DeltaPrem{Multiplikativer Anteil}{%
    \MonoidHom{f}{R,\odot,1_R}{S,\boxtimes,1_S}}
  \DeltaRow{\textbf{Neues Symbol}}{}
  \DeltaPrem{Symbol}{\mathsf{RingHom}}
}

\FormulaAxiomDeltaK[Trägerabbildung eines Ringhomomorphismus]{%
  \RingHom{f}{R,\oplus,\odot,0_R,1_R}
    {S,\boxplus,\boxtimes,0_S,1_S}
  \vdash f\colon R\to S%
}{RingHomFunctionAxiom}{%
  \DeltaRow{Mengen}{R\dsep S}
  \DeltaRow{Funktionen}{f\dsep\oplus\dsep\odot\dsep\boxplus\dsep\boxtimes}
}
\par\smallskip

\FormulaAxiomDeltaK[Additionserhaltung eines Ringhomomorphismus]{%
  \begin{aligned}[t]
    &\RingHom{f}{R,\oplus,\odot,0_R,1_R}
      {S,\boxplus,\boxtimes,0_S,1_S}\dsep x,y\in R\\
    &\vdash f(x\oplus y)=f(x)\boxplus f(y)
  \end{aligned}
}{RingHomAdditionAxiom}{%
  \DeltaRow{Mengen}{R\dsep S\dsep x\dsep y}
  \DeltaRow{Funktionen}{f\dsep\oplus\dsep\odot\dsep\boxplus\dsep\boxtimes}
}
\par\smallskip

\FormulaAxiomDeltaK[Multiplikationserhaltung eines Ringhomomorphismus]{%
  \begin{aligned}[t]
    &\RingHom{f}{R,\oplus,\odot,0_R,1_R}
      {S,\boxplus,\boxtimes,0_S,1_S}\dsep x,y\in R\\
    &\vdash f(x\odot y)=f(x)\boxtimes f(y)
  \end{aligned}
}{RingHomMultiplicationAxiom}{%
  \DeltaRow{Mengen}{R\dsep S\dsep x\dsep y}
  \DeltaRow{Funktionen}{f\dsep\oplus\dsep\odot\dsep\boxplus\dsep\boxtimes}
}
\par\smallskip

\FormulaAxiomDeltaK[Einserhaltung eines Ringhomomorphismus]{%
  \RingHom{f}{R,\oplus,\odot,0_R,1_R}
    {S,\boxplus,\boxtimes,0_S,1_S}
  \vdash f(1_R)=1_S%
}{RingHomOneAxiom}{%
  \DeltaRow{Mengen}{R\dsep S\dsep1_R\dsep1_S}
  \DeltaRow{Funktionen}{f\dsep\oplus\dsep\odot\dsep\boxplus\dsep\boxtimes}
}

\FormulaThmDeltaK[Ringhomomorphismen erhalten die Null]{%
  \RingHom{f}{R,\oplus,\odot,0_R,1_R}
    {S,\boxplus,\boxtimes,0_S,1_S}
  \vdash f(0_R)=0_S%
}{RingHomZeroPreservation}{%
  \DeltaRow{Mengen}{R\dsep S\dsep0_R\dsep0_S}
  \DeltaRow{Funktionen}{f\dsep\oplus\dsep\odot\dsep\boxplus\dsep\boxtimes}
}
\begin{tabproof}
  \proofstep{1}{%
    \RingHom{f}{R,\oplus,\odot,0_R,1_R}
      {S,\boxplus,\boxtimes,0_S,1_S}}{\rA}
  \proofstep{1}{\GroupHom{f}{R,\oplus,0_R}{S,\boxplus,0_S}}{%
    \FormulaRefAuto{RingHomDef}[def]{1}}
  \proofstep{1}{f(0_R)=0_S}{%
    \FormulaRefAuto{GroupHomIdentityPreservation}[thm]{2}}
\end{tabproof}

\FormulaThmDeltaK[Ringhomomorphismen erhalten additive Inverse]{%
  \RingHom{f}{R,\oplus,\odot,0_R,1_R}
    {S,\boxplus,\boxtimes,0_S,1_S}\dsep x\in R
  \vdash f(-x)=-f(x)%
}{RingHomNegativePreservation}{%
  \DeltaRow{Mengen}{R\dsep S\dsep x}
  \DeltaRow{Funktionen}{f\dsep\oplus\dsep\odot\dsep\boxplus\dsep\boxtimes}
}
\begin{tabproof}
  \proofstep{1}{%
    \RingHom{f}{R,\oplus,\odot,0_R,1_R}
      {S,\boxplus,\boxtimes,0_S,1_S}}{\rA}
  \proofstep{2}{x\in R}{\rA}
  \proofstep{1}{\GroupHom{f}{R,\oplus,0_R}{S,\boxplus,0_S}}{%
    \FormulaRefAuto{RingHomDef}[def]{1}}
  \proofstep{1,2}{f(-x)=-f(x)}{%
    \FormulaRefAuto{GroupHomInversePreservation}[thm]{3,2}}
\end{tabproof}

\section{Ringisomorphismen und Komposition}

\FormulaDefDeltaK[Begriff des Ringisomorphismus]{%
  \RingIso{f}{R,\oplus,\odot,0_R,1_R}
    {S,\boxplus,\boxtimes,0_S,1_S}%
}{RingIsoDef}{%
  \DeltaRow{Mengen}{R\dsep S}
  \DeltaRow{Funktionen}{f\dsep\oplus\dsep\odot\dsep\boxplus\dsep\boxtimes}
  \DeltaRow{\textbf{Axiome}}{}
  \DeltaPrem{Ringhomomorphismus}{%
    \RingHom{f}{R,\oplus,\odot,0_R,1_R}
      {S,\boxplus,\boxtimes,0_S,1_S}}
  \DeltaPrem{Bijektivität}{f\colon R\bij S}
  \DeltaRow{\textbf{Neues Symbol}}{}
  \DeltaPrem{Ringisomorphismus}{%
    \RingIso{f}{R,\oplus,\odot,0_R,1_R}
      {S,\boxplus,\boxtimes,0_S,1_S}}
}

\FormulaThmDeltaK[Komposition von Ringhomomorphismen]{%
  \begin{aligned}[t]
    &\RingHom{f}{R,\oplus,\odot,0_R,1_R}
      {S,\boxplus,\boxtimes,0_S,1_S}\dsep{}\\[-2pt]
    &\RingHom{g}{S,\boxplus,\boxtimes,0_S,1_S}
      {T,\uplus,\otimes,0_T,1_T}\\
    &\vdash
      \RingHom{g\circ f}{R,\oplus,\odot,0_R,1_R}
      {T,\uplus,\otimes,0_T,1_T}
  \end{aligned}
}{RingHomComposition}{%
  \DeltaRow{Mengen}{R\dsep S\dsep T}
  \DeltaRow{Funktionen}{f\dsep g\dsep\oplus\dsep\odot\dsep
    \boxplus\dsep\boxtimes\dsep\uplus\dsep\otimes}
}
\begin{tabproofwide}
  \proofstepwidestar[1]{%
    \RingHom{f}{R,\oplus,\odot,0_R,1_R}
      {S,\boxplus,\boxtimes,0_S,1_S}}{\rA}
  \proofstepwidestar[2]{%
    \RingHom{g}{S,\boxplus,\boxtimes,0_S,1_S}
      {T,\uplus,\otimes,0_T,1_T}}{\rA}
  \proofstepwidestar[1]{%
    \GroupHom{f}{R,\oplus,0_R}{S,\boxplus,0_S}}{%
    \FormulaRefAuto{RingHomDef}[def]{1}}
  \proofstepwidestar[2]{%
    \GroupHom{g}{S,\boxplus,0_S}{T,\uplus,0_T}}{%
    \FormulaRefAuto{RingHomDef}[def]{2}}
  \proofstepwidestar[1,2]{%
    \GroupHom{g\circ f}{R,\oplus,0_R}{T,\uplus,0_T}}{%
    \FormulaRefAuto{GroupHomComposition}[thm]{3,4}}
  \proofstepwidestar[1]{%
    \MonoidHom{f}{R,\odot,1_R}{S,\boxtimes,1_S}}{%
    \FormulaRefAuto{RingHomDef}[def]{1}}
  \proofstepwidestar[2]{%
    \MonoidHom{g}{S,\boxtimes,1_S}{T,\otimes,1_T}}{%
    \FormulaRefAuto{RingHomDef}[def]{2}}
  \proofstepwidestar[1,2]{%
    \MonoidHom{g\circ f}{R,\odot,1_R}{T,\otimes,1_T}}{%
    \FormulaRefAuto{MonoidHomComposition}[thm]{6,7}}
  \proofstepwidestar[1]{%
    \RingAlg{R}{\oplus}{\odot}{0_R}{1_R}}{%
    \FormulaRefAuto{RingHomDef}[def]{1}}
  \proofstepwidestar[2]{%
    \RingAlg{T}{\uplus}{\otimes}{0_T}{1_T}}{%
    \FormulaRefAuto{RingHomDef}[def]{2}}
  \proofstepwidestar[1,2]{%
    \RingHom{g\circ f}{R,\oplus,\odot,0_R,1_R}
      {T,\uplus,\otimes,0_T,1_T}}{%
    \FormulaRefAuto{RingHomDef}[def]{9,10,5,8}}
\end{tabproofwide}

\section{Spezielle Definitionen}

Für kommutative Ringe bleiben die Erhaltungsgleichungen unverändert.
Die Kommutativität gehört zu den beiden beteiligten Ringstrukturen und
ist keine zusätzliche Forderung an die Abbildung.

\FormulaDefDeltaK[Homomorphismen kommutativer Ringe]{%
  \begin{aligned}[t]
    &\CommutativeRingHom{f}{R,\oplus,\odot,0_R,1_R}
      {S,\boxplus,\boxtimes,0_S,1_S}\\
    &\quad\coloneqq
      \CommutativeRingAlg{R}{\oplus}{\odot}{0_R}{1_R}\land
      \CommutativeRingAlg{S}{\boxplus}{\boxtimes}{0_S}{1_S}\\[-2pt]
    &\qquad\land
      \RingHom{f}{R,\oplus,\odot,0_R,1_R}
        {S,\boxplus,\boxtimes,0_S,1_S}
  \end{aligned}
}{CommutativeRingHomDef}{%
  \DeltaRow{Mengen}{R\dsep S\dsep0_R\dsep1_R\dsep0_S\dsep1_S}
  \DeltaRow{Funktionen}{f\dsep\oplus\dsep\odot\dsep\boxplus\dsep\boxtimes}
  \DeltaRow{\textbf{Neues Symbol}}{}
  \DeltaPrem{Symbol}{\mathsf{KRingHom}}
}

\FormulaAxiomDeltaK[Quellstruktur eines Homomorphismus kommutativer Ringe]{%
  \CommutativeRingHom{f}{R,\oplus,\odot,0_R,1_R}
    {S,\boxplus,\boxtimes,0_S,1_S}
  \vdash \CommutativeRingAlg{R}{\oplus}{\odot}{0_R}{1_R}%
}{CommutativeRingHomSourceAxiom}{%
  \DeltaRow{Mengen}{R\dsep S\dsep0_R\dsep1_R\dsep0_S\dsep1_S}
  \DeltaRow{Funktionen}{f\dsep\oplus\dsep\odot\dsep\boxplus\dsep\boxtimes}
}

\FormulaAxiomDeltaK[Zielstruktur eines Homomorphismus kommutativer Ringe]{%
  \CommutativeRingHom{f}{R,\oplus,\odot,0_R,1_R}
    {S,\boxplus,\boxtimes,0_S,1_S}
  \vdash \CommutativeRingAlg{S}{\boxplus}{\boxtimes}{0_S}{1_S}%
}{CommutativeRingHomTargetAxiom}{%
  \DeltaRow{Mengen}{R\dsep S\dsep0_R\dsep1_R\dsep0_S\dsep1_S}
  \DeltaRow{Funktionen}{f\dsep\oplus\dsep\odot\dsep\boxplus\dsep\boxtimes}
}

\FormulaAxiomDeltaK[Ringhomomorphismusanteil eines Homomorphismus kommutativer Ringe]{%
  \CommutativeRingHom{f}{R,\oplus,\odot,0_R,1_R}
    {S,\boxplus,\boxtimes,0_S,1_S}
  \vdash
  \RingHom{f}{R,\oplus,\odot,0_R,1_R}
    {S,\boxplus,\boxtimes,0_S,1_S}%
}{CommutativeRingHomBaseHomAxiom}{%
  \DeltaRow{Mengen}{R\dsep S\dsep0_R\dsep1_R\dsep0_S\dsep1_S}
  \DeltaRow{Funktionen}{f\dsep\oplus\dsep\odot\dsep\boxplus\dsep\boxtimes}
}

\FormulaDefDeltaK[Isomorphismen kommutativer Ringe]{%
  \begin{aligned}[t]
    &\CommutativeRingIso{f}{R,\oplus,\odot,0_R,1_R}
      {S,\boxplus,\boxtimes,0_S,1_S}\\
    &\quad\coloneqq
      \CommutativeRingHom{f}{R,\oplus,\odot,0_R,1_R}
        {S,\boxplus,\boxtimes,0_S,1_S}\\[-2pt]
    &\qquad\land f\colon R\bij S
  \end{aligned}
}{CommutativeRingIsoDef}{%
  \DeltaRow{Mengen}{R\dsep S\dsep0_R\dsep1_R\dsep0_S\dsep1_S}
  \DeltaRow{Funktionen}{f\dsep\oplus\dsep\odot\dsep\boxplus\dsep\boxtimes}
  \DeltaRow{\textbf{Neues Symbol}}{}
  \DeltaPrem{Symbol}{\mathsf{KRingIso}}
}

\FormulaAxiomDeltaK[Homomorphismusanteil eines Isomorphismus kommutativer Ringe]{%
  \CommutativeRingIso{f}{R,\oplus,\odot,0_R,1_R}
    {S,\boxplus,\boxtimes,0_S,1_S}
  \vdash
  \CommutativeRingHom{f}{R,\oplus,\odot,0_R,1_R}
    {S,\boxplus,\boxtimes,0_S,1_S}%
}{CommutativeRingIsoHomAxiom}{%
  \DeltaRow{Mengen}{R\dsep S\dsep0_R\dsep1_R\dsep0_S\dsep1_S}
  \DeltaRow{Funktionen}{f\dsep\oplus\dsep\odot\dsep\boxplus\dsep\boxtimes}
}

\FormulaAxiomDeltaK[Bijektivität eines Isomorphismus kommutativer Ringe]{%
  \CommutativeRingIso{f}{R,\oplus,\odot,0_R,1_R}
    {S,\boxplus,\boxtimes,0_S,1_S}
  \vdash f\colon R\bij S%
}{CommutativeRingIsoBijectiveAxiom}{%
  \DeltaRow{Mengen}{R\dsep S\dsep0_R\dsep1_R\dsep0_S\dsep1_S}
  \DeltaRow{Funktionen}{f\dsep\oplus\dsep\odot\dsep\boxplus\dsep\boxtimes}
}

Die Kommutativität der Multiplikation ist wiederum kein zusätzliches
Erhaltungsgesetz. Sie wird durch einen Ringisomorphismus in beide Richtungen
transportiert.

\chapter{Die algebraische Struktur der ganzen Zahlen}

In den folgenden Beweisen bezeichnen \(\IntZero\) und
\(\IntOne\) die in Band~17 konstruierten neutralen Elemente. Sämtliche
Rechengesetze werden aus den dort registrierten Sätzen übernommen; die
Schlusszeilen verwenden ausschließlich die abstrakten Strukturdefinitionen
dieses und des vorangehenden Bandes.

\section{Die additive abelsche Gruppe}

\FormulaThmDeltaK[Additive abelsche Gruppe der ganzen Zahlen]{%
  \AbelianGroupAlg{\mathbb Z}{+}{\IntZero}%
}{IntegerAdditiveAbelianGroup}{%
  \DeltaRow{Mengen}{\mathbb Z\dsep \IntZero\dsep z\dsep w\dsep u}
  \DeltaRow{Funktionen}{+}
}
\begin{tabproofsplitwide}
  \proofpartwideindR[Additive Halbgruppe]{%
    \SemigroupAlg{\mathbb Z}{+}%
  }{\SemigroupAlg{\mathbb Z}{+}}[IntegerAdditiveSemigroup]
    \proofstepwidestar[1]{z\in\mathbb Z}{\rA}
    \proofstepwidestar[2]{w\in\mathbb Z}{\rA}
    \proofstepwidestar[1,2]{z+w\in\mathbb Z}{%
      \FormulaRefAuto{IntegerAddClosureAxiom}[ax]{1,2}}
    \proofstepwidestar[4]{u\in\mathbb Z}{\rA}
    \proofstepwidestar[1,2,4]{(z+w)+u=z+(w+u)}{%
      \FormulaRefAuto{IntegerAddAssociativeAxiom}[ax]{1,2,4}}
    \proofstepwidestar[]{\SemigroupAlg{\mathbb Z}{+}}{%
      \FormulaRefAuto{SemigroupDef}[def]{3,5}}
  \closeproofpartwideindR

  \proofpartwideindR[Additives Monoid]{%
    \MonoidAlg{\mathbb Z}{+}{\IntZero}%
  }{\MonoidAlg{\mathbb Z}{+}{\IntZero}}[IntegerAdditiveMonoid]
    \proofstepwidestar[]{\SemigroupAlg{\mathbb Z}{+}}{%
      \FormulaRefAuto{IntegerAdditiveSemigroup}[thm]}
    \proofstepwidestar[]{\IntZero\in\mathbb Z}{%
      \FormulaRefAuto{IntZeroCarrier}[thm]}
    \proofstepwidestar[3]{z\in\mathbb Z}{\rA}
    \proofstepwidestar[3]{\IntZero+z=z}{%
      \rAEb{\FormulaRefAuto{IntegerAddZeroAxiom}[ax]{3}}}
    \proofstepwidestar[3]{z+\IntZero=z}{%
      \rAEa{\FormulaRefAuto{IntegerAddZeroAxiom}[ax]{3}}}
    \proofstepwidestar[]{\MonoidAlg{\mathbb Z}{+}{\IntZero}}{%
      \FormulaRefAuto{MonoidDef}[def]{1,2,4,5}}
  \closeproofpartwideindR

  \proofpartwideindR[Additive Gruppe]{%
    \GroupAlg{\mathbb Z}{+}{\IntZero}%
  }{\GroupAlg{\mathbb Z}{+}{\IntZero}}[IntegerAdditiveGroup]
    \proofstepwidestar[]{\MonoidAlg{\mathbb Z}{+}{\IntZero}}{%
      \FormulaRefAuto{IntegerAdditiveMonoid}[thm]}
    \proofstepwidestar[1]{z\in\mathbb Z}{\rA}
    \proofstepwidestar[1]{-z\in\mathbb Z}{%
      \FormulaRefAuto{IntNegClosure}[thm]{1}}
    \proofstepwidestar[1]{%
      z+(-z)=\IntZero\land(-z)+z=\IntZero}{%
      \FormulaRefAuto{IntegerAddInverseAxiom}[ax]{1}}
    \proofstepwidestar[1]{%
      \exists w\in\mathbb Z\,
      \bigl(z+w=\IntZero\land w+z=\IntZero\bigr)}{%
      \rEI{\rAI{3,4}}}
    \proofstepwidestar[]{\GroupAlg{\mathbb Z}{+}{\IntZero}}{%
      \FormulaRefAuto{GroupDef}[def]{1,5}}
  \closeproofpartwideindR

  \proofpartwide{Kommutativität}
    \proofstepwidestar[]{\GroupAlg{\mathbb Z}{+}{\IntZero}}{%
      \FormulaRefAuto{IntegerAdditiveGroup}[thm]}
    \proofstepwidestar[2]{z\in\mathbb Z}{\rA}
    \proofstepwidestar[3]{w\in\mathbb Z}{\rA}
    \proofstepwidestar[2,3]{z+w=w+z}{%
      \FormulaRefAuto{IntegerAddCommutativeAxiom}[ax]{2,3}}
    \proofstepwidestar[]{%
      \AbelianGroupAlg{\mathbb Z}{+}{\IntZero}}{%
      \FormulaRefAuto{AbelianGroupDef}[def]{1,4}}
  \closeproofpartwide
\end{tabproofsplitwide}

\section{Das multiplikative kommutative Monoid}

\FormulaThmDeltaK[Multiplikatives kommutatives Monoid der ganzen Zahlen]{%
  \CommutativeMonoidAlg{\mathbb Z}{\cdot}{\IntOne}%
}{IntegerMultiplicativeCommutativeMonoid}{%
  \DeltaRow{Mengen}{\mathbb Z\dsep \IntOne\dsep z\dsep w\dsep u}
  \DeltaRow{Funktionen}{\cdot}
}
\begin{tabproofsplitwide}
  \proofpartwideindR[Multiplikative Halbgruppe]{%
    \SemigroupAlg{\mathbb Z}{\cdot}%
  }{\SemigroupAlg{\mathbb Z}{\cdot}}[IntegerMultiplicativeSemigroup]
    \proofstepwidestar[1]{z\in\mathbb Z}{\rA}
    \proofstepwidestar[2]{w\in\mathbb Z}{\rA}
    \proofstepwidestar[1,2]{z\cdot w\in\mathbb Z}{%
      \FormulaRefAuto{IntegerMulClosureAxiom}[ax]{1,2}}
    \proofstepwidestar[4]{u\in\mathbb Z}{\rA}
    \proofstepwidestar[1,2,4]{%
      (z\cdot w)\cdot u=z\cdot(w\cdot u)}{%
      \FormulaRefAuto{IntegerMulAssociativeAxiom}[ax]{1,2,4}}
    \proofstepwidestar[]{\SemigroupAlg{\mathbb Z}{\cdot}}{%
      \FormulaRefAuto{SemigroupDef}[def]{3,5}}
  \closeproofpartwideindR

  \proofpartwideindR[Multiplikatives Monoid]{%
    \MonoidAlg{\mathbb Z}{\cdot}{\IntOne}%
  }{\MonoidAlg{\mathbb Z}{\cdot}{\IntOne}}[IntegerMultiplicativeMonoid]
    \proofstepwidestar[]{\SemigroupAlg{\mathbb Z}{\cdot}}{%
      \FormulaRefAuto{IntegerMultiplicativeSemigroup}[thm]}
    \proofstepwidestar[]{\IntOne\in\mathbb Z}{%
      \FormulaRefAuto{IntOneCarrier}[thm]}
    \proofstepwidestar[3]{z\in\mathbb Z}{\rA}
    \proofstepwidestar[3]{\IntOne\cdot z=z}{%
      \rAEb{\FormulaRefAuto{IntegerMulOneAxiom}[ax]{3}}}
    \proofstepwidestar[3]{z\cdot\IntOne=z}{%
      \rAEa{\FormulaRefAuto{IntegerMulOneAxiom}[ax]{3}}}
    \proofstepwidestar[]{\MonoidAlg{\mathbb Z}{\cdot}{\IntOne}}{%
      \FormulaRefAuto{MonoidDef}[def]{1,2,4,5}}
  \closeproofpartwideindR

  \proofpartwide{Kommutativität}
    \proofstepwidestar[]{\MonoidAlg{\mathbb Z}{\cdot}{\IntOne}}{%
      \FormulaRefAuto{IntegerMultiplicativeMonoid}[thm]}
    \proofstepwidestar[2]{z\in\mathbb Z}{\rA}
    \proofstepwidestar[3]{w\in\mathbb Z}{\rA}
    \proofstepwidestar[2,3]{z\cdot w=w\cdot z}{%
      \FormulaRefAuto{IntegerMulCommutativeAxiom}[ax]{2,3}}
    \proofstepwidestar[]{%
      \CommutativeMonoidAlg{\mathbb Z}{\cdot}{\IntOne}}{%
      \FormulaRefAuto{CommutativeMonoidDef}[def]{1,4}}
  \closeproofpartwide
\end{tabproofsplitwide}

\section{Der kommutative Ring mit Eins}

\FormulaThmDeltaK[Kommutativer Ring der ganzen Zahlen]{%
  \CommutativeRingAlg{\mathbb Z}{+}{\cdot}
    {\IntZero}{\IntOne}%
}{IntegerCommutativeRing}{%
  \DeltaRow{Mengen}{%
    \mathbb Z\dsep\IntZero\dsep\IntOne
    \dsep z\dsep w\dsep u}
  \DeltaRow{Funktionen}{+\dsep\cdot}
}
\begin{tabproofsplitwide}
  \proofpartwideindR[Ring mit Eins]{%
    \RingAlg{\mathbb Z}{+}{\cdot}{\IntZero}{\IntOne}%
  }{%
    \RingAlg{\mathbb Z}{+}{\cdot}{\IntZero}{\IntOne}%
  }[IntegerUnitalRing]
    \proofstepwidestar[]{%
      \AbelianGroupAlg{\mathbb Z}{+}{\IntZero}}{%
      \FormulaRefAuto{IntegerAdditiveAbelianGroup}[thm]}
    \proofstepwidestar[]{%
      \CommutativeMonoidAlg{\mathbb Z}{\cdot}{\IntOne}}{%
      \FormulaRefAuto{IntegerMultiplicativeCommutativeMonoid}[thm]}
    \proofstepwidestar[]{%
      \MonoidAlg{\mathbb Z}{\cdot}{\IntOne}}{%
      \FormulaRefAuto{CommutativeMonoidBaseMonoidAxiom}[ax]{2}}
    \proofstepwidestar[4]{z\in\mathbb Z}{\rA}
    \proofstepwidestar[5]{w\in\mathbb Z}{\rA}
    \proofstepwidestar[6]{u\in\mathbb Z}{\rA}
    \proofstepwidestar[4,5,6]{%
      z\cdot(w+u)=z\cdot w+z\cdot u}{%
      \FormulaRefAuto{IntegerMulLeftDistributiveAxiom}[ax]{4,5,6}}
    \proofstepwidestar[4,5,6]{%
      (z+w)\cdot u=z\cdot u+w\cdot u}{%
      \FormulaRefAuto{IntegerMulRightDistributiveAxiom}[ax]{4,5,6}}
    \proofstepwidestar[]{%
      \RingAlg{\mathbb Z}{+}{\cdot}{\IntZero}{\IntOne}}{%
      \FormulaRefAuto{RingDef}[def]{1,3,7,8}}
  \closeproofpartwideindR

  \proofpartwide{Kommutative Multiplikation}
    \proofstepwidestar[]{%
      \RingAlg{\mathbb Z}{+}{\cdot}{\IntZero}{\IntOne}}{%
      \FormulaRefAuto{IntegerUnitalRing}[thm]}
    \proofstepwidestar[2]{z\in\mathbb Z}{\rA}
    \proofstepwidestar[3]{w\in\mathbb Z}{\rA}
    \proofstepwidestar[2,3]{z\cdot w=w\cdot z}{%
      \FormulaRefAuto{IntegerMulCommutativeAxiom}[ax]{2,3}}
    \proofstepwidestar[]{%
      \CommutativeRingAlg{\mathbb Z}{+}{\cdot}
        {\IntZero}{\IntOne}}{%
      \FormulaRefAuto{CommutativeRingDef}[def]{1,4}}
  \closeproofpartwide
\end{tabproofsplitwide}

\section{Der kanonische Ringhomomorphismus aus den ganzen Zahlen}

Sei \(\RingAlg{R}{\oplus}{\odot}{0_R}{1_R}\). Nach
\FormulaRefAuto{RingUnderlyingSemiring}[thm] ist \(R\) auch ein Halbring,
also steht die kanonische Abbildung
\(\NaturalSemiringMap{R,\oplus,\odot,0_R,1_R}\colon\mathbb N\to R\)
aus Band~32 zur Verfügung. Ein Differenzenpaar \((a,b)\) soll auf die
Differenz seiner beiden Bilder abgebildet werden.

\subsection{Repräsentantenunabhängigkeit}

\FormulaThmDeltaK[Kreuzsummen liefern gleiche Differenzen im Ring]{%
  \begin{aligned}[t]
    &\RingAlg{R}{\oplus}{\odot}{0_R}{1_R}\dsep
      p,q,r,s\in R\dsep p\oplus s=q\oplus r\\
    &\vdash p\oplus(-q)=r\oplus(-s)
  \end{aligned}
}{RingCrossSumGivesDifferenceEquality}{%
  \DeltaRow{Mengen}{R\dsep0_R\dsep1_R\dsep p\dsep q\dsep r\dsep s}
  \DeltaRow{Funktionen}{\oplus\dsep\odot}
}
Im folgenden Beweis kürzen wir
\(\bar q:=-q\) und \(\bar s:=-s\) ab.
\begin{tabproofwide}
  \proofstepwidestar[1]{\RingAlg{R}{\oplus}{\odot}{0_R}{1_R}}{\rA}
  \proofstepwidestar[2]{p\in R}{\rA}
  \proofstepwidestar[3]{q\in R}{\rA}
  \proofstepwidestar[4]{r\in R}{\rA}
  \proofstepwidestar[5]{s\in R}{\rA}
  \proofstepwidestar[6]{p\oplus s=q\oplus r}{\rA}
  \proofstepwidestar[1]{\AbelianGroupAlg{R}{\oplus}{0_R}}{%
    \FormulaRefAuto{RingAdditiveAbelianGroupAxiom}[ax]{1}}
  \proofstepwidestar[1]{\GroupAlg{R}{\oplus}{0_R}}{%
    \FormulaRefAuto{AbelianGroupBaseGroupAxiom}[ax]{7}}
  \proofstepwidestar[1]{\CommutativeMonoidAlg{R}{\oplus}{0_R}}{%
    \FormulaRefAuto{AbelianGroupIsCommutativeMonoid}[thm]{7}}
  \proofstepwidestar[1]{\CommutativeSemigroupAlg{R}{\oplus}}{%
    \FormulaRefAuto{CommutativeMonoidIsCommutativeSemigroup}[thm]{9}}
  \proofstepwidestar[1,3]{\bar q\in R}{%
    \FormulaRefAuto{RingNegativeCarrier}[thm]{1,3}}
  \proofstepwidestar[1,5]{\bar s\in R}{%
    \FormulaRefAuto{RingNegativeCarrier}[thm]{1,5}}
  \proofstepwide[1,2,11]{p\oplus\bar q}{=}{%
    (p\oplus\bar q)\oplus0_R}{%
    \parbox[t]{.22\textwidth}{\raggedright
      \FormulaRefAuto{a=b\vdash b=a}\\
      \FormulaRefAuto{GroupRightIdentity}[thm]\\
      \FormulaRefAuto{GroupClosure}[thm]}}
  \proofstepwide[1,5]{(p\oplus\bar q)\oplus0_R}{=}{%
    (p\oplus\bar q)\oplus(s\oplus\bar s)}{%
    \parbox[t]{.22\textwidth}{\raggedright
      \FormulaRefAuto{a=b\vdash b=a}\\
      \FormulaRefAuto{RingNegativeRight}[thm]}}
  \proofstepwide[1,2,5,11,12]{%
    (p\oplus\bar q)\oplus(s\oplus\bar s)}{=}{%
    \bigl((p\oplus s)\oplus\bar q\bigr)\oplus\bar s}{%
    \parbox[t]{.22\textwidth}{\raggedright
      \FormulaRefAuto{SemigroupAssociativityAxiom}[ax]\\
      \FormulaRefAuto{CommutativeSemigroupCommutativityAxiom}[ax]}}
  \proofstepwide[6]{%
    \bigl((p\oplus s)\oplus\bar q\bigr)\oplus\bar s}{=}{%
    \bigl((q\oplus r)\oplus\bar q\bigr)\oplus\bar s}{\rIE{6}}
  \proofstepwide[1,3,4,11,12]{%
    \bigl((q\oplus r)\oplus\bar q\bigr)\oplus\bar s}{=}{%
    \bigl((q\oplus\bar q)\oplus r\bigr)\oplus\bar s}{%
    \FormulaRefAuto{CommutativeSemigroupThreeTermExchange}[thm]{%
      10,3,4,11}}
  \proofstepwide[1,3,4,12]{%
    \bigl((q\oplus\bar q)\oplus r\bigr)\oplus\bar s}{=}{%
    (0_R\oplus r)\oplus\bar s}{%
    \FormulaRefAuto{RingNegativeRight}[thm]{1,3}}
  \proofstepwide[1,4,12]{(0_R\oplus r)\oplus\bar s}{=}{%
    r\oplus\bar s}{%
    \FormulaRefAuto{GroupLeftIdentity}[thm]{8,4}}
  \proofstepwidestar[1,2,3,4,5,6]{%
    p\oplus\bar q=r\oplus\bar s}{\rChain{13,14,15,16,17,18,19}}
\end{tabproofwide}

\FormulaDefDeltaK[Wert eines Differenzenpaares im Zielring]{%
  \begin{aligned}[t]
    &\RingAlg{R}{\oplus}{\odot}{0_R}{1_R}\dsep a,b\in\mathbb N
    \vdash{}\\[-2pt]
    &\IntegerRingPairValue{R,\oplus,\odot,0_R,1_R}{a}{b}
      \coloneqq
      \NaturalSemiringMap{R,\oplus,\odot,0_R,1_R}(a)
      \oplus
      \bigl(
        -\NaturalSemiringMap{R,\oplus,\odot,0_R,1_R}(b)
      \bigr)
  \end{aligned}
}{IntegerRingPairValueDef}{%
  \DeltaRow{Mengen}{R\dsep0_R\dsep1_R\dsep a\dsep b}
  \DeltaRow{Funktionen}{\oplus\dsep\odot}
  \DeltaRow{\textbf{Neues Symbol}}{}
  \DeltaPrem{Wert eines Differenzenpaares}{%
    \IntegerRingPairValue{R,\oplus,\odot,0_R,1_R}{a}{b}}
}

\FormulaThmDeltaK[Der Wert eines Differenzenpaares liegt im Zielring]{%
  \RingAlg{R}{\oplus}{\odot}{0_R}{1_R}\dsep a,b\in\mathbb N
  \vdash
  \IntegerRingPairValue{R,\oplus,\odot,0_R,1_R}{a}{b}\in R%
}{IntegerRingPairValueCarrier}{%
  \DeltaRow{Mengen}{R\dsep a\dsep b}
  \DeltaRow{Funktionen}{\oplus\dsep\odot}
}
Für diesen Beweis setzen wir
\[
  \begin{aligned}
    \eta_R(n)
      &:=\NaturalSemiringMap{R,\oplus,\odot,0_R,1_R}(n),\\
    v_R(a,b)
      &:=\IntegerRingPairValue{R,\oplus,\odot,0_R,1_R}{a}{b}.
  \end{aligned}
\]
\begin{tabproof}
  \proofstep{1}{\RingAlg{R}{\oplus}{\odot}{0_R}{1_R}}{\rA}
  \proofstep{2}{a\in\mathbb N}{\rA}
  \proofstep{3}{b\in\mathbb N}{\rA}
  \proofstep{1}{\SemiringAlg{R}{\oplus}{\odot}{0_R}{1_R}}{%
    \FormulaRefAuto{RingUnderlyingSemiring}[thm]{1}}
  \proofstep{1}{\eta_R\colon\mathbb N\to R}{%
    \FormulaRefAuto{NaturalSemiringMapFunction}[thm]{4}}
  \proofstep{1,2}{\eta_R(a)\in R}{%
    \FormulaRefAuto{F\colon A\to B\dsep a\in A\vdash F(a)\in B}{5,2}}
  \proofstep{1,3}{\eta_R(b)\in R}{%
    \FormulaRefAuto{F\colon A\to B\dsep a\in A\vdash F(a)\in B}{5,3}}
  \proofstep{1,3}{-\eta_R(b)\in R}{%
    \FormulaRefAuto{RingNegativeCarrier}[thm]{1,7}}
  \proofstep{1,2,3}{%
    \eta_R(a)\oplus\bigl(-\eta_R(b)\bigr)\in R}{%
    \begin{gathered}
      \FormulaRefAuto{GroupClosure}[thm]\\
      \FormulaRefAuto{RingAdditiveGroup}[thm]
    \end{gathered}}
  \proofstep{1,2,3}{v_R(a,b)=\eta_R(a)\oplus\bigl(-\eta_R(b)\bigr)}{%
    \FormulaRefAuto{IntegerRingPairValueDef}[def]{1,2,3}}
  \proofstep{1,2,3}{v_R(a,b)\in R}{%
    \rIE{10,9}}
\end{tabproof}

\FormulaThmDeltaK[Repräsentantenunabhängigkeit des Ganzzahlbildes]{%
  \begin{aligned}[t]
    &\RingAlg{R}{\oplus}{\odot}{0_R}{1_R}\dsep
      a,b,c,d\in\mathbb N\dsep
      \IntClass{a}{b}=\IntClass{c}{d}\\
    &\vdash
      \IntegerRingPairValue{R,\oplus,\odot,0_R,1_R}{a}{b}
      =
      \IntegerRingPairValue{R,\oplus,\odot,0_R,1_R}{c}{d}
  \end{aligned}
}{IntegerRingMapRepresentativeIndependent}{%
  \DeltaRow{Mengen}{R\dsep a\dsep b\dsep c\dsep d}
  \DeltaRow{Funktionen}{\oplus\dsep\odot}
}
Zur Entlastung der Beweiszeilen schreiben wir hier
\[
  \begin{aligned}
    \eta_R(n)
      &:=\NaturalSemiringMap{R,\oplus,\odot,0_R,1_R}(n),\\
    v_R(x,y)
      &:=\IntegerRingPairValue{R,\oplus,\odot,0_R,1_R}{x}{y}.
  \end{aligned}
\]
\begin{tabproofwide}
  \proofstepwidestar[1]{\RingAlg{R}{\oplus}{\odot}{0_R}{1_R}}{\rA}
  \proofstepwidestar[2]{a,b,c,d\in\mathbb N}{\rA}
  \proofstepwidestar[3]{\IntClass{a}{b}=\IntClass{c}{d}}{\rA}
  \proofstepwidestar[2,3]{a+d=b+c}{%
    \FormulaRefAuto{IntegerClassEquality}[thm]{2,3}}
  \proofstepwidestar[1]{\SemiringAlg{R}{\oplus}{\odot}{0_R}{1_R}}{%
    \FormulaRefAuto{RingUnderlyingSemiring}[thm]{1}}
  \proofstepwidestar[1,2]{%
    \eta_R(a+d)=\eta_R(a)\oplus\eta_R(d)}{%
    \FormulaRefAuto{NaturalSemiringMapAddition}[thm]{5,2}}
  \proofstepwidestar[1,2]{%
    \eta_R(b+c)=\eta_R(b)\oplus\eta_R(c)}{%
    \FormulaRefAuto{NaturalSemiringMapAddition}[thm]{5,2}}
  \proofstepwidestar[1,2,3]{%
    \eta_R(a+d)=\eta_R(b+c)}{\rIE{4}}
  \proofstepwidestar[1,2,3]{%
    \eta_R(a)\oplus\eta_R(d)
    =
    \eta_R(b)\oplus\eta_R(c)}{%
    \rIE{6,7,8}}
  \proofstepwidestar[1,2,3]{%
    \eta_R(a)\oplus\bigl(-\eta_R(b)\bigr)
    =
    \eta_R(c)\oplus\bigl(-\eta_R(d)\bigr)}{%
    \parbox[t]{.25\textwidth}{\raggedright
      \FormulaRefAuto{RingCrossSumGivesDifferenceEquality}[thm]\\
      mit \FormulaRefAuto{NaturalSemiringMapFunction}[thm]\\
      und \FormulaRefAuto{F\colon A\to B\dsep a\in A\vdash F(a)\in B}}}
  \proofstepwidestar[1,2]{v_R(a,b)=\eta_R(a)\oplus\bigl(-\eta_R(b)\bigr)}{%
    \FormulaRefAuto{IntegerRingPairValueDef}[def]{1,2}}
  \proofstepwidestar[1,2]{v_R(c,d)=\eta_R(c)\oplus\bigl(-\eta_R(d)\bigr)}{%
    \FormulaRefAuto{IntegerRingPairValueDef}[def]{1,2}}
  \proofstepwidestar[1,2,3]{v_R(a,b)=v_R(c,d)}{%
    \rIE{11,12,10}}
\end{tabproofwide}

\subsection{Quotientenabstieg}

\FormulaThmDeltaK[Eindeutiger Quotientenabstieg des Ganzzahlbildes]{%
  \begin{aligned}[t]
    &\RingAlg{R}{\oplus}{\odot}{0_R}{1_R}\vdash{}\\[-2pt]
    &\exists! f\Bigl(
      f\colon\mathbb Z\to R\land
      \forall a,b\in\mathbb N\,
      f(\IntClass{a}{b})
      =
      \IntegerRingPairValue{R,\oplus,\odot,0_R,1_R}{a}{b}
    \Bigr)
  \end{aligned}
}{IntegerRingMapQuotientDescent}{%
  \DeltaRow{Mengen}{%
    R\dsep a\dsep b\dsep c\dsep d\dsep z\dsep u\dsep u_1\dsep u_2}
  \DeltaRow{Funktionen}{f\dsep g\dsep\oplus\dsep\odot}
  \DeltaRow{Repräsentanten}{%
    \FormulaRefAuto{IntegerPairRepresentative}[thm]}
  \DeltaRow{Werte im Zielring}{%
    \FormulaRefAuto{IntegerRingPairValueCarrier}[thm]}
  \DeltaRow{Repräsentantenunabhängigkeit}{%
    \FormulaRefAuto{IntegerRingMapRepresentativeIndependent}[thm]}
}

Für den folgenden Beweis setzen wir lokal
\[
  \begin{aligned}
    v_R(c,d)
      &:=
        \IntegerRingPairValue{R,\oplus,\odot,0_R,1_R}{c}{d},\\
    G_R
      :=\Bigl\{(z,u)\in\mathbb Z\times R\ \Bigm|\
        &\exists c,d\in\mathbb N\,
        \Bigl(z=\IntClass{c}{d}\land{}\\[-2pt]
        &\qquad
        u=v_R(c,d)
        \Bigr)
      \Bigr\}.
  \end{aligned}
\]

\begin{tabproofsplitwide}
  \proofpartwideindR[Eindeutiger Zielwert jeder Ganzzahl]{%
    \begin{aligned}[t]
      &\RingAlg{R}{\oplus}{\odot}{0_R}{1_R}\dsep z\in\mathbb Z\\
      &\vdash
      \exists! u\in R\,\exists a,b\in\mathbb N\,
      \Bigl(
        z=\IntClass{a}{b}\land
        u=v_R(a,b)
      \Bigr)
    \end{aligned}
  }{%
    \begin{aligned}[t]
      &\RingAlg{R}{\oplus}{\odot}{0_R}{1_R}\dsep z\in\mathbb Z\\
      &\vdash
      \exists! u\in R\,\exists a,b\in\mathbb N\,
      \Bigl(
        z=\IntClass{a}{b}\land
        u=v_R(a,b)
      \Bigr)
    \end{aligned}
  }[IntegerRingMapTargetValueUnique]
    \proofstepwidestar[1]{%
      \RingAlg{R}{\oplus}{\odot}{0_R}{1_R}}{\rA}
    \proofstepwidestar[2]{z\in\mathbb Z}{\rA}
    \proofstepwidestar[2]{%
      \exists a,b\in\mathbb N\,z=\IntClass{a}{b}}{%
      \FormulaRefAuto{IntegerPairRepresentative}[thm]{2}}
    \proofstepwidestar[4]{%
      a,b\in\mathbb N\land z=\IntClass{a}{b}}{\rA}
    \proofstepwidestar[1,4]{%
      v_R(a,b)\in R}{%
      \FormulaRefAuto{IntegerRingPairValueCarrier}[thm]{1,4}}
    \proofstepwidestar[]{%
      v_R(a,b)
      =
      v_R(a,b)}{\rII}
    \proofstepwidestar[4]{%
      \exists c,d\in\mathbb N\,
      \Bigl(
        z=\IntClass{c}{d}\land
          v_R(a,b)
        =
          v_R(c,d)
      \Bigr)}{\rEI{\rAI{4,6}}}
    \proofstepwidestar[1,4]{%
      \begin{aligned}[t]
        &v_R(a,b)\in R
        \land{}\\[-2pt]
        &\exists c,d\in\mathbb N\,
        \Bigl(
          z=\IntClass{c}{d}\land
          v_R(a,b)
          =
          v_R(c,d)
        \Bigr)
      \end{aligned}}{\rAI{5,7}}
    \proofstepwidestar[1,4]{%
      \exists u\in R\,\exists c,d\in\mathbb N\,
      \Bigl(
        z=\IntClass{c}{d}\land
        u=v_R(c,d)
      \Bigr)}{\rEI{8}}
    \proofstepwidestar[1,2]{%
      \exists u\in R\,\exists a,b\in\mathbb N\,
      \Bigl(
        z=\IntClass{a}{b}\land
        u=v_R(a,b)
      \Bigr)}{\rEE{3,4,9}}
    \proofstepwidestar[11]{%
      u_1\in R\land
      \exists a,b\in\mathbb N\,
      \Bigl(
        z=\IntClass{a}{b}\land
        u_1=v_R(a,b)
      \Bigr)}{\rA}
    \proofstepwidestar[12]{%
      u_2\in R\land
      \exists c,d\in\mathbb N\,
      \Bigl(
        z=\IntClass{c}{d}\land
        u_2=v_R(c,d)
      \Bigr)}{\rA}
    \proofstepwidestar[11]{%
      \exists a,b\in\mathbb N\,
      \Bigl(
        z=\IntClass{a}{b}\land
        u_1=v_R(a,b)
      \Bigr)}{\rAEb{11}}
    \proofstepwidestar[12]{%
      \exists c,d\in\mathbb N\,
      \Bigl(
        z=\IntClass{c}{d}\land
        u_2=v_R(c,d)
      \Bigr)}{\rAEb{12}}
    \proofstepwidestar[15]{%
      a,b\in\mathbb N\land z=\IntClass{a}{b}\land
      u_1=v_R(a,b)}{\rA}
    \proofstepwidestar[16]{%
      c,d\in\mathbb N\land z=\IntClass{c}{d}\land
      u_2=v_R(c,d)}{\rA}
    \proofstepwidestar[15]{a,b\in\mathbb N}{\rAEa{15}}
    \proofstepwidestar[15]{z=\IntClass{a}{b}}{\rAEa{\rAEb{15}}}
    \proofstepwidestar[15]{%
      u_1=v_R(a,b)}{%
      \rAEb{\rAEb{15}}}
    \proofstepwidestar[16]{c,d\in\mathbb N}{\rAEa{16}}
    \proofstepwidestar[16]{z=\IntClass{c}{d}}{\rAEa{\rAEb{16}}}
    \proofstepwidestar[16]{%
      u_2=v_R(c,d)}{%
      \rAEb{\rAEb{16}}}
    \proofstepwidestar[15,16]{a,b,c,d\in\mathbb N}{\rAI{17,20}}
    \proofstepwidestar[15,16]{%
      \IntClass{a}{b}=\IntClass{c}{d}}{\rIE{18,21}}
    \proofstepwidestar[1,15,16]{%
      v_R(a,b)
      =
      v_R(c,d)}{%
      \FormulaRefAuto{IntegerRingMapRepresentativeIndependent}[thm]{%
        1,23,24}}
    \proofstepwidestar[1,15,16]{u_1=u_2}{\rIE{19,22,25}}
    \proofstepwidestar[1,12,15]{u_1=u_2}{\rEE{14,16,26}}
    \proofstepwidestar[1,11,12]{u_1=u_2}{\rEE{13,15,27}}
    \proofstepwidestar[1,2]{%
      \exists! u\in R\,\exists a,b\in\mathbb N\,
      \Bigl(
        z=\IntClass{a}{b}\land
        u=v_R(a,b)
      \Bigr)}{\UEI{10,11,12,28}}
  \closeproofpartwideindR

  \proofpartwideindR[Existenz der Klassenabbildung]{%
    \begin{aligned}[t]
      &\RingAlg{R}{\oplus}{\odot}{0_R}{1_R}\vdash{}\\[-2pt]
      &\exists f\Bigl(
        f\colon\mathbb Z\to R\land
        \forall a,b\in\mathbb N\,
        f(\IntClass{a}{b})
        =
        v_R(a,b)
      \Bigr)
    \end{aligned}
  }{%
    \begin{aligned}[t]
      &\RingAlg{R}{\oplus}{\odot}{0_R}{1_R}\vdash{}\\[-2pt]
      &\exists f\Bigl(
        f\colon\mathbb Z\to R\land
        \forall a,b\in\mathbb N\,
        f(\IntClass{a}{b})
        =
        v_R(a,b)
      \Bigr)
    \end{aligned}
  }[IntegerRingMapDescentExistence]
    \proofstepwidestar[1]{%
      \RingAlg{R}{\oplus}{\odot}{0_R}{1_R}}{\rA}
    \proofstepwidestar[]{G_R\subseteq\mathbb Z\times R}{%
      \FormulaRefAuto{\{x\in A\mid P(x)\}\subseteq A}[thm]}
    \proofstepwidestar[3]{z\in\mathbb Z}{\rA}
    \proofstepwidestar[1,3]{%
      \exists! u\in R\,\exists a,b\in\mathbb N\,
      \Bigl(
        z=\IntClass{a}{b}\land
        u=v_R(a,b)
      \Bigr)}{%
      \FormulaRefAuto{IntegerRingMapTargetValueUnique}[thm]{1,3}}
    \proofstepwidestar[1,3]{%
      \exists u\in R\,\exists a,b\in\mathbb N\,
      \Bigl(
        z=\IntClass{a}{b}\land
        u=v_R(a,b)
      \Bigr)}{%
      \FormulaRefAuto{\exists! x\,P(x)\vdash\exists x\,P(x)}{4}}
    \proofstepwidestar[6]{%
      u\in R\land
      \exists a,b\in\mathbb N\,
      \Bigl(
        z=\IntClass{a}{b}\land
        u=v_R(a,b)
      \Bigr)}{\rA}
    \proofstepwidestar[3,6]{(z,u)\in\mathbb Z\times R}{%
      \FormulaRefAuto{a\in A\dsep b\in B\vdash(a,b)\in A\times B}{%
        3,\rAEa{6}}}
    \proofstepwidestar[3,6]{(z,u)\in G_R}{%
      \FormulaRefAuto{x\in A\dsep P(x)\vdash
        x\in\{x\in A\mid P(x)\}}{7,\rAEb{6}}}
    \proofstepwidestar[3,6]{%
      u\in R\land (z,u)\in G_R}{\rAI{\rAEa{6},8}}
    \proofstepwidestar[3,6]{%
      \exists u\in R\,(z,u)\in G_R}{\rEI{9}}
    \proofstepwidestar[1,3]{%
      \exists u\in R\,(z,u)\in G_R}{\rEE{5,6,10}}
    \proofstepwidestar[12]{%
      u_1\in R\land (z,u_1)\in G_R}{\rA}
    \proofstepwidestar[13]{%
      u_2\in R\land (z,u_2)\in G_R}{\rA}
    \proofstepwidestar[12]{%
      \begin{aligned}[t]
        &\exists a,b\in\mathbb N\,\Bigl(
          z=\IntClass{a}{b}\land{}\\[-2pt]
        &\qquad u_1=v_R(a,b)\Bigr)
      \end{aligned}}{%
      \rAEb{\FormulaRefAuto{x\in\{u\in A\mid P(u)\}
        \eqvdash x\in A\land P(x)}{\rAEb{12}}}}
    \proofstepwidestar[13]{%
      \begin{aligned}[t]
        &\exists a,b\in\mathbb N\,\Bigl(
          z=\IntClass{a}{b}\land{}\\[-2pt]
        &\qquad u_2=v_R(a,b)\Bigr)
      \end{aligned}}{%
      \rAEb{\FormulaRefAuto{x\in\{u\in A\mid P(u)\}
        \eqvdash x\in A\land P(x)}{\rAEb{13}}}}
    \proofstepwidestar[12]{%
      \begin{aligned}[t]
        &u_1\in R\land
          \exists a,b\in\mathbb N\,\Bigl(
          z=\IntClass{a}{b}\land{}\\[-2pt]
        &\qquad u_1=v_R(a,b)\Bigr)
      \end{aligned}}{\rAI{\rAEa{12},14}}
    \proofstepwidestar[13]{%
      \begin{aligned}[t]
        &u_2\in R\land
          \exists a,b\in\mathbb N\,\Bigl(
          z=\IntClass{a}{b}\land{}\\[-2pt]
        &\qquad u_2=v_R(a,b)\Bigr)
      \end{aligned}}{\rAI{\rAEa{13},15}}
    \proofstepwidestar[1,3,12,13]{u_1=u_2}{%
      \FormulaRefAuto{\exists! x\,P(x)\dsep P(a)\dsep P(b)
        \vdash a=b}{4,16,17}}
    \proofstepwidestar[1,3]{%
      \exists! u\in R\,(z,u)\in G_R}{\UEI{11,12,13,18}}
    \proofstepwidestar[1]{%
      \forall z\in\mathbb Z\,\exists! u\in R\,(z,u)\in G_R}{%
      \rUI{\rRI{3,19}}}
    \proofstepwidestar[1]{G_R\colon\mathbb Z\to R}{%
      \FormulaRefAuto{UniqueValuedGraphFunction}[thm]{2,20}}
    \proofstepwidestar[22]{a,b\in\mathbb N}{\rA}
    \proofstepwidestar[22]{\IntClass{a}{b}\in\mathbb Z}{%
      \FormulaRefAuto{IntegerClassInZ}[thm]{22}}
    \proofstepwidestar[1,22]{%
      v_R(a,b)\in R}{%
      \FormulaRefAuto{IntegerRingPairValueCarrier}[thm]{1,22}}
    \proofstepwidestar[1,22]{%
      \left(
        \IntClass{a}{b},
        v_R(a,b)
      \right)\in\mathbb Z\times R}{%
      \FormulaRefAuto{a\in A\dsep b\in B\vdash(a,b)\in A\times B}{23,24}}
    \proofstepwidestar[22]{%
      \begin{aligned}[t]
        &\exists c,d\in\mathbb N\,\Bigl(
          \IntClass{a}{b}=\IntClass{c}{d}\\[-2pt]
        &\qquad{}\land v_R(a,b)=v_R(c,d)
          \Bigr)
      \end{aligned}}{\rEI{\rAI{22,\rAI{\rII,\rII}}}}
    \proofstepwidestar[1,22]{%
      \left(
        \IntClass{a}{b},
        v_R(a,b)
      \right)\in G_R}{%
      \FormulaRefAuto{x\in A\dsep P(x)\vdash
        x\in\{x\in A\mid P(x)\}}{25,26}}
    \proofstepwidestar[1,22]{%
      G_R(\IntClass{a}{b})
      =
      v_R(a,b)}{%
      \FormulaRefAuto{F\colon A\to B\dsep (x,y)\in F
        \vdash F(x)=y}{21,27}}
    \proofstepwidestar[1]{%
      \begin{aligned}[t]
        &\forall a,b\in\mathbb N\,\\[-2pt]
        &G_R(\IntClass{a}{b})\\[-2pt]
        &\qquad=
        v_R(a,b)
      \end{aligned}}{\rUI{\rRI{22,28}}}
    \proofstepwidestar[1]{%
      \begin{aligned}[t]
        &G_R\colon\mathbb Z\to R\\[-2pt]
        &\land
        \forall a,b\in\mathbb N\,\\[-2pt]
        &G_R(\IntClass{a}{b})\\[-2pt]
        &\qquad=
        v_R(a,b)
      \end{aligned}}{\rAI{21,29}}
    \proofstepwidestar[1]{%
      \exists f\Bigl(
        \begin{aligned}[t]
          &f\colon\mathbb Z\to R\land{}\\[-2pt]
          &\forall a,b\in\mathbb N\,
            f(\IntClass{a}{b})=v_R(a,b)
        \end{aligned}
      \Bigr)}{\rEI{30}}
  \closeproofpartwideindR

  \proofpartwideindR[Eindeutigkeit der Klassenabbildung]{%
    \begin{aligned}[t]
      &\RingAlg{R}{\oplus}{\odot}{0_R}{1_R}\dsep{}\\[-2pt]
      &\Bigl(
        f\colon\mathbb Z\to R\land
        \forall a,b\in\mathbb N\,
        f(\IntClass{a}{b})
        =
        v_R(a,b)
      \Bigr)\dsep{}\\[-2pt]
      &\Bigl(
        g\colon\mathbb Z\to R\land
        \forall a,b\in\mathbb N\,
        g(\IntClass{a}{b})
        =
        v_R(a,b)
      \Bigr)
      \vdash f=g
    \end{aligned}
  }{%
    \begin{aligned}[t]
      &\RingAlg{R}{\oplus}{\odot}{0_R}{1_R}\dsep{}\\[-2pt]
      &\Bigl(
        f\colon\mathbb Z\to R\land
        \forall a,b\in\mathbb N\,
        f(\IntClass{a}{b})
        =
        v_R(a,b)
      \Bigr)\dsep{}\\[-2pt]
      &\Bigl(
        g\colon\mathbb Z\to R\land
        \forall a,b\in\mathbb N\,
        g(\IntClass{a}{b})
        =
        v_R(a,b)
      \Bigr)
      \vdash f=g
    \end{aligned}
  }[IntegerRingMapDescentUniqueness]
    \proofstepwidestar[1]{%
      \RingAlg{R}{\oplus}{\odot}{0_R}{1_R}}{\rA}
    \proofstepwidestar[2]{%
      \begin{aligned}[t]
        &f\colon\mathbb Z\to R\land{}\\[-2pt]
        &\forall a,b\in\mathbb N\,
          f(\IntClass{a}{b})=v_R(a,b)
      \end{aligned}}{\rA}
    \proofstepwidestar[3]{%
      \begin{aligned}[t]
        &g\colon\mathbb Z\to R\land{}\\[-2pt]
        &\forall a,b\in\mathbb N\,
          g(\IntClass{a}{b})=v_R(a,b)
      \end{aligned}}{\rA}
    \proofstepwidestar[2]{f\colon\mathbb Z\to R}{\rAEa{2}}
    \proofstepwidestar[2]{%
      \forall a,b\in\mathbb N\,
      f(\IntClass{a}{b})
      =
      v_R(a,b)}{\rAEb{2}}
    \proofstepwidestar[3]{g\colon\mathbb Z\to R}{\rAEa{3}}
    \proofstepwidestar[3]{%
      \forall a,b\in\mathbb N\,
      g(\IntClass{a}{b})
      =
      v_R(a,b)}{\rAEb{3}}
    \proofstepwidestar[8]{z\in\mathbb Z}{\rA}
    \proofstepwidestar[8]{%
      \exists a,b\in\mathbb N\,z=\IntClass{a}{b}}{%
      \FormulaRefAuto{IntegerPairRepresentative}[thm]{8}}
    \proofstepwidestar[10]{%
      a,b\in\mathbb N\land z=\IntClass{a}{b}}{\rA}
    \proofstepwidestar[10]{a,b\in\mathbb N}{\rAEa{10}}
    \proofstepwidestar[10]{z=\IntClass{a}{b}}{\rAEb{10}}
    \proofstepwidestar[10]{a\in\mathbb N}{\rAEa{11}}
    \proofstepwidestar[10]{b\in\mathbb N}{\rAEb{11}}
    \proofstepwidestar[2,10]{%
      \forall b\in\mathbb N\,
      f(\IntClass{a}{b})
      =
      v_R(a,b)}{%
      \rRE{13,\rUE{5}}}
    \proofstepwidestar[2,10]{%
      f(\IntClass{a}{b})
      =
      v_R(a,b)}{%
      \rRE{14,\rUE{15}}}
    \proofstepwidestar[3,10]{%
      \forall b\in\mathbb N\,
      g(\IntClass{a}{b})
      =
      v_R(a,b)}{%
      \rRE{13,\rUE{7}}}
    \proofstepwidestar[3,10]{%
      g(\IntClass{a}{b})
      =
      v_R(a,b)}{%
      \rRE{14,\rUE{17}}}
    \proofstepwidestar[2,10]{%
      f(z)=v_R(a,b)}{%
      \rIE{12,16}}
    \proofstepwidestar[3,10]{%
      g(z)=v_R(a,b)}{%
      \rIE{12,18}}
    \proofstepwidestar[2,3,10]{f(z)=g(z)}{%
      \FormulaRefAuto{a=b\dsep c=b\vdash a=c}{19,20}}
    \proofstepwidestar[2,3,8]{f(z)=g(z)}{\rEE{9,10,21}}
    \proofstepwidestar[2,3]{%
      \forall z\in\mathbb Z\,f(z)=g(z)}{\rUI{\rRI{8,22}}}
    \proofstepwidestar[2,3]{f=g}{%
      \FormulaRefAuto{FunctionExtensionality}[thm]{4,6,23}}
  \closeproofpartwideindR

  \proofpartwide{Abschluss}
    \proofstepwidestar[1]{%
      \RingAlg{R}{\oplus}{\odot}{0_R}{1_R}}{\rA}
    \proofstepwidestar[1]{%
      \exists f\Bigl(
        \begin{aligned}[t]
          &f\colon\mathbb Z\to R\land{}\\[-2pt]
          &\forall a,b\in\mathbb N\,
            f(\IntClass{a}{b})=v_R(a,b)
        \end{aligned}
      \Bigr)}{%
      \FormulaRefAuto{IntegerRingMapDescentExistence}[thm]{1}}
    \proofstepwidestar[3]{%
      \begin{aligned}[t]
        &f\colon\mathbb Z\to R\land{}\\[-2pt]
        &\forall a,b\in\mathbb N\,
          f(\IntClass{a}{b})=v_R(a,b)
      \end{aligned}}{\rA}
    \proofstepwidestar[4]{%
      \begin{aligned}[t]
        &g\colon\mathbb Z\to R\land{}\\[-2pt]
        &\forall a,b\in\mathbb N\,
          g(\IntClass{a}{b})=v_R(a,b)
      \end{aligned}}{\rA}
    \proofstepwidestar[1,3,4]{f=g}{%
      \FormulaRefAuto{IntegerRingMapDescentUniqueness}[thm]{1,3,4}}
    \proofstepwidestar[1]{%
      \exists! f\Bigl(
        \begin{aligned}[t]
          &f\colon\mathbb Z\to R\land{}\\[-2pt]
          &\forall a,b\in\mathbb N\,
            f(\IntClass{a}{b})=v_R(a,b)
        \end{aligned}
      \Bigr)}{\UEI{2,3,4,5}}
  \closeproofpartwide
\end{tabproofsplitwide}

\FormulaDefDeltaK[Kanonische Abbildung aus den ganzen Zahlen]{%
  \begin{aligned}[t]
    &\RingAlg{R}{\oplus}{\odot}{0_R}{1_R}\vdash{}\\[-2pt]
    &\IntegerRingMap{R,\oplus,\odot,0_R,1_R}
      \colon\mathbb Z\to R,\\[-2pt]
    &\forall a,b\in\mathbb N\,
      \Bigl(
        \IntegerRingMap{R,\oplus,\odot,0_R,1_R}(\IntClass{a}{b})
        \coloneqq
        \IntegerRingPairValue{R,\oplus,\odot,0_R,1_R}{a}{b}
      \Bigr)
  \end{aligned}
}{IntegerRingMapDef}{%
  \DeltaRow{Mengen}{R\dsep a\dsep b}
  \DeltaRow{Funktionen}{\oplus\dsep\odot}
  \DeltaRow{Repräsentanten}{%
    \FormulaRefAuto{IntegerPairRepresentative}[thm]}
  \DeltaRow{Repräsentantenunabhängigkeit}{%
    \FormulaRefAuto{IntegerRingMapRepresentativeIndependent}[thm]}
  \DeltaRow{Existenz und Eindeutigkeit}{%
    \FormulaRefAuto{IntegerRingMapQuotientDescent}[thm]}
  \DeltaRow{\textbf{Neues Symbol}}{}
  \DeltaPrem{Kanonische Abbildung}{%
    \IntegerRingMap{R,\oplus,\odot,0_R,1_R}}
}

\FormulaThmDeltaK[Trägerabbildung der kanonischen Ganzzahlabbildung]{%
  \RingAlg{R}{\oplus}{\odot}{0_R}{1_R}
  \vdash
  \IntegerRingMap{R,\oplus,\odot,0_R,1_R}
  \colon\mathbb Z\to R%
}{IntegerRingMapFunction}{%
  \DeltaRow{Mengen}{R}
  \DeltaRow{Funktionen}{\oplus\dsep\odot}
}
\begin{tabproof}
  \proofstep{1}{\RingAlg{R}{\oplus}{\odot}{0_R}{1_R}}{\rA}
  \proofstep{1}{%
    \IntegerRingMap{R,\oplus,\odot,0_R,1_R}
    \colon\mathbb Z\to R}{%
    \FormulaRefAuto{IntegerRingMapDef}[def]{1}}
\end{tabproof}
\par\smallskip

\FormulaThmDeltaK[Auswertung an einer Ganzzahlklasse]{%
  \begin{aligned}[t]
    &\RingAlg{R}{\oplus}{\odot}{0_R}{1_R}\dsep a,b\in\mathbb N\\
    &\vdash
      \IntegerRingMap{R,\oplus,\odot,0_R,1_R}(\IntClass{a}{b})
      =
      \NaturalSemiringMap{R,\oplus,\odot,0_R,1_R}(a)
      \oplus
      \bigl(-\NaturalSemiringMap{R,\oplus,\odot,0_R,1_R}(b)\bigr)
  \end{aligned}
}{IntegerRingMapClassValue}{%
  \DeltaRow{Mengen}{R\dsep a\dsep b}
  \DeltaRow{Funktionen}{\oplus\dsep\odot}
}
\begin{tabproofwide}
  \proofstepwidestar[1]{\RingAlg{R}{\oplus}{\odot}{0_R}{1_R}}{\rA}
  \proofstepwidestar[2]{a,b\in\mathbb N}{\rA}
  \proofstepwidestar[1,2]{%
    \begin{aligned}[t]
      &\IntegerRingMap{R,\oplus,\odot,0_R,1_R}(\IntClass{a}{b})\\[-2pt]
      &\quad=
        \NaturalSemiringMap{R,\oplus,\odot,0_R,1_R}(a)
        \oplus
        \bigl(-\NaturalSemiringMap{R,\oplus,\odot,0_R,1_R}(b)\bigr)
    \end{aligned}}{%
    \parbox[t]{.25\textwidth}{\raggedright
      \FormulaRefAuto{IntegerRingMapDef}[def]\\
      \FormulaRefAuto{IntegerRingPairValueDef}[def]}}
\end{tabproofwide}

\subsection{Erhaltung der Ringoperationen}

\FormulaThmDeltaK[Die kanonische Ganzzahlabbildung erhält die Addition]{%
  \begin{aligned}[t]
    &\RingAlg{R}{\oplus}{\odot}{0_R}{1_R}\dsep z,w\in\mathbb Z\\
    &\vdash
      \IntegerRingMap{R,\oplus,\odot,0_R,1_R}(z+w)
      =
      \IntegerRingMap{R,\oplus,\odot,0_R,1_R}(z)
      \oplus
      \IntegerRingMap{R,\oplus,\odot,0_R,1_R}(w)
  \end{aligned}
}{IntegerRingMapAddition}{%
  \DeltaRow{Mengen}{R\dsep z\dsep w\dsep a\dsep b\dsep c\dsep d}
  \DeltaRow{Funktionen}{\oplus\dsep\odot}
}

\begin{tabproofwide}
  \proofstepwidestar[1]{%
    \RingAlg{R}{\oplus}{\odot}{0_R}{1_R}}{\rA}
  \proofstepwidestar[2]{z\in\mathbb Z}{\rA}
  \proofstepwidestar[3]{w\in\mathbb Z}{\rA}
  \proofstepwidestar[2]{%
    \exists a,b\in\mathbb N\,z=\IntClass{a}{b}}{%
    \FormulaRefAuto{IntegerPairRepresentative}[thm]{2}}
  \proofstepwidestar[3]{%
    \exists c,d\in\mathbb N\,w=\IntClass{c}{d}}{%
    \FormulaRefAuto{IntegerPairRepresentative}[thm]{3}}
  \proofstepwidestar[6]{%
    a,b\in\mathbb N\land z=\IntClass{a}{b}}{\rA}
  \proofstepwidestar[7]{%
    c,d\in\mathbb N\land w=\IntClass{c}{d}}{\rA}
  \proofstepwidestar[6,7]{%
    z+w=\IntClass{a+c}{b+d}}{%
    \FormulaRefAuto{IntegerAdditionDef}[def]{6,7}}
  \proofstepwidestar[1,6,7]{%
    \begin{aligned}[t]
      &\IntegerRingMap{R,\oplus,\odot,0_R,1_R}(z+w)\\[-2pt]
      &\quad=
      \IntegerRingPairValue{R,\oplus,\odot,0_R,1_R}{a+c}{b+d}
    \end{aligned}}{%
    \parbox[t]{.25\textwidth}{\raggedright
      \FormulaRefAuto{IntegerAdditionDef}[def]\\
      \FormulaRefAuto{IntegerRingMapClassValue}[thm]\\
      \FormulaRefAuto{IntegerRingPairValueDef}[def]\\
      \FormulaRefAuto{PeanoAddClosure}[thm]}}
  \proofstepwidestar[1,6,7]{%
    \begin{aligned}[t]
      &\IntegerRingPairValue{R,\oplus,\odot,0_R,1_R}{a+c}{b+d}\\[-2pt]
      &\quad=
      \IntegerRingPairValue{R,\oplus,\odot,0_R,1_R}{a}{b}
      \oplus
      \IntegerRingPairValue{R,\oplus,\odot,0_R,1_R}{c}{d}
    \end{aligned}}{%
    \parbox[t]{.25\textwidth}{\raggedright
      \FormulaRefAuto{IntegerRingPairValueDef}[def]\\
      \FormulaRefAuto{RingUnderlyingSemiring}[thm]\\
      \FormulaRefAuto{NaturalSemiringMapAddition}[thm]\\
      \FormulaRefAuto{GroupProductInverse}[thm]\\
      \FormulaRefAuto{RingAdditiveAbelianGroupAxiom}[ax]\\
      \FormulaRefAuto{SemigroupAssociativityAxiom}[ax]\\
      \FormulaRefAuto{AbelianGroupCommutativityAxiom}[ax]}}
  \proofstepwidestar[1,6]{%
    \IntegerRingMap{R,\oplus,\odot,0_R,1_R}(z)
    =
    \IntegerRingPairValue{R,\oplus,\odot,0_R,1_R}{a}{b}}{%
    \parbox[t]{.25\textwidth}{\raggedright
      \FormulaRefAuto{IntegerRingMapClassValue}[thm]\\
      \FormulaRefAuto{IntegerRingPairValueDef}[def]}}
  \proofstepwidestar[1,7]{%
    \IntegerRingMap{R,\oplus,\odot,0_R,1_R}(w)
    =
    \IntegerRingPairValue{R,\oplus,\odot,0_R,1_R}{c}{d}}{%
    \parbox[t]{.25\textwidth}{\raggedright
      \FormulaRefAuto{IntegerRingMapClassValue}[thm]\\
      \FormulaRefAuto{IntegerRingPairValueDef}[def]}}
  \proofstepwidestar[1,6,7]{%
    \begin{aligned}[t]
      &\IntegerRingPairValue{R,\oplus,\odot,0_R,1_R}{a}{b}
      \oplus
      \IntegerRingPairValue{R,\oplus,\odot,0_R,1_R}{c}{d}\\[-2pt]
      &\quad=
      \IntegerRingMap{R,\oplus,\odot,0_R,1_R}(z)
      \oplus
      \IntegerRingMap{R,\oplus,\odot,0_R,1_R}(w)
    \end{aligned}}{\rIE{11,12,\rII}}
  \proofstepwidestar[1,6,7]{%
    \begin{aligned}[t]
      &\IntegerRingMap{R,\oplus,\odot,0_R,1_R}(z+w)\\[-2pt]
      &\quad=
      \IntegerRingMap{R,\oplus,\odot,0_R,1_R}(z)
      \oplus
      \IntegerRingMap{R,\oplus,\odot,0_R,1_R}(w)
    \end{aligned}}{\rChain{9,10,13}}
  \proofstepwidestar[1,3,6]{%
    \begin{aligned}[t]
      &\IntegerRingMap{R,\oplus,\odot,0_R,1_R}(z+w)\\[-2pt]
      &\quad=
      \IntegerRingMap{R,\oplus,\odot,0_R,1_R}(z)
      \oplus
      \IntegerRingMap{R,\oplus,\odot,0_R,1_R}(w)
    \end{aligned}}{\rEE{5,7,14}}
  \proofstepwidestar[1,2,3]{%
    \begin{aligned}[t]
      &\IntegerRingMap{R,\oplus,\odot,0_R,1_R}(z+w)\\[-2pt]
      &\quad=
      \IntegerRingMap{R,\oplus,\odot,0_R,1_R}(z)
      \oplus
      \IntegerRingMap{R,\oplus,\odot,0_R,1_R}(w)
    \end{aligned}}{\rEE{4,6,15}}
\end{tabproofwide}

\FormulaThmDeltaK[Die kanonische Ganzzahlabbildung erhält die Multiplikation]{%
  \begin{aligned}[t]
    &\RingAlg{R}{\oplus}{\odot}{0_R}{1_R}\dsep z,w\in\mathbb Z\\
    &\vdash
      \IntegerRingMap{R,\oplus,\odot,0_R,1_R}(z\cdot w)
      =
      \IntegerRingMap{R,\oplus,\odot,0_R,1_R}(z)
      \odot
      \IntegerRingMap{R,\oplus,\odot,0_R,1_R}(w)
  \end{aligned}
}{IntegerRingMapMultiplication}{%
  \DeltaRow{Mengen}{R\dsep z\dsep w\dsep a\dsep b\dsep c\dsep d}
  \DeltaRow{Funktionen}{\oplus\dsep\odot}
}

\begin{tabproofwide}
  \proofstepwidestar[1]{%
    \RingAlg{R}{\oplus}{\odot}{0_R}{1_R}}{\rA}
  \proofstepwidestar[2]{z\in\mathbb Z}{\rA}
  \proofstepwidestar[3]{w\in\mathbb Z}{\rA}
  \proofstepwidestar[2]{%
    \exists a,b\in\mathbb N\,z=\IntClass{a}{b}}{%
    \FormulaRefAuto{IntegerPairRepresentative}[thm]{2}}
  \proofstepwidestar[3]{%
    \exists c,d\in\mathbb N\,w=\IntClass{c}{d}}{%
    \FormulaRefAuto{IntegerPairRepresentative}[thm]{3}}
  \proofstepwidestar[6]{%
    a,b\in\mathbb N\land z=\IntClass{a}{b}}{\rA}
  \proofstepwidestar[7]{%
    c,d\in\mathbb N\land w=\IntClass{c}{d}}{\rA}
  \proofstepwidestar[6,7]{%
    z\cdot w=\IntClass{ac+bd}{ad+bc}}{%
    \FormulaRefAuto{IntegerMultiplicationDef}[def]{6,7}}
  \proofstepwidestar[1,6,7]{%
    \begin{aligned}[t]
      &\IntegerRingMap{R,\oplus,\odot,0_R,1_R}(z\cdot w)\\[-2pt]
      &\quad=
      \IntegerRingPairValue{R,\oplus,\odot,0_R,1_R}{ac+bd}{ad+bc}
    \end{aligned}}{%
    \parbox[t]{.25\textwidth}{\raggedright
      \FormulaRefAuto{IntegerMultiplicationDef}[def]\\
      \FormulaRefAuto{IntegerRingMapClassValue}[thm]\\
      \FormulaRefAuto{IntegerRingPairValueDef}[def]\\
      \FormulaRefAuto{PeanoAddClosure}[thm]\\
      \FormulaRefAuto{PeanoMulClosure}[thm]}}
  \proofstepwidestar[1,6,7]{%
    \begin{aligned}[t]
      &\IntegerRingPairValue{R,\oplus,\odot,0_R,1_R}{ac+bd}{ad+bc}\\[-2pt]
      &\quad=
      \IntegerRingPairValue{R,\oplus,\odot,0_R,1_R}{a}{b}
      \odot
      \IntegerRingPairValue{R,\oplus,\odot,0_R,1_R}{c}{d}
    \end{aligned}}{%
    \parbox[t]{.25\textwidth}{\raggedright
      \FormulaRefAuto{IntegerRingPairValueDef}[def]\\
      \FormulaRefAuto{RingUnderlyingSemiring}[thm]\\
      \FormulaRefAuto{NaturalSemiringMapAddition}[thm]\\
      \FormulaRefAuto{NaturalSemiringMapMultiplication}[thm]\\
      \FormulaRefAuto{RingLeftDistributivityAxiom}[ax]\\
      \FormulaRefAuto{RingRightDistributivityAxiom}[ax]\\
      \FormulaRefAuto{GroupProductInverse}[thm]\\
      \FormulaRefAuto{RingLeftNegativeProduct}[thm]\\
      \FormulaRefAuto{RingRightNegativeProduct}[thm]\\
      \FormulaRefAuto{RingNegativeTimesNegative}[thm]\\
      \FormulaRefAuto{RingAdditiveAbelianGroupAxiom}[ax]\\
      \FormulaRefAuto{SemigroupAssociativityAxiom}[ax]\\
      \FormulaRefAuto{AbelianGroupCommutativityAxiom}[ax]}}
  \proofstepwidestar[1,6]{%
    \IntegerRingMap{R,\oplus,\odot,0_R,1_R}(z)
    =
    \IntegerRingPairValue{R,\oplus,\odot,0_R,1_R}{a}{b}}{%
    \parbox[t]{.25\textwidth}{\raggedright
      \FormulaRefAuto{IntegerRingMapClassValue}[thm]\\
      \FormulaRefAuto{IntegerRingPairValueDef}[def]}}
  \proofstepwidestar[1,7]{%
    \IntegerRingMap{R,\oplus,\odot,0_R,1_R}(w)
    =
    \IntegerRingPairValue{R,\oplus,\odot,0_R,1_R}{c}{d}}{%
    \parbox[t]{.25\textwidth}{\raggedright
      \FormulaRefAuto{IntegerRingMapClassValue}[thm]\\
      \FormulaRefAuto{IntegerRingPairValueDef}[def]}}
  \proofstepwidestar[1,6,7]{%
    \begin{aligned}[t]
      &\IntegerRingPairValue{R,\oplus,\odot,0_R,1_R}{a}{b}
      \odot
      \IntegerRingPairValue{R,\oplus,\odot,0_R,1_R}{c}{d}\\[-2pt]
      &\quad=
      \IntegerRingMap{R,\oplus,\odot,0_R,1_R}(z)
      \odot
      \IntegerRingMap{R,\oplus,\odot,0_R,1_R}(w)
    \end{aligned}}{\rIE{11,12,\rII}}
  \proofstepwidestar[1,6,7]{%
    \begin{aligned}[t]
      &\IntegerRingMap{R,\oplus,\odot,0_R,1_R}(z\cdot w)\\[-2pt]
      &\quad=
      \IntegerRingMap{R,\oplus,\odot,0_R,1_R}(z)
      \odot
      \IntegerRingMap{R,\oplus,\odot,0_R,1_R}(w)
    \end{aligned}}{\rChain{9,10,13}}
  \proofstepwidestar[1,3,6]{%
    \begin{aligned}[t]
      &\IntegerRingMap{R,\oplus,\odot,0_R,1_R}(z\cdot w)\\[-2pt]
      &\quad=
      \IntegerRingMap{R,\oplus,\odot,0_R,1_R}(z)
      \odot
      \IntegerRingMap{R,\oplus,\odot,0_R,1_R}(w)
    \end{aligned}}{\rEE{5,7,14}}
  \proofstepwidestar[1,2,3]{%
    \begin{aligned}[t]
      &\IntegerRingMap{R,\oplus,\odot,0_R,1_R}(z\cdot w)\\[-2pt]
      &\quad=
      \IntegerRingMap{R,\oplus,\odot,0_R,1_R}(z)
      \odot
      \IntegerRingMap{R,\oplus,\odot,0_R,1_R}(w)
    \end{aligned}}{\rEE{4,6,15}}
\end{tabproofwide}

\FormulaThmDeltaK[Null- und Einserhaltung der kanonischen Ganzzahlabbildung]{%
  \RingAlg{R}{\oplus}{\odot}{0_R}{1_R}
  \vdash
  \IntegerRingMap{R,\oplus,\odot,0_R,1_R}(\IntZero)=0_R
  \land
  \IntegerRingMap{R,\oplus,\odot,0_R,1_R}(\IntOne)=1_R%
}{IntegerRingMapConstants}{%
  \DeltaRow{Mengen}{R\dsep\IntZero\dsep\IntOne}
  \DeltaRow{Funktionen}{\oplus\dsep\odot}
}

\begin{tabproofwide}
  \proofstepwidestar[1]{%
    \RingAlg{R}{\oplus}{\odot}{0_R}{1_R}}{\rA}
  \proofstepwidestar[1]{%
    \IntegerRingMap{R,\oplus,\odot,0_R,1_R}(\IntZero)=0_R}{%
    \parbox[t]{.25\textwidth}{\raggedright
      \FormulaRefAuto{IntegerZeroClass}[thm]\\
      \FormulaRefAuto{IntegerRingMapClassValue}[thm]\\
      \FormulaRefAuto{NaturalSemiringMapFunction}[thm]\\
      \FormulaRefAuto{NaturalSemiringMapZero}[thm]\\
      \FormulaRefAuto{RingNegativeRight}[thm]}}
  \proofstepwidestar[1]{%
    \IntegerRingMap{R,\oplus,\odot,0_R,1_R}(\IntOne)=1_R}{%
    \parbox[t]{.25\textwidth}{\raggedright
      \FormulaRefAuto{IntegerOneDef}[def]\\
      \FormulaRefAuto{NaturalIntegerEmbedding}[def]\\
      \FormulaRefAuto{PeanoOneInN}[thm]\\
      \FormulaRefAuto{IntegerRingMapClassValue}[thm]\\
      \FormulaRefAuto{NaturalSemiringMapFunction}[thm]\\
      \FormulaRefAuto{NaturalSemiringMapZero}[thm]\\
      \FormulaRefAuto{NaturalSemiringMapOne}[thm]\\
      \FormulaRefAuto{RingNegativeRight}[thm]\\
      \FormulaRefAuto{GroupLeftIdentity}[thm]\\
      \FormulaRefAuto{GroupRightIdentity}[thm]}}
  \proofstepwidestar[1]{%
    \IntegerRingMap{R,\oplus,\odot,0_R,1_R}(\IntZero)=0_R
    \land
    \IntegerRingMap{R,\oplus,\odot,0_R,1_R}(\IntOne)=1_R}{%
    \rAI{2,3}}
\end{tabproofwide}

\FormulaThmDeltaK[Die kanonische Ganzzahlabbildung ist ein Ringhomomorphismus]{%
  \RingAlg{R}{\oplus}{\odot}{0_R}{1_R}
  \vdash
  \RingHom{\IntegerRingMap{R,\oplus,\odot,0_R,1_R}}
    {\mathbb Z,+,\cdot,\IntZero,\IntOne}
    {R,\oplus,\odot,0_R,1_R}%
}{IntegerRingCanonicalHomomorphism}{%
  \DeltaRow{Mengen}{R\dsep z\dsep w}
  \DeltaRow{Funktionen}{\oplus\dsep\odot}
}

Für die beiden folgenden Beweise setzen wir lokal
\[
  \mathfrak Z
    \coloneqq(\mathbb Z,+,\cdot,\IntZero,\IntOne),
  \qquad
  \mathfrak R
    \coloneqq(R,\oplus,\odot,0_R,1_R),
  \qquad
  \Phi_R
    \coloneqq\IntegerRingMap{R,\oplus,\odot,0_R,1_R}.
\]

\begin{tabproofwide}
  \proofstepwidestar[1]{%
    \RingAlg{R}{\oplus}{\odot}{0_R}{1_R}}{\rA}
  \proofstepwidestar[]{%
    \RingAlg{\mathbb Z}{+}{\cdot}{\IntZero}{\IntOne}}{%
    \FormulaRefAuto{IntegerUnitalRing}[thm]}
  \proofstepwidestar[1]{%
    \Phi_R\colon\mathbb Z\to R}{%
    \FormulaRefAuto{IntegerRingMapFunction}[thm]{1}}
  \proofstepwidestar[]{%
    \GroupAlg{\mathbb Z}{+}{\IntZero}}{%
    \FormulaRefAuto{IntegerAdditiveGroup}[thm]}
  \proofstepwidestar[1]{\GroupAlg{R}{\oplus}{0_R}}{%
    \FormulaRefAuto{RingAdditiveGroup}[thm]{1}}
  \proofstepwidestar[]{\SemigroupAlg{\mathbb Z}{+}}{%
    \FormulaRefAuto{GroupBaseSemigroup}[thm]{4}}
  \proofstepwidestar[1]{\SemigroupAlg{R}{\oplus}}{%
    \FormulaRefAuto{GroupBaseSemigroup}[thm]{5}}
  \proofstepwidestar[8]{z\in\mathbb Z}{\rA}
  \proofstepwidestar[9]{w\in\mathbb Z}{\rA}
  \proofstepwidestar[1,8,9]{%
    \begin{aligned}[t]
      &\Phi_R(z+w)\\[-2pt]
      &\quad=
      \Phi_R(z)
      \oplus
      \Phi_R(w)
    \end{aligned}}{%
    \FormulaRefAuto{IntegerRingMapAddition}[thm]{1,8,9}}
  \proofstepwidestar[1]{%
    \SemigroupHom{\Phi_R}
      {\mathbb Z,+}{R,\oplus}}{%
    \FormulaRefAuto{SemigroupHomDef}[def]{6,7,3,10}}
  \proofstepwidestar[1]{%
    \GroupHom{\Phi_R}
      {\mathbb Z,+,\IntZero}{R,\oplus,0_R}}{%
    \FormulaRefAuto{GroupHomDef}[def]{4,5,11}}
  \proofstepwidestar[]{%
    \MonoidAlg{\mathbb Z}{\cdot}{\IntOne}}{%
    \FormulaRefAuto{IntegerMultiplicativeMonoid}[thm]}
  \proofstepwidestar[1]{\MonoidAlg{R}{\odot}{1_R}}{%
    \FormulaRefAuto{RingMultiplicativeMonoidAxiom}[ax]{1}}
  \proofstepwidestar[]{\SemigroupAlg{\mathbb Z}{\cdot}}{%
    \FormulaRefAuto{MonoidBaseSemigroupAxiom}[ax]{13}}
  \proofstepwidestar[1]{\SemigroupAlg{R}{\odot}}{%
    \FormulaRefAuto{MonoidBaseSemigroupAxiom}[ax]{14}}
  \proofstepwidestar[1,8,9]{%
    \begin{aligned}[t]
      &\Phi_R(z\cdot w)\\[-2pt]
      &\quad=
      \Phi_R(z)
      \odot
      \Phi_R(w)
    \end{aligned}}{%
    \FormulaRefAuto{IntegerRingMapMultiplication}[thm]{1,8,9}}
  \proofstepwidestar[1]{%
    \SemigroupHom{\Phi_R}
      {\mathbb Z,\cdot}{R,\odot}}{%
    \FormulaRefAuto{SemigroupHomDef}[def]{15,16,3,17}}
  \proofstepwidestar[1]{%
    \Phi_R(\IntZero)=0_R
    \land
    \Phi_R(\IntOne)=1_R}{%
    \FormulaRefAuto{IntegerRingMapConstants}[thm]{1}}
  \proofstepwidestar[1]{%
    \Phi_R(\IntOne)=1_R}{%
    \rAEb{19}}
  \proofstepwidestar[1]{%
    \MonoidHom{\Phi_R}
      {\mathbb Z,\cdot,\IntOne}{R,\odot,1_R}}{%
    \FormulaRefAuto{MonoidHomDef}[def]{13,14,18,20}}
  \proofstepwidestar[1]{%
    \RingHom{\Phi_R}{\mathfrak Z}{\mathfrak R}}{%
    \FormulaRefAuto{RingHomDef}[def]{2,1,12,21}}
\end{tabproofwide}

\FormulaThmDeltaK[Universelle Eigenschaft der ganzen Zahlen]{%
  \begin{aligned}[t]
    &\RingAlg{R}{\oplus}{\odot}{0_R}{1_R}\vdash{}\\[-2pt]
    &\exists! f\,
      \RingHom{f}
        {\mathbb Z,+,\cdot,\IntZero,\IntOne}
        {R,\oplus,\odot,0_R,1_R}
  \end{aligned}
}{IntegerRingUniversalHomomorphism}{%
  \DeltaRow{Mengen}{R\dsep z\dsep m\dsep n}
  \DeltaRow{Funktionen}{f\dsep\oplus\dsep\odot}
}

\begin{tabproofsplitwide}
  \proofpartwideindR[Eindeutigkeit auf der natürlichen Einbettung]{%
    \begin{aligned}[t]
      &\operatorname{Ring}(\mathfrak R)\dsep\\[-2pt]
      &\RingHom{f}{\mathfrak Z}{\mathfrak R}\dsep\\[-2pt]
      &n\in\mathbb N\vdash
        f(\IntEmbed(n))={}\\[-2pt]
      &\quad
        \NaturalSemiringMap{R,\oplus,\odot,0_R,1_R}(n)
    \end{aligned}
  }{%
    \begin{aligned}[t]
      &\operatorname{Ring}(\mathfrak R)\dsep\\[-2pt]
      &\RingHom{f}{\mathfrak Z}{\mathfrak R}\dsep\\[-2pt]
      &n\in\mathbb N\vdash
        f(\IntEmbed(n))={}\\[-2pt]
      &\quad
        \NaturalSemiringMap{R,\oplus,\odot,0_R,1_R}(n)
    \end{aligned}
  }[IntegerRingHomEmbeddingValue]
    \proofstepwidestar[1]{%
      \operatorname{Ring}(\mathfrak R)}{\rA}
    \proofstepwidestar[2]{%
      \RingHom{f}{\mathfrak Z}{\mathfrak R}}{\rA}
    \proofstepwidestar[2]{f(\IntZero)=0_R}{%
      \FormulaRefAuto{RingHomZeroPreservation}[thm]{2}}
    \proofstepwidestar[]{\IntZero=\IntEmbed(0)}{%
      \FormulaRefAuto{IntegerZeroDef}[def]}
    \proofstepwidestar[2]{f(\IntEmbed(0))=0_R}{\rIE{4,3}}
    \proofstepwidestar[1]{%
      \NaturalSemiringMap{R,\oplus,\odot,0_R,1_R}(0)=0_R}{%
      \FormulaRefAuto{NaturalSemiringMapZero}[thm]{%
        \FormulaRefAuto{RingUnderlyingSemiring}[thm]{1}}}
    \proofstepwidestar[1,2]{%
      \begin{aligned}[t]
        &f(\IntEmbed(0))={}\\[-2pt]
        &\quad\NaturalSemiringMap{R,\oplus,\odot,0_R,1_R}(0)
      \end{aligned}}{%
      \rIE{5,\FormulaRefAuto{a=b\vdash b=a}{6}}}
    \proofstepwidestar[8]{m\in\mathbb N}{\rA}
    \proofstepwidestar[9]{%
      \begin{aligned}[t]
        &f(\IntEmbed(m))={}\\[-2pt]
        &\quad\NaturalSemiringMap{R,\oplus,\odot,0_R,1_R}(m)
      \end{aligned}}{\rA}
    \proofstepwidestar[8]{%
      \begin{aligned}[t]
        &f(\IntEmbed(\Succ(m)))={}\\[-2pt]
        &\quad f(\IntEmbed(m+1))
      \end{aligned}}{%
      \BandXXIXProofRefStack{%
        \FormulaRefAuto{a=b\vdash b=a}[thm]\\
        \bigl(\FormulaRefAuto{PeanoAddOneSucc}[thm](8)\bigr)}}
    \proofstepwidestar[8]{%
      \IntEmbed(m+1)=\IntEmbed(m)+\IntEmbed(1)}{%
      \BandXXIXProofRefStack{%
        \FormulaRefAuto{NaturalIntegerEmbeddingAdditive}[thm]\\
        \bigl(8,\FormulaRefAuto{PeanoOneInN}[thm]\bigr)}}
    \proofstepwidestar[8]{%
      \begin{aligned}[t]
        &f(\IntEmbed(m+1))={}\\[-2pt]
        &\quad f\bigl(\IntEmbed(m)+\IntEmbed(1)\bigr)
      \end{aligned}}{\rIE{11,\rII}}
    \proofstepwidestar[2,8]{%
      \begin{aligned}[t]
        &f\bigl(\IntEmbed(m)+\IntEmbed(1)\bigr)={}\\[-2pt]
        &\quad f(\IntEmbed(m))\oplus f(\IntEmbed(1))
      \end{aligned}}{%
      \BandXXIXProofRefStack{%
        \FormulaRefAuto{RingHomAdditionAxiom}[ax]\\
        \left(
          \begin{gathered}
            2,\,
            \FormulaRefAuto{F\colon A\to B\dsep x\in A\vdash F(x)\in B}{%
              \FormulaRefAuto{NaturalIntegerEmbeddingFunction}[thm],8},\\[-2pt]
            \FormulaRefAuto{F\colon A\to B\dsep x\in A\vdash F(x)\in B}{%
              \FormulaRefAuto{NaturalIntegerEmbeddingFunction}[thm],
              \FormulaRefAuto{PeanoOneInN}[thm]}
          \end{gathered}
        \right)}}
    \proofstepwidestar[2]{f(\IntEmbed(1))=1_R}{%
      \rIE{%
        \FormulaRefAuto{IntegerOneDef}[def],
        \FormulaRefAuto{RingHomOneAxiom}[ax]{2}}}
    \proofstepwidestar[1,2,9]{%
      \begin{aligned}[t]
        &f(\IntEmbed(m))\oplus f(\IntEmbed(1))={}\\[-2pt]
        &\quad\NaturalSemiringMap{R,\oplus,\odot,0_R,1_R}(m)\oplus1_R
      \end{aligned}}{%
      \rIE{9,14}}
    \proofstepwidestar[1,8]{%
      \begin{aligned}[t]
        &\NaturalSemiringMap{R,\oplus,\odot,0_R,1_R}(m)\oplus1_R={}\\[-2pt]
        &\quad\NaturalSemiringMap{R,\oplus,\odot,0_R,1_R}(\Succ(m))
      \end{aligned}}{%
      \BandXXIXProofRefStack{%
        \FormulaRefAuto{a=b\vdash b=a}[thm]\\
        \left(
          \FormulaRefAuto{NaturalSemiringMapSuccessor}[thm]
          \bigl(\FormulaRefAuto{RingUnderlyingSemiring}[thm](1),8\bigr)
        \right)}}
    \proofstepwidestar[1,2,8,9]{%
      \begin{aligned}[t]
        &f(\IntEmbed(\Succ(m)))={}\\[-2pt]
        &\quad\NaturalSemiringMap{R,\oplus,\odot,0_R,1_R}(\Succ(m))
      \end{aligned}}{%
      \rChain{10,12,13,15,16}}
    \proofstepwidestar[18]{n\in\mathbb N}{\rA}
    \proofstepwidestar[1,2,18]{%
      \begin{aligned}[t]
        &f(\IntEmbed(n))={}\\[-2pt]
        &\quad\NaturalSemiringMap{R,\oplus,\odot,0_R,1_R}(n)
      \end{aligned}}{%
      \FormulaRefAuto{PeanoInductionRule}[ax]{7,17,18}}
  \closeproofpartwideindR

  \proofpartwideindR[Punktweise Eindeutigkeit auf den ganzen Zahlen]{%
    \begin{aligned}[t]
      &\operatorname{Ring}(\mathfrak R)\dsep\\[-2pt]
      &\RingHom{f}{\mathfrak Z}{\mathfrak R}\dsep\\[-2pt]
      &z\in\mathbb Z\vdash f(z)=\Phi_R(z)
    \end{aligned}
  }{%
    \begin{aligned}[t]
      &\operatorname{Ring}(\mathfrak R)\dsep\\[-2pt]
      &\RingHom{f}{\mathfrak Z}{\mathfrak R}\dsep\\[-2pt]
      &z\in\mathbb Z\vdash f(z)=\Phi_R(z)
    \end{aligned}
  }[IntegerRingHomPointwiseUnique]
    \proofstepwidestar[1]{%
      \operatorname{Ring}(\mathfrak R)}{\rA}
    \proofstepwidestar[2]{%
      \RingHom{f}{\mathfrak Z}{\mathfrak R}}{\rA}
    \proofstepwidestar[3]{z\in\mathbb Z}{\rA}
    \proofstepwidestar[1]{%
      \RingHom{\Phi_R}{\mathfrak Z}{\mathfrak R}}{%
      \FormulaRefAuto{IntegerRingCanonicalHomomorphism}[thm]{1}}
    \proofstepwidestar[3]{%
      \exists n\in\mathbb N\,
      \bigl(z=\IntEmbed(n)\lor z=-\IntEmbed(n)\bigr)}{%
      \FormulaRefAuto{IntegerSignedNormalForm}[thm]{3}}
    \proofstepwidestar[6]{%
      n\in\mathbb N\land
      \bigl(z=\IntEmbed(n)\lor z=-\IntEmbed(n)\bigr)}{\rA}
    \proofstepwidestar[6]{n\in\mathbb N}{\rAEa{6}}
    \proofstepwidestar[6]{%
      z=\IntEmbed(n)\lor z=-\IntEmbed(n)}{\rAEb{6}}
    \proofstepwidestar[1,2,6]{%
      \begin{aligned}[t]
        &f(\IntEmbed(n))={}\\[-2pt]
        &\quad\NaturalSemiringMap{R,\oplus,\odot,0_R,1_R}(n)
      \end{aligned}}{%
      \FormulaRefAuto{IntegerRingHomEmbeddingValue}[thm]{1,2,7}}
    \proofstepwidestar[1,6]{%
      \begin{aligned}[t]
        &\Phi_R(\IntEmbed(n))={}\\[-2pt]
        &\quad\NaturalSemiringMap{R,\oplus,\odot,0_R,1_R}(n)
      \end{aligned}}{%
      \FormulaRefAuto{IntegerRingHomEmbeddingValue}[thm]{1,4,7}}
    \proofstepwidestar[1,2,6]{%
      f(\IntEmbed(n))=\Phi_R(\IntEmbed(n))}{%
      \FormulaRefAuto{a=b\dsep c=b\vdash a=c}{9,10}}
    \proofstepwidestar[12]{z=\IntEmbed(n)}{\rA}
    \proofstepwidestar[1,2,6,12]{%
      f(z)=\Phi_R(z)}{%
      \rIE{12,11}}
    \proofstepwidestar[14]{z=-\IntEmbed(n)}{\rA}
    \proofstepwidestar[6]{\IntEmbed(n)\in\mathbb Z}{%
      \BandXXIXProofRefStack{%
        \FormulaRefAuto{F\colon A\to B\dsep x\in A\vdash F(x)\in B}[thm]\\
        \bigl(\FormulaRefAuto{NaturalIntegerEmbeddingFunction}[thm],7\bigr)}}
    \proofstepwidestar[2,6]{%
      f(-\IntEmbed(n))=-f(\IntEmbed(n))}{%
      \FormulaRefAuto{RingHomNegativePreservation}[thm]{2,15}}
    \proofstepwidestar[1,6]{%
      \Phi_R(-\IntEmbed(n))=-\Phi_R(\IntEmbed(n))}{%
      \FormulaRefAuto{RingHomNegativePreservation}[thm]{4,15}}
    \proofstepwidestar[1,2,6]{%
      -f(\IntEmbed(n))=-\Phi_R(\IntEmbed(n))}{%
      \rIE{11,\rII}}
    \proofstepwidestar[1,6]{%
      -\Phi_R(\IntEmbed(n))=\Phi_R(-\IntEmbed(n))}{%
      \FormulaRefAuto{a=b\vdash b=a}{17}}
    \proofstepwidestar[1,2,6]{%
      f(-\IntEmbed(n))=\Phi_R(-\IntEmbed(n))}{%
      \rChain{16,18,19}}
    \proofstepwidestar[1,2,6,14]{%
      f(z)=\Phi_R(z)}{%
      \rIE{14,20}}
    \proofstepwidestar[1,2,6]{%
      f(z)=\Phi_R(z)}{%
      \rOE{8,12,13,14,21}}
    \proofstepwidestar[1,2,3]{%
      f(z)=\Phi_R(z)}{%
      \rEE{5,6,22}}
  \closeproofpartwideindR

  \proofpartwideindR[Eindeutigkeit der Ringabbildung]{%
    \begin{aligned}[t]
      &\operatorname{Ring}(\mathfrak R)\dsep\\[-2pt]
      &\RingHom{f}{\mathfrak Z}{\mathfrak R}\\
      &\vdash f=\Phi_R
    \end{aligned}
  }{%
    \begin{aligned}[t]
      &\operatorname{Ring}(\mathfrak R)\dsep\\[-2pt]
      &\RingHom{f}{\mathfrak Z}{\mathfrak R}\\
      &\vdash f=\Phi_R
    \end{aligned}
  }[IntegerRingHomUnique]
    \proofstepwidestar[1]{%
      \operatorname{Ring}(\mathfrak R)}{\rA}
    \proofstepwidestar[2]{%
      \RingHom{f}{\mathfrak Z}{\mathfrak R}}{\rA}
    \proofstepwidestar[2]{f\colon\mathbb Z\to R}{%
      \FormulaRefAuto{RingHomFunctionAxiom}[ax]{2}}
    \proofstepwidestar[1]{\Phi_R\colon\mathbb Z\to R}{%
      \FormulaRefAuto{IntegerRingMapFunction}[thm]{1}}
    \proofstepwidestar[5]{z\in\mathbb Z}{\rA}
    \proofstepwidestar[1,2,5]{%
      f(z)=\Phi_R(z)}{%
      \FormulaRefAuto{IntegerRingHomPointwiseUnique}[thm]{1,2,5}}
    \proofstepwidestar[1,2]{%
      \begin{aligned}[t]
        &\forall z\in\mathbb Z\,\\[-2pt]
        &f(z)=\Phi_R(z)
      \end{aligned}}{%
      \rUI{\rRI{5,6}}}
    \proofstepwidestar[1,2]{%
      f=\Phi_R}{%
      \FormulaRefAuto{FunctionExtensionality}[thm]{3,4,7}}
  \closeproofpartwideindR

  \proofpartwide{Existenz und Eindeutigkeit}
    \proofstepwidestar[1]{%
      \operatorname{Ring}(\mathfrak R)}{\rA}
    \proofstepwidestar[1]{%
      \RingHom{\Phi_R}{\mathfrak Z}{\mathfrak R}}{%
      \FormulaRefAuto{IntegerRingCanonicalHomomorphism}[thm]{1}}
    \proofstepwidestar[3]{%
      \RingHom{f}{\mathfrak Z}{\mathfrak R}}{\rA}
    \proofstepwidestar[1,3]{%
      f=\Phi_R}{%
      \FormulaRefAuto{IntegerRingHomUnique}[thm]{1,3}}
    \proofstepwidestar[1]{%
      \begin{aligned}[t]
        &\exists! f\,\\[-2pt]
        &\RingHom{f}{\mathfrak Z}{\mathfrak R}
      \end{aligned}}{\UEI{2,3,4}}
  \closeproofpartwide
\end{tabproofsplitwide}

\end{document}
